\documentclass[11pt,oneside,openany]{book}

\usepackage[T1]{fontenc}
\usepackage{lmodern}
\usepackage[letterpaper,margin=1.10in,includeheadfoot]{geometry}
\usepackage{amsmath,amssymb,amsthm,mathtools}
\usepackage{microtype}
\usepackage{xcolor}
\usepackage{titlesec}
\usepackage{fancyhdr}
\usepackage{hyperref}

\ifdefined\pdftrailerid
  \pdftrailerid{MathGenome Advanced Commutative Algebra 2026-11-17}
\fi

\definecolor{ink}{HTML}{172923}
\definecolor{deepgreen}{HTML}{2B5242}
\definecolor{rulegreen}{HTML}{8A9F92}
\hypersetup{
  colorlinks=true,linkcolor=deepgreen,urlcolor=deepgreen,citecolor=deepgreen,
  pdftitle={Advanced Commutative Algebra},
  pdfsubject={Differentials, Smoothness, and Local Algebra},
  pdfauthor={MathGenome},bookmarksnumbered=true
}
\setcounter{tocdepth}{1}
\setcounter{secnumdepth}{2}
\setlength{\headheight}{15pt}
\setlength{\emergencystretch}{2em}

\titleformat{\chapter}[display]
  {\normalfont\color{ink}}{\color{deepgreen}\large\scshape Chapter \thechapter}{9pt}
  {\huge\bfseries}
\titlespacing*{\chapter}{0pt}{4pt}{25pt}
\titleformat{\section}{\normalfont\Large\bfseries\color{ink}}{\thesection}{0.7em}{}
\titleformat{\subsection}{\normalfont\large\bfseries\color{ink}}{\thesubsection}{0.7em}{}

\pagestyle{fancy}
\fancyhf{}
\fancyhead[L]{\small\sffamily\color{deepgreen}MathGenome}
\fancyhead[R]{\small\nouppercase{\leftmark}}
\fancyfoot[C]{\small\thepage}
\renewcommand{\headrulewidth}{0.35pt}
\renewcommand{\chaptermark}[1]{\markboth{\thechapter.\ #1}{}}
\fancypagestyle{plain}{%
  \fancyhf{}\fancyfoot[C]{\small\thepage}\renewcommand{\headrulewidth}{0pt}%
}

\newtheorem{theorem}{Theorem}[chapter]
\newtheorem{lemma}[theorem]{Lemma}
\newtheorem{proposition}[theorem]{Proposition}
\newtheorem{corollary}[theorem]{Corollary}
\theoremstyle{definition}
\newtheorem{definition}[theorem]{Definition}
\newtheorem{example}[theorem]{Example}
\theoremstyle{remark}
\newtheorem{remark}[theorem]{Remark}
\allowdisplaybreaks[1]

\begin{document}
\frontmatter
\hypersetup{pageanchor=false}
\begin{titlepage}
\thispagestyle{empty}
\thispdfpagelabel{Cover}
\begin{flushleft}
\vspace*{0.65in}
{\large\scshape\color{deepgreen}MathGenome\par}
\vspace{0.35in}
{\color{rulegreen}\rule{0.94\textwidth}{0.65pt}\par}
\vspace{0.43in}
{\color{ink}\fontsize{36}{42}\selectfont\bfseries
Advanced\\[4pt] Commutative Algebra\par}
\vspace{0.25in}
{\color{deepgreen}\Large Differentials, Smoothness, and Local Algebra\par}
\vfill
{\large Mathematical lecture notes\par}
\vspace{0.10in}
{\normalsize MathGenome v0.21.0 dependency baseline\par}
\vspace*{0.45in}
\end{flushleft}
\end{titlepage}
\hypersetup{pageanchor=true}

\tableofcontents

\chapter*{Scope and conventions}
\addcontentsline{toc}{chapter}{Scope and conventions}

These notes develop the affine algebra needed for relative differentials and
smooth and \'{e}tale morphisms. Chapter~\ref{ch:differentials} constructs
K\"ahler differentials and computes them under polynomial presentation,
localization, and base change. Chapter~\ref{ch:lifting} uses square-zero
extensions to express lifting and uniqueness. Chapters~\ref{ch:smooth}
and~\ref{ch:etale} add finite presentation, prove the Jacobian criteria,
and give local algebra forms that can be transported to affine schemes.
Part~\ref{part:local} develops regular sequences, the Koszul complex
and depth, followed by Cohen--Macaulay rings, Serre's normality
criterion, and Gorenstein rings with their canonical modules.
The Noetherian local arguments in this part are based on
\cite[Chapters 17, 18, and 21]{eisenbud},
\cite[Chapters 1--3]{brunsherzog}, and the corresponding Stacks chapters.

Rings are commutative with identity, ring maps preserve identity, and the zero
ring is allowed. In particular its spectrum is empty. An $A$-algebra is a ring
equipped with a specified structural map out of $A$. A tower of algebras
always uses compatible structural maps. Every tensor product carries the
displayed scalar structure. The symbols $\operatorname{Spec}(B)$,
$B_{\mathfrak q}$, and $\kappa(\mathfrak q)$ denote the prime spectrum, the
local ring at $\mathfrak q$, and its residue field. For $b\in B$, $B_b$
denotes localization at the powers of $b$. If $B$ is a localization of some
$A$-algebra of finite type, it is \emph{essentially of finite type} over $A$.

A finite-type algebra has a finite set of algebra generators. A finitely
presented algebra is a quotient of a polynomial algebra in finitely many
variables by an ideal with finitely many generators. A finitely presented
module is the cokernel of a map between finite free modules. Finiteness of
an algebra as a module is a separate condition. The base-change convention
for $A\to A'$ is $B'=A'\otimes_A B$.

The preceding MathGenome notes \cite{p103} provide quotients, tensor products,
exact sequences, localization, Nakayama's lemma, and the relevant flatness
results. The abstract algebra notes \cite{d07} supply the corresponding
field-theoretic preliminaries. The separate source map records the edition
and location of each mathematical result, cross-chapter proof dependencies,
and candidate formalization interfaces. The MathGenome and Mathlib versions
used in that map are specified in the accompanying build and review record.
Geometric regularity, relative differential sheaves, the full cotangent
complex, and the development of local cohomology and dualizing complexes
have their own treatments. Finite flat refinements of \'{e}tale covers
and \'{e}tale-local lifting of coprime factorizations belong to the
descent and Henselian topics, respectively \cite[Tags 00UF--00UH]{stacks}.

\mainmatter
\part{Differentials, Smoothness, and \'{E}tale Algebras}
\label{part:differentials}
\chapter{K\"ahler Differentials}\label{ch:differentials}

Differentials record first-order variation of an algebra over its base. The universal property makes them usable with arbitrary coefficient modules, including modules on a quotient, an open subset, or a fibre. Keeping track of the module structure is what distinguishes the two fundamental exact sequences from the base change isomorphism. The main references for this chapter are \cite[Chapter 16]{eisenbud}, \cite[Section 25]{matsumura}, and \cite[Section 10.131]{stacks}.

Throughout the chapter, $A\to B$ is a map of commutative rings with identity, $M$ is a $B$-module, and all algebra maps preserve the identity. The zero ring is allowed. In a tower of rings we use the induced actions without displaying the structure maps where the source and target determine them. For instance, $b_C$ is the image of $b\in B$ in a specified $B$-algebra $C$. We use the universal properties of quotients and tensor products, scalar extension and restriction, and the right exactness of tensoring as in \cite[Chapters 1 and 4]{p103}.

\section{Derivations and the universal differential}

\begin{definition}[Relative derivations]\label{res:KD-01}
An $A$-derivation $D:B\to M$ is an $A$-linear map satisfying
\[
 D(bb')=bD(b')+b'D(b).
\]
Equivalently, $D$ is additive, satisfies this product rule, and vanishes on the image of $A$. The derivations form a $B$-module $\operatorname{Der}_A(B,M)$. A $B$-linear map $M\to N$ induces a map of derivation modules by composition. Two $A$-derivations into the same $B$-module coincide if they coincide on an $A$-algebra generating family of $B$.
\end{definition}

The product rule at $1$ gives $D(1)=2D(1)$ and hence $D(1)=0$. An $A$-linear derivation therefore vanishes on the base. Conversely, an additive derivation vanishing on the base satisfies $D(ab)=aD(b)$. The $B$-module operations are pointwise, and $B$-linear composition preserves the rule. For the last assertion, the difference of the two derivations vanishes on the base and the generators. Its zero set is closed under addition, negation, and multiplication, so it contains the subalgebra they generate.

\begin{theorem}[Construction and universal property]\label{res:KD-02}
There is a $B$-module $\Omega_{B/A}$ with an $A$-derivation $d=d_{B/A}:B\to\Omega_{B/A}$ for which composition with $d$ gives a $B$-linear isomorphism
\[
 \operatorname{Hom}_B(\Omega_{B/A},M)\ \xrightarrow{\ \sim\ }\ \operatorname{Der}_A(B,M)
\]
natural in $M$. One construction takes the free $B$-module on symbols $[b]$ and divides by the submodule generated by
\[
 [b+b']-[b]-[b'],\qquad [bb']-b[b']-b'[b],\qquad [a_B].
\]
The class of $[b]$ is $db$. In particular, the $db$ generate $\Omega_{B/A}$.
\end{theorem}

\begin{proof}
The relations make $d$ a derivation by Definition~\ref{res:KD-01}. Given $D$, the $B$-linear map from the free module sending $[b]$ to $D(b)$ annihilates all three types of relations and descends uniquely. Its inverse correspondence sends a linear map $\lambda$ to $\lambda d$. Both operations are $B$-linear, and postcomposing with a coefficient-module map commutes with the correspondence. This also shows that any two universal pairs have a unique isomorphism carrying one universal derivation to the other.
\end{proof}

\begin{proposition}[The diagonal ideal]\label{res:KD-03}
Let $T=B\otimes_A B$, let $\mu:T\to B$ be multiplication, and set $J=\ker\mu$. The two $B$-actions on $J/J^2$ from the tensor factors agree. There is a $B$-linear isomorphism
\[
 \Omega_{B/A}\ \xrightarrow{\ \sim\ }\ J/J^2,\qquad
 db\longmapsto [1\otimes b-b\otimes1].
\]
\end{proposition}

\begin{proof}
Write $l(b)=b\otimes 1$, $r(b)=1\otimes b$, and $\Delta(b)=r(b)-l(b)$. Then $(r(b)-l(b))u\in J^2$ for $u\in J$, which gives the equality of the actions. The identity
\[
 \Delta(bb')=l(b)\Delta(b')+l(b')\Delta(b)+\Delta(b)\Delta(b')
\]
and $\Delta(a_B)=0$ show that the class of $\Delta$ is a derivation. Theorem~\ref{res:KD-02} gives the displayed map. In the other direction the rule $b\otimes c\mapsto b\,dc$ defines an additive map $F:T\to\Omega_{B/A}$ by $A$-balance. On pure tensors the Leibniz rule gives
\[
 F(uv)=\mu(u)F(v)+\mu(v)F(u).
\]
Both sides are additive in each variable, so this holds on $T$. Thus $F$ kills $J^2$ and restricts to a $B$-linear map $J/J^2\to\Omega_{B/A}$ with $F(\Delta(b))=db$. For $u=\sum_i b_i\otimes c_i\in J$ one has $\sum_i b_ic_i=0$, so $u=\sum_i l(b_i)\Delta(c_i)$. The inverse identities now follow on generators.
\end{proof}

\begin{proposition}[Derivations and square-zero lifts]\label{res:KD-04}
Give $B\oplus M$ the ring law
\[
 (b,m)(b',m')=(bb',bm'+b'm)
\]
and its $A$-algebra structure through $B\to B\oplus M$, $b\mapsto(b,0)$. The $A$-algebra sections of $B\oplus M\to B$ are the maps $b\mapsto (b,D(b))$ for $D\in\operatorname{Der}_A(B,M)$.

More generally, let $C$ be an $A$-algebra and $I\subset C$ an ideal with $I^2=0$. Fix $\bar f:B\to C/I$ over $A$. The quotient $C/I$ acts on $I$ by lifting scalars to $C$. If $f:B\to C$ is a lift of $\bar f$, all lifts are uniquely of the form $f+D$ for $D\in\operatorname{Der}_A(B,I)$, with the $B$-action induced by $\bar f$.
\end{proposition}

\begin{proof}
In the first assertion, comparing the second coordinates in the product law gives exactly the product rule, and comparing the images of $A$ gives vanishing on the base. For the second, multiplication on $I$ is independent of the scalar lift since $I^2=0$. If $g$ is a second lift then $D=g-f$ has image in $I$. Expanding $(f(b)+D(b))(f(b')+D(b'))$ leaves no term $D(b)D(b')$. The remaining terms show that $D$ satisfies the asserted product rule. Conversely the expansion shows that adding such a derivation to $f$ gives a multiplicative map. It is additive, unital because $D(1)=0$, and agrees with the structure map on $A$.
\end{proof}

\begin{proposition}[Naturality in the algebra]\label{res:KD-05}
For a commutative square of ring maps $(A\to B)\to(A'\to B')$ with target map $f:B\to B'$, there is a natural $B$-linear map
\[
 \Omega_{B/A}\longrightarrow\Omega_{B'/A'},\qquad db\longmapsto d(f(b)),
\]
where the codomain uses restriction along $f$. These maps respect identity squares and composition.
\end{proposition}

\begin{proof}
The map $d_{B'/A'} f$ is an $A$-derivation when the coefficients are restricted along $f$. Theorem~\ref{res:KD-02} gives its unique linear factorization. For the identities and compositions, both candidate maps take the same value on every generator $db$.
\end{proof}

\begin{corollary}[Surjective structural maps]\label{res:KD-06}
If $A\to B$ is surjective then $\Omega_{B/A}=0$. In particular $\Omega_{A/A}=0$ and $\Omega_{(A/I)/A}=0$ for every ideal $I$ of $A$.
\end{corollary}

\begin{proof}
Each $b$ comes from the base, so every generator $db$ vanishes.
\end{proof}

\section{Free algebras, towers, and quotients}

\begin{proposition}[Polynomial differentials]\label{res:KD-07}
For the polynomial algebra $P=A[X_i:i\in\Lambda]$, the elements $dX_i$ form a free $P$-basis of $\Omega_{P/A}$. For every $f\in P$,
\[
 df=\sum_{i\in\Lambda} \frac{\partial f}{\partial X_i}\,dX_i,
\]
with finite support. In one variable this is $df=f'(X)dX$.
\end{proposition}

\begin{proof}
Take the free $P$-module with basis $e_i$. Formal differentiation on monomials, extended $A$-linearly, gives the derivation $f\mapsto\sum_i(\partial_i f)e_i$. Indeed the exponents add when two monomials are multiplied, which gives the product rule on monomials and hence on their finite sums. For any coefficient module and any family $m_i$, this derivation followed by $e_i\mapsto m_i$ assigns $X_i\mapsto m_i$. The generating family criterion of Definition~\ref{res:KD-01} makes that derivation unique. Thus the free module has the universal property.
\end{proof}

\begin{proposition}[Maps in a tower]\label{res:KD-08}
For $A\to B\to C$ there are canonical $C$-linear maps
\[
 u:C\otimes_B\Omega_{B/A}\longrightarrow\Omega_{C/A},
 \qquad v:\Omega_{C/A}\longrightarrow\Omega_{C/B}
\]
with $u(c\otimes db)=c\,d_{C/A}(b_C)$ and $v(d_{C/A}c)=d_{C/B}c$.
\end{proposition}

\begin{proof}
Use Proposition~\ref{res:KD-05} for the top map $B\to C$ with unchanged base $A$ and then extend scalars to get $u$. To get $v$, use the identity as top map and $A\to B$ as base map. Both formulas then follow from the values on generators.
\end{proof}

\begin{theorem}[The first fundamental exact sequence]\label{res:KD-09}
With $u,v$ as in Proposition~\ref{res:KD-08}, the sequence
\[
 C\otimes_B\Omega_{B/A}\xrightarrow{\ u\ }\Omega_{C/A}
 \xrightarrow{\ v\ }\Omega_{C/B}\longrightarrow 0
\]
is exact. The first arrow can have nonzero kernel.
\end{theorem}

\begin{proof}
The derivation $d_{C/B}$ vanishes on $B$, so $v u=0$. The quotient $Q=\Omega_{C/A}/\operatorname{im}(u)$ receives a derivation from $C$ vanishing on $B$. For any $C$-module $N$, linear maps $Q\to N$ correspond to $A$-derivations $C\to N$ vanishing on the image of $B$. These are precisely the $B$-derivations. The induced map $Q\to\Omega_{C/B}$ therefore identifies their universal pairs and is an isomorphism. For the kernel example take $A=C=k$ for a field $k$, $B=k[t]$, and $B\to C$ given by $t\mapsto 0$. The first module is the nonzero vector space $k\,dt$ by Proposition~\ref{res:KD-07}, while the middle module is zero by Corollary~\ref{res:KD-06}.
\end{proof}

\begin{proposition}[A surjective map of targets]\label{res:KD-10}
In the square of Proposition~\ref{res:KD-05}, suppose that $f:B\to B'$ is surjective. The induced map of differentials is surjective. Its kernel is generated over $B$ by the $db$ with $f(b)\in\operatorname{im}(A'\to B')$.
\end{proposition}

\begin{proof}
Every generator $db'$ lifts, which proves surjectivity. Let $L$ be the proposed kernel and put $Q=\Omega_{B/A}/L$. By naturality $L$ maps to zero. If $i\in\ker f$ and $b\in B$, then $i\,db=d(ib)-b\,di$ vanishes in $Q$. Since the differentials generate, $Q$ is a $B'$-module. The derivation $B\to Q$ vanishes on $\ker f$ and therefore descends additively to $B'$. The descended map satisfies the product rule by choosing lifts, and it annihilates the image of $A'$ by the definition of $L$. Its linear factorization $\Omega_{B'/A'}\to Q$ is inverse to $Q\to\Omega_{B'/A'}$ on generators on both sides.
\end{proof}

\begin{theorem}[The conormal exact sequence]\label{res:KD-11}
Let $A\to P\twoheadrightarrow B$ be a tower with $I=\ker(P\to B)$. There is a sequence of $B$-modules
\[
 I/I^2\xrightarrow{\ j\ } B\otimes_P\Omega_{P/A}
 \longrightarrow\Omega_{B/A}\longrightarrow0,\qquad
 j([f])=1\otimes df .
\]
It is exact at the middle and last terms. The first arrow need not be injective.
\end{theorem}

\begin{proof}
The action on $I/I^2$ factors through $B$. In the middle term, $1\otimes d(fg)=0$ for $f,g\in I$ by the product rule, and for $p\in P$ and $f\in I$ one has $1\otimes d(pf)=\bar p\otimes df$. Thus $j$ is well defined and $B$-linear. For a $B$-module $N$, a map $B\otimes_P\Omega_{P/A}\to N$ corresponds to an $A$-derivation $P\to N$ by scalar extension and Theorem~\ref{res:KD-02}. It vanishes on the image of $j$ exactly if its derivation kills $I$. In that case the derivation descends uniquely to $B$ since the action on $N$ is already through $B$. Hence the cokernel of $j$ represents $\operatorname{Der}_A(B,-)$, and its canonical map to $\Omega_{B/A}$ is an isomorphism. For the failure of injectivity take $A=P=\mathbb Z$ and $B=\mathbb F_p$ for a prime $p$. Then $I/I^2=(p)/(p^2)$ is nonzero, while $\Omega_{\mathbb Z/\mathbb Z}=0$.
\end{proof}

\begin{corollary}[A presentation by the Jacobian matrix]\label{res:KD-12}
Suppose $B$ has a presentation
\[
 B=A[X_1,\ldots,X_n]/(f_1,\ldots,f_m).
\]
As a $B$-module, $\Omega_{B/A}$ is the cokernel of
\[
 B^m\longrightarrow B^n,\qquad
 e_j\longmapsto \sum_{i=1}^n\overline{\frac{\partial f_j}{\partial X_i}}\,e_i .
\]
The $i$th basis vector on the right maps to $d\bar X_i$. Thus rows are indexed by variables and columns by relations.
\end{corollary}

\begin{proof}
The images of the $f_j$ generate $I/I^2$ over $B$: write each $f\in I$ as a finite sum of multiple terms $h_j f_j$ and reduce the coefficients modulo $I$. Composing the resulting surjection $B^m\to I/I^2$ with the first map of Theorem~\ref{res:KD-11} has the same image. Proposition~\ref{res:KD-07} computes this composite by partial derivatives and identifies its target with $B^n$.
\end{proof}

\begin{corollary}[Finite generation and finite presentation]\label{res:KD-13}
If $B$ is of finite type over $A$, then $\Omega_{B/A}$ is finitely generated over $B$. If $B$ is finitely presented as an $A$-algebra, then $\Omega_{B/A}$ is finitely presented as a $B$-module.
\end{corollary}

\begin{proof}
Choose a surjection $A[X_1,\ldots,X_n]\to B$. The conormal sequence and the polynomial basis give a surjection $B^n\to\Omega_{B/A}$. With finitely many defining relations, Corollary~\ref{res:KD-12} presents its kernel as the image of a finite free module. Here finite presentation of the algebra determines the relations and finite presentation of the module is the conclusion.
\end{proof}

\begin{proposition}[A split quotient]\label{res:KD-14}
If $P\twoheadrightarrow B$ in Theorem~\ref{res:KD-11} has an $A$-algebra section $s$, then its conormal sequence is split short exact. The left inverse of $j$ corresponds to the derivation
\[
 P\longrightarrow I/I^2,\qquad p\longmapsto[p-s(\bar p)] .
\]
\end{proposition}

\begin{proof}
The difference lies in $I$. To verify the product rule, write $p=s(\bar p)+i$ and $q=s(\bar q)+i'$ with $i,i'\in I$. Modulo $I^2$, the difference $pq-s(\overline{pq})$ has class $\bar p[i']+\bar q[i]$. The assignment is additive and kills the base since $s$ is an $A$-algebra map. Its linear factorization out of $B\otimes_P\Omega_{P/A}$ sends $1\otimes df$ to $[f]$ for $f\in I$, so it retracts $j$. Letting $r$ denote this retraction, the endomorphism $1-jr$ of the middle term kills $\operatorname{im}(j)$ and induces a section from its cokernel. By Theorem~\ref{res:KD-11} the cokernel is $\Omega_{B/A}$.
\end{proof}

\begin{proposition}[Restriction from a thicker quotient]\label{res:KD-15}
Let $I\subset B$ and $n\ge 1$. Set $B_i=B/I^i$. The map of differentials gives an isomorphism
\[
 B_n\otimes_B\Omega_{B/A}\ \xrightarrow{\ \sim\ }\
 B_n\otimes_{B_{n+1}}\Omega_{B_{n+1}/A}.
\]
\end{proposition}

\begin{proof}
With $P=B$ in Theorem~\ref{res:KD-11}, the quotient by $I^{n+1}$ expresses $\Omega_{B_{n+1}/A}$ as the quotient of $B_{n+1}\otimes_B\Omega_{B/A}$ by the image of $I^{n+1}/I^{2n+2}$. The product rule gives $d(I^{n+1})\subset I^n\Omega_{B/A}$ by differentiating a product of $n+1$ elements of $I$. Finite sums of these products generate the ideal, so the same containment holds for all of $I^{n+1}$. After applying the right exact functor $B_n\otimes_{B_{n+1}}-$, each image of a generator of the relation submodule is zero. The displayed map is the resulting cokernel isomorphism. For $B=k[t]$, $I=(t)$, and $n=1$, the common module is $k\,dt$ by Proposition~\ref{res:KD-07}, whereas $\Omega_{B_1/k}=0$ by Corollary~\ref{res:KD-06}.
\end{proof}

\section{Localization, base change, and operations on algebras}

\begin{proposition}[Localization of derivations and differentials]\label{res:KD-16}
Let $S$ be a multiplicative subset of $B$, and set $T=S^{-1}B$. Every $A$-derivation $D:B\to N$ into a $T$-module extends uniquely to an $A$-derivation on $T$. The formula is
\[
 \widetilde D(b/s)=s^{-1}D(b)-bs^{-2}D(s).
\]
In particular $S^{-1}\Omega_{B/A}\simeq \Omega_{T/A}$, with $db/1\mapsto d(b/1)$. If $U$ is a multiplicative subset of $A$ whose image is invertible in $B$, then $\Omega_{B/A}\to\Omega_{B/U^{-1}A}$ is an isomorphism.
\end{proposition}

\begin{proof}
The map $B\to T\oplus N$ given by $b\mapsto(b/1,D(b))$ is an $A$-algebra map by Proposition~\ref{res:KD-04}. The images of elements $s\in S$ are units, with inverses $(s^{-1},-s^{-2}D(s))$ under the square-zero product law. The ring localization universal property therefore gives a map $T\to T\oplus N$. Its first component is the identity by uniqueness of localization. Its second component is the extended derivation, and multiplying by the displayed inverse computes it on fractions. Any extension must define the same localization map, which also establishes well-definedness when $B$ has zero divisors or $0\in S$. To identify the modules, for every $T$-module $N$ use the two isomorphisms
\[
 \operatorname{Hom}_T(T\otimes_B\Omega_{B/A},N)
 \simeq\operatorname{Der}_A(B,N)
 \simeq\operatorname{Der}_A(T,N).
\]
For the last claim, an $A$-derivation kills the images of $U$ and their inverses by applying the product rule to $u\,u^{-1}=1$. It consequently kills the image of $U^{-1}A$. The coefficient derivation modules for the two bases agree, and Theorem~\ref{res:KD-02} identifies the universal modules.
\end{proof}

\begin{theorem}[Arbitrary base change]\label{res:KD-17}
Let $A\to A'$ be any ring map and $B'=A'\otimes_A B$. There is a canonical isomorphism of $B'$-modules
\[
 B'\otimes_B\Omega_{B/A}\ \xrightarrow{\ \sim\ }\ \Omega_{B'/A'},
 \qquad b'\otimes db\longmapsto b'\,d(1\otimes b).
\]
It is compatible with the identity base change and composition of base changes via the canonical tensor associativity isomorphisms.
\end{theorem}

\begin{proof}
For a $B'$-module $N$, restricting an $A'$-derivation of $B'$ along $b\mapsto 1\otimes b$ gives an $A$-derivation of $B$. Conversely, $D:B\to N$ extends by the rule
\[
 a'\otimes b\longmapsto a'D(b).
\]
It is $A$-balanced since $D(ab)=aD(b)$, and the product rule on pure tensors follows from that of $D$. It then holds on sums by additivity. Also $D(1)=0$ ensures the extension kills $A'$. Any $A'$-derivation is determined by its values on $1\otimes b$ since it is $A'$-linear. These inverse correspondences are natural in $N$. Scalar extension and Theorem~\ref{res:KD-02} identify the representing modules. The resulting map is the one obtained by extending Proposition~\ref{res:KD-05} to $B'$, so the formula follows. Under successive base changes both paths agree on the extended $db$ generators. No flatness assumption enters the argument.
\end{proof}

\begin{proposition}[Tensor products of algebras]\label{res:KD-18}
For $C=B_1\otimes_A B_2$ there is a canonical $C$-linear isomorphism
\[
 \Omega_{C/A}\ \simeq\
 (C\otimes_{B_1}\Omega_{B_1/A})\oplus
 (C\otimes_{B_2}\Omega_{B_2/A}).
\]
It sends $d(b_1\otimes b_2)$ to the pair whose terms are $b_2(1\otimes db_1)$ and $b_1(1\otimes db_2)$, with the coefficients mapped to $C$. In particular for $C=B_1[X_i:i\in\Lambda]$ the second summand can be identified with $\bigoplus_i C\,dX_i$.
\end{proposition}

\begin{proof}
Let $N$ be a $C$-module. Restrict a derivation along the two factors to obtain $D_j:B_j\to N$. The inverse correspondence is
\[
 D(b_1\otimes b_2)=b_2D_1(b_1)+b_1D_2(b_2).
\]
Each term is additive and $A$-balanced since the two $A$-actions on $N$ agree. Expanding the product rule in each factor gives the product rule on pure tensors and hence on all of $C$. The formula kills the base, and any derivation with these restrictions must satisfy the formula. The representing module is the direct sum by the universal property and scalar extension. For the polynomial case apply Proposition~\ref{res:KD-07} to the second factor.
\end{proof}

\begin{proposition}[Directed colimits and coequalizers]\label{res:KD-19}
Let $(A_i\to B_i)$ be a system of ring maps indexed by a nonempty directed set, and put $A=\varinjlim A_i$ and $B=\varinjlim B_i$. The transition maps give an isomorphism of $B$-modules
\[
 \varinjlim_i (B\otimes_{B_i}\Omega_{B_i/A_i})\ \simeq\ \Omega_{B/A}.
\]
For two $A$-algebra maps $u,v:P\rightrightarrows Q$ with coequalizer $C$, there is also a natural presentation of $\Omega_{C/A}$ as the cokernel of the difference of the two induced maps
\[
 C\otimes_P\Omega_{P/A}\rightrightarrows C\otimes_Q\Omega_{Q/A}.
\]
The $P$-action on $C$ is the common composite through $u$ and $v$.
\end{proposition}

\begin{proof}
Fix a $B$-module $N$. A compatible family of $A_i$-derivations $B_i\to N$ defines an additive function on the colimit. Indeed values agree for representatives that become equal at a common stage, and addition can also be evaluated at a common stage. Products at a common stage satisfy Leibniz there. The resulting map kills the image of $A$ since every element of $A$ is represented in some $A_i$. The construction is inverse to restriction. Applying Theorem~\ref{res:KD-02}, scalar extension, and the colimit universal property gives the first isomorphism. For the second, $C$ is $Q$ modulo the ideal generated by the elements $u(p)-v(p)$. A derivation $Q\to N$ to a $C$-module descends exactly when the two restrictions to $P$ coincide. Vanishing on the generating differences suffices since the ideal itself annihilates $N$ and Leibniz extends the vanishing to all their multiples. Translating this condition through the universal property gives the claimed cokernel.
\end{proof}

\begin{proposition}[Finite products of algebras]\label{res:KD-20}
For a finite family of $A$-algebras $(B_i)$, the projections from $B=\prod_i B_i$ induce an isomorphism of $B$-modules
\[
 \Omega_{B/A}\ \simeq\ \prod_i\Omega_{B_i/A}.
\]
\end{proposition}

\begin{proof}
Let $e_i\in B$ be the coordinate idempotents. For any derivation $D$, one has $(1-2e_i)D(e_i)=0$ by differentiating $e_i^2=e_i$. The coefficient is a unit since its square is $1$, so $D(e_i)=0$. For a $B$-module $N$, decompose $N=\prod_i e_iN$ using $\sum_i e_i=1$ and orthogonality. The component $e_iD:B\to e_iN$ depends only on the $i$th coordinate since $e_iD(b)=D(e_i b)$. Together these components identify $A$-derivations from $B$ with families of $A$-derivations $B_i\to e_iN$. In the inverse direction the formula forms a derivation coordinatewise. The universal property now gives the displayed map, including the empty family case since every module over the zero ring is zero.
\end{proof}

\section{Calculations and cotangent fibres}

\begin{example}[Hypersurface calculations]\label{res:KD-21}
Let $P$ be the polynomial ring on $X_1,\ldots,X_n$ over $A$ and set $B=P/(f)$. The module $\Omega_{B/A}$ has generators $d\bar X_i$ and a single defining relation
\[
 \sum_i \overline{\partial_i f}\,d\bar X_i=0.
\]
In one variable the module is $B/(f'(\bar X))\,d\bar X$. For $k[x,y]/(xy)$ the relation is $y\,dx+x\,dy=0$. In $k[\epsilon]/(\epsilon^2)$ it is $2\epsilon\,d\epsilon=0$. Over a field of characteristic two this last module is free of rank one and the algebra contains the nonzero nilpotent $\epsilon$.
\end{example}

These follow by applying Corollary~\ref{res:KD-12} with one relation. In the latter example, when $k$ is a field and its characteristic is different from two, the module is instead the copy of $k\,d\epsilon$ with $\epsilon$ acting by zero.

\begin{proposition}[Differentials of field extensions]\label{res:KD-22}
Let $L=K[\alpha]$ be a finite simple algebraic field extension with minimal polynomial $f$. Then $\Omega_{L/K}\simeq L/(f'(\alpha))\,d\alpha$. In particular $\Omega_{L/K}=0$ if the extension is separable. If $\operatorname{char}K=p>0$ and $X^p-a$ is irreducible, the extension obtained by adjoining a root has a one-dimensional module of differentials. More generally, for any separable algebraic extension $L/K$ and any ring map $A\to K$, the canonical map is an isomorphism
\[
 L\otimes_K\Omega_{K/A}\ \xrightarrow{\ \sim\ }\ \Omega_{L/A},
\]
so $\Omega_{L/K}=0$ also in the infinite separable algebraic case.
\end{proposition}

\begin{proof}
Apply Example~\ref{res:KD-21} to $K[X]/(f)$. In the separable case $f'(\alpha)$ is a nonzero element of a field, so is invertible. For $X^p-a$ its derivative is identically zero. For the comparison, first take $L/K$ finite separable and use the primitive element theorem \cite[Chapter 12, Primitive Element Theorem]{d07}. For an $L$-module $N$ and $A$-derivation $D:K\to N$, choose the only possible value
\[
 z=-f'(\alpha)^{-1}\sum_i \alpha^i D(f_i)
 \qquad \text{where}\qquad f(X)=\sum_i f_i X^i.
\]
For arbitrary $g(X)=\sum_i g_i X^i$ define $D_P(g)=\sum_i\alpha^i D(g_i)+g'(\alpha)z$. The formula is additive. On a product, the coefficient term satisfies Leibniz by expanding $D(g_i h_j)$ and the polynomial derivative term satisfies Leibniz by formal differentiation. It kills $A$ and $f$, and the product rule then shows that it kills $(f)$ since $f(\alpha)=0$ in $L$. Therefore it gives a derivation on $L$ extending $D$, and generation by $K$ and $\alpha$ shows that it is the unique such extension. An arbitrary separable algebraic extension is the directed union of its finite separable subextensions, since finitely many separable algebraic elements generate a finite separable extension. The unique extensions to each stage agree on overlaps and define the unique derivation on the union, as in Proposition~\ref{res:KD-19}. Scalar extension and the universal property now give the stated comparison. Finally set $A=K$ and apply Corollary~\ref{res:KD-06}.
\end{proof}

\begin{proposition}[Cotangent and tangent at a field-valued point]\label{res:KD-23}
Let $x:B\to K$ be an $A$-algebra point to a field $K$. The $K$-dual of $K\otimes_B\Omega_{B/A}$ is naturally $\operatorname{Der}_A(B,K)$, where $B$ acts on $K$ via $x$. Equivalently, it is the space of $A$-algebra lifts to $K[\epsilon]/(\epsilon^2)$ reducing to $x$, with the constant lift as zero. If $A=K=k$ and $x$ is a $k$-algebra point with $\mathfrak m=\ker x$, there is a canonical isomorphism of $k$-vector spaces
\[
 \mathfrak m/\mathfrak m^2\ \xrightarrow{\ \sim\ }\ k\otimes_B\Omega_{B/k},
 \qquad [b]\longmapsto 1\otimes db .
\]
\end{proposition}

\begin{proof}
The dual description follows from scalar extension and Theorem~\ref{res:KD-02}. The constant lift through $K$ and Proposition~\ref{res:KD-04} identify a tangent vector $D$ with $b\mapsto x(b)+\epsilon D(b)$. For a $k$-point, the structural map $k\to B$ splits $x$, and $x$ is surjective. Proposition~\ref{res:KD-14} gives a split conormal sequence whose last term $\Omega_{k/k}$ is zero. This makes its first map the required isomorphism.
\end{proof}

\begin{corollary}[Compatibility with the fibre over the base point]\label{res:KD-24}
Let $\mathfrak q\in\operatorname{Spec}B$ and $\mathfrak p$ be its inverse image in $A$. Set $A'=\kappa(\mathfrak p)$, $B'=A'\otimes_A B$, and $K=\kappa(\mathfrak q)$. The natural map $B'\to K$ over $A'$ gives a canonical isomorphism
\[
 K\otimes_B\Omega_{B/A}\ \xrightarrow{\ \sim\ }\
 K\otimes_{B'}\Omega_{B'/A'} .
\]
One can also replace $B$ by $B_{\mathfrak q}$ on the left, using $\Omega_{B_{\mathfrak q}/A_{\mathfrak p}}$.
\end{corollary}

\begin{proof}
The quotient of $B_{\mathfrak q}$ by its maximal ideal is $K$. The map $A\to K$ kills $\mathfrak p$ and inverts its complement, so it factors through $A'$. This and the map $B\to K$ give the displayed point on the fibre. Tensor Theorem~\ref{res:KD-17} over $B'$ with $K$, and use associativity of tensor product. To compare with the local rings, first apply Proposition~\ref{res:KD-16} to $B_{\mathfrak q}$. The complement of $\mathfrak p$ becomes invertible there, so the same proposition lets us replace $A$ by $A_{\mathfrak p}$. The resulting maps send the image of each $db$ to its corresponding differential under the point and so give the claimed compatibility.
\end{proof}

\begin{example}[Three changes of ring]\label{res:KD-25}
For $f\in A$, the relative module $\Omega_{A_f/A}$ is zero, by localization of $\Omega_{A/A}$. For a field $K$, the module $\Omega_{K(t)/K}$ is free on $dt$, by localizing $\Omega_{K[t]/K}$ at the nonzero polynomials. Finally $\Omega_{\mathbb F_p/\mathbb Z}=0$ because the structural map is surjective. The latter map fails to be flat: multiplication by $p$ on $\mathbb Z$ is injective, but tensoring it with $\mathbb F_p$ gives the zero endomorphism of a nonzero module.
\end{example}

The zero module in the first example measures a change whose source is the ring being localized. In the rational function example the base stays $K$, so the differential of the polynomial variable survives. The third example will be used when comparing infinitesimal uniqueness with infinitesimal existence.
\chapter{Square-Zero Lifting and Formal Properties}
\label{ch:lifting}

An algebra map into a quotient by a square-zero ideal specifies the values of
functions on a first-order thickening up to errors in the ideal. The lifting
problem separates into existence and uniqueness. The differential module
controls the freedom to change a lift, and a polynomial presentation identifies
the equations that a lift has to satisfy.

For the classical infinitesimal criteria and the presentation formulation,
see \cite[Sections 10.138, 10.148, and 10.150]{stacks}.
For the relation to the difference of lifts and the obstruction group
in deformation theory, see \cite[Proposition 3.10 and Lemma 4.5]{hartshorne}.

\section{Infinitesimal lifting and derivations}

\begin{definition}[Formal lifting properties]\label{res:FL-01}
Let $A\to B$ be a ring map. For every $A$-algebra $C$ and ideal $N\subseteq C$
with $N^2=0$, consider the reduction map
\[
  \operatorname{Hom}_{A\text{-alg}}(B,C)\longrightarrow
  \operatorname{Hom}_{A\text{-alg}}(B,C/N).
\]
The map $A\to B$ is \emph{formally smooth} if every such reduction map is
surjective, \emph{formally unramified} if every such map is injective, and
\emph{formally \'{e}tale} if every such map is bijective. No finite generation
condition enters these definitions.
\end{definition}

\begin{lemma}[Lifting through a nilpotent ideal]\label{res:FL-02}
Each of the properties in Definition~\ref{res:FL-01} is equivalent to the same
condition with $N$ ranging over all nilpotent ideals.
\end{lemma}
\begin{proof}
If $N^m=0$, the reduction $C\to C/N$ is the composite of the maps
$C/N^{j+1}\to C/N^j$ for $j=1,\ldots,m-1$. Their kernels $N^j/N^{j+1}$
have square zero because $2j\ge j+1$. Apply the appropriate condition at each
step, preserving existence, uniqueness, or both under composition. Conversely,
square-zero ideals are nilpotent.
\end{proof}

\begin{proposition}[The set of lifts]\label{res:FL-03}
Fix an $A$-algebra map $u:B\to C/N$ with $N^2=0$. Regard $N$ as a
$B$-module via $u$ and the $C/N$-module structure obtained from multiplication
in $C$. If $u$ has a lift $v$, the map
\[
  \operatorname{Hom}_B(\Omega_{B/A},N)\longrightarrow
  \{\text{$A$-algebra lifts of }u\},\qquad h\longmapsto v+h\,d_{B/A}
\]
is a bijection. Changing $v$ translates the parametrization. In particular the
group $\operatorname{Hom}_B(\Omega_{B/A},N)$ acts freely and transitively on
every nonempty fiber of the reduction map.
\end{proposition}
\begin{proof}
Result~\ref{res:KD-04} proves that multiplication gives the quotient action
on $N$ and that the difference with $v$ gives a bijection from lifts to
derivations into this module. Apply the universal property
in Result~\ref{res:KD-02}. If the reference lift is changed to
$v+h_0d_{B/A}$, the expression $v+h\,d_{B/A}$ has new coordinate $h-h_0$.
This proves both the bijection and its translation rule.
\end{proof}

\begin{proposition}[The unramified criterion]\label{res:FL-04}
For any $A$-algebra $B$, formal unramifiedness is equivalent to
$\Omega_{B/A}=0$.
\end{proposition}
\begin{proof}
If the module vanishes, Proposition~\ref{res:FL-03} gives at most one lift for
every square-zero problem. Conversely, put $C=B\ltimes\Omega_{B/A}$, the
square-zero extension with product $(b,m)(b',m')=(bb',bm'+b'm)$.
Modulo the second summand, both $b\mapsto(b,0)$ and
$b\mapsto(b,d_{B/A}b)$ lift $\operatorname{id}_B$. Formal unramifiedness forces
these lifts to coincide. Then $d_{B/A}=0$, and the classes $db$
generate $\Omega_{B/A}$ by Result~\ref{res:KD-02}.
\end{proof}

\section{The obstruction from a presentation}

\begin{proposition}[Polynomial lifts]\label{res:FL-05}
A polynomial $A$-algebra $A[X_\lambda:\lambda\in\Lambda]$ on any set of
variables is formally smooth. In fact it lifts along every surjection of
$A$-algebras.
\end{proposition}
\begin{proof}
For each variable choose a preimage of its prescribed value under the
surjection. Evaluation at these preimages defines the lift, because every
polynomial involves only finitely many variables.
\end{proof}

\begin{proposition}[The linear lifting equation]\label{res:FL-06}
Let $P\twoheadrightarrow B$ have kernel $I$ and suppose $P$ is formally smooth
over $A$. Put $L=I/I^2$ and $F=B\otimes_P\Omega_{P/A}$, and write
$\delta:L\to F$ for the conormal map. For a square-zero lifting problem
$u:B\to C/N$, choose a lift $\phi:P\to C$ of $P\to B\xrightarrow{u} C/N$.
The map
\[
  o_\phi:L\longrightarrow N,\qquad [f]\longmapsto\phi(f)
\]
is $B$-linear. Its class in
\[
  \operatorname{coker}\!\bigl(
     \operatorname{Hom}_B(F,N)\xrightarrow{\delta^*}
       \operatorname{Hom}_B(L,N)\bigr)
\]
is independent of $\phi$ and vanishes exactly when $u$ lifts to $C$.
\end{proposition}
\begin{proof}
$\phi(I)\subseteq N$ and hence $\phi(I^2)=0$. For $a\in P$ and $f\in I$,
$\phi(af)=\phi(a)\phi(f)$, and the action of $\phi(a)$ on $N$ agrees
with the action of $u(\bar a)$. This proves the linearity claim. By
Proposition~\ref{res:FL-03} applied to $P$, every other choice of
$\phi$ is $\phi+h\,d_{P/A}$ for a $P$-linear map $\Omega_{P/A}\to N$.
Since the $P$-action factors through $B$, we view $h$ as a $B$-linear map
$F\to N$ by extension of scalars. Evaluating on $f\in I$ changes $o_\phi$ by
$h\delta$, as $\delta([f])=1\otimes df$ in Result~\ref{res:KD-11}.
If $o_\phi=h\delta$, the adjusted lift $\phi-h\,d_{P/A}$ kills $I$ and
factors through $B$. Conversely, a lift of $u$ composed with $P\to B$ has
zero obstruction map.
\end{proof}

For a finite presentation $P=A[X_1,\ldots,X_n]$ with
$I=(f_1,\ldots,f_m)$, choose representatives $x_i\in C$ of $u(\bar X_i)$.
Correcting $x_i$ by $t_i\in N$ changes $f_j(x_1,\ldots,x_n)$ by
$\sum_i u(\overline{\partial_i f_j})t_i$. This follows from
Result~\ref{res:KD-12} and the formula for $h\delta$ above, or by expanding
monomials and discarding terms of order two in $N$. The lifting equation is
therefore to solve, for every $j$,
\[
  \sum_i u(\overline{\partial_i f_j})t_i=-f_j(x_1,\ldots,x_n)
  \quad\text{in }N .
\]

\begin{theorem}[The splitting criterion]\label{res:FL-07}
Let $P\twoheadrightarrow B$ with kernel $I$, where $P$ is a formally smooth
$A$-algebra. The following conditions are equivalent.
\begin{enumerate}
  \item $B$ is formally smooth over $A$.
  \item $P/I^2\to B$ admits an $A$-algebra section.
  \item The conormal sequence
  \[
    0\longrightarrow I/I^2\xrightarrow{\delta}
    B\otimes_P\Omega_{P/A}\longrightarrow\Omega_{B/A}\longrightarrow 0
  \]
  is split exact as a sequence of $B$-modules.
\end{enumerate}
\end{theorem}
\begin{proof}
The kernel of $P/I^2\to B$ is $I/I^2$, which has square zero. The identity
of $B$ lifts to a section when $B$ is formally smooth. For a given section
$s:B\to P/I^2$, the formula
\[
  D(p)=(p\bmod I^2)-s(p\bmod I)
\]
defines an $A$-derivation into $I/I^2$ by Proposition~\ref{res:FL-03}:
the two maps from $P$ used in the difference reduce to the same map to $B$.
This derivation yields a $B$-linear $\rho:B\otimes_P\Omega_{P/A}\to I/I^2$.
On $f\in I$ it gives $\rho\delta([f])=[f]$, so $\rho$ retracts $\delta$.
The right exact sequence of Result~\ref{res:KD-11} is therefore split short
exact. Finally, if such a retraction exists, every obstruction map
$o_\phi$ equals $(o_\phi\rho)\delta$. Proposition~\ref{res:FL-06} then
provides a lift for every square-zero test. The implication from a section to
the split sequence works for an arbitrary surjection onto $B$.
\end{proof}

\begin{corollary}[Projectivity and the two-term criterion]\label{res:FL-08}
For a polynomial presentation $P\twoheadrightarrow B$, the algebra $B$ is
formally smooth over $A$ if and only if $\delta$ is injective and
$\Omega_{B/A}$ is projective over $B$. In that case $I/I^2$ and
$\Omega_{B/A}$ are direct summands of the free $B$-module
$B\otimes_P\Omega_{P/A}$. If $B$ is finitely presented over $A$, both
summands are finite projective for a finite polynomial presentation.
\end{corollary}
\begin{proof}
By Proposition~\ref{res:FL-05} and Theorem~\ref{res:FL-07}, formal smoothness
gives a split conormal sequence. By Result~\ref{res:KD-07}, its middle term
is free on the polynomial variables. The converse follows because a
surjection onto a projective module splits, and the injective conormal map
then identifies its domain with the kernel. A finite polynomial presentation
makes the middle term finite free, whose summands are finitely generated.
\end{proof}

\section{Change of rings and local tests}

\begin{proposition}[Composition and arbitrary base change]\label{res:FL-09}
Each property in Definition~\ref{res:FL-01} is preserved by composition and
by arbitrary base change. In a tower $A\to B\to D$, if $A\to B$ is
formally unramified and $A\to D$ is formally smooth, then $B\to D$ is
formally smooth. If $A\to D$ is formally unramified, so is $B\to D$.
\end{proposition}
\begin{proof}
For existence under composition, lift the restriction of a test map
$D\to C/N$ to $B$. Use that lift to give $C$ its $B$-algebra structure and
then lift $D\to C/N$ over $B$. For uniqueness, any two lifts first agree on
$B$ by formal unramifiedness of $B$ over $A$, and then agree on $D$ by formal
unramifiedness over $B$.

For base change along $A\to A'$, suppose $C$ is an $A'$-algebra. Restriction
and the coproduct property identify maps from $A'\otimes_A B$ over $A'$ with
maps from $B$ over $A$, both for target $C$ and for $C/N$. These identifications
commute with reduction. This proves all three base-change statements.

For the smoothness assertion over $B$, let $C$ be the $B$-algebra in a
square-zero test. A lift over $A$ exists, and its restriction to $B$ equals
the prescribed map $B\to C$ by formal unramifiedness over $A$. It is therefore
a lift over $B$. Uniqueness over $A$ implies uniqueness among maps over $B$.
\end{proof}

\begin{proposition}[Localization and finite principal covers]\label{res:FL-10}
For any multiplicative set $S\subseteq B$, the map $B\to S^{-1}B$ is
formally \'{e}tale. For each formal property in Definition~\ref{res:FL-01},
localizing the target preserves the property. One may also localize the
source at elements already acting as units in the target.

Suppose $B$ is finitely presented over $A$ and $b_1,\ldots,b_s$ generate the
unit ideal of $B$. Then each property holds for $A\to B$ if it holds for all
$A\to B_{b_i}$. For formal unramifiedness the same assertion holds without
finite presentation, and the condition can also be checked on all
$A\to B_{\mathfrak q}$ for $\mathfrak q\in\operatorname{Spec}B$.
\end{proposition}
\begin{proof}
In a square-zero test over $B$, any image of an element of $S$ is a unit
in $C$ if its reduction is a unit in $C/N$. Indeed, multiplying by a
lift of the inverse gives $1+n$ for $n\in N$, with inverse $1-n$.
The universal property of localization then gives exactly one lift.
This also covers a localization at a set containing zero: if the test map
exists, the conclusion about the unit $0$ forces $C$ to be the zero ring.
Target localization follows by Proposition~\ref{res:FL-09}. If $S$
is instead a multiplicative set in $A$ already mapping to units in $B$,
an $S^{-1}A$-algebra test over $B$ is also a test over $A$. Any $A$-algebra
lift automatically respects $S^{-1}A$, since homomorphisms preserve inverses.
Conversely, in any test over $A$ the image of $S$ is invertible after
reduction because the test map factors through $B$. Its image is then
invertible in $C$, giving $C$ the unique required $S^{-1}A$-algebra structure.

For formal unramifiedness, Proposition~\ref{res:FL-04} and
Result~\ref{res:KD-16} reduce the test to vanishing of a module on a cover.
If a module element dies in every $B_{b_i}$, some power of each $b_i$
annihilates it. These powers still generate the unit ideal because no maximal
ideal contains them all. For the prime-local test, the annihilator of any
nonzero element is proper, and a maximal ideal containing it gives a
localization in which that element is still nonzero.

For formal smoothness, fix a finite polynomial presentation with ideal $I$
finitely generated and let $L$, $F$, and $\delta$ be as in
Proposition~\ref{res:FL-06}. Choose lifts $g_i\in P$ of $b_i$. Over $B_{b_i}$
the conormal map is the localization $\delta_{b_i}$: indeed
$P_{g_i}\twoheadrightarrow B_{b_i}$ has localized kernel, localization
commutes with $I/I^2$, and Result~\ref{res:KD-16} compares the middle terms.
The algebra $P_{g_i}$ is formally smooth by the first part and
Proposition~\ref{res:FL-09}. Theorem~\ref{res:FL-07} gives a retraction
of $\delta_{b_i}$ over $B_{b_i}$. As $F$ is finite free, clear the
denominators of its finitely many basis images to find $h_i:F\to L$ whose
localization is a power of $b_i$ times that retraction. The equality
$(h_i\delta)_{b_i}=b_i^{a_i}\operatorname{id}_{L_{b_i}}$ then holds.
Since $L$ is finitely generated, a further power of $b_i$ clears the error on
all its generators. After adjusting $h_i$ and $a_i$ we have
$h_i\delta=b_i^{a_i}\operatorname{id}_{L}$. Write $1=\sum_i c_i b_i^{a_i}$.
Then $\sum_i c_i h_i$ is a global retraction of $\delta$, and the theorem
applies. An empty list of generators of the unit ideal can occur only for
the zero ring, in which case the conormal modules are zero and the same
criterion applies. The formally \'{e}tale case combines the two arguments.
\end{proof}

\begin{theorem}[Formally \'{e}tale intermediate maps]\label{res:FL-11}
The map $A\to B$ is formally \'{e}tale if and only if it is formally smooth
and $\Omega_{B/A}=0$. If $A\to B\to D$ and $B\to D$ is formally \'{e}tale,
there is a canonical isomorphism
\[
  D\otimes_B\Omega_{B/A}\xrightarrow{\ \sim\ }\Omega_{D/A},
  \qquad 1\otimes db\longmapsto d(b_D),
\]
where $b_D$ is the image of $b$ in $D$.
\end{theorem}
\begin{proof}
Bijectivity of the reduction map means existence and uniqueness, and
Proposition~\ref{res:FL-04} identifies uniqueness. For the comparison,
let $M$ be any $D$-module. An $A$-derivation $d_0:B\to M$ defines an
$A$-algebra map $B\to D\ltimes M$ by $b\mapsto(b_D,d_0(b))$.
Give $D\ltimes M$ this $B$-algebra structure and consider
$\operatorname{id}_D$ as a map to $(D\ltimes M)/M$. Its unique lift has
first component $\operatorname{id}_D$ and second component an $A$-derivation
extending $d_0$. Conversely every such extension defines a lift. Hence
restriction induces an isomorphism
$\operatorname{Der}_A(D,M)\cong\operatorname{Der}_A(B,M)$, natural in $M$.
Via the universal property and extension of scalars, restriction corresponds
to precomposing with the displayed canonical map. A map of modules that
induces these bijections for every $M$ is an isomorphism: take $M$ to be its
source to obtain a left inverse and $M$ to be its target to check the right
inverse by injectivity.
\end{proof}

\begin{proposition}[Quotients and two examples]\label{res:FL-12}
For any ideal $I\subseteq A$, the quotient $A\to A/I$ is formally unramified.
It is formally smooth, equivalently formally \'{e}tale, precisely when
$I=I^2$. Over a field $k$ of characteristic $p>0$, the algebra
$k[X]/(X^p)$ has free differential module of rank one but is not formally
smooth over $k$.
\end{proposition}
\begin{proof}
A map over $A$ out of $A/I$ is fixed by the structure map of the target,
which gives uniqueness. For $A\to A/I$, use $P=A$ in
Theorem~\ref{res:FL-07}. The middle conormal term is zero, so its first
map splits injectively exactly when $I/I^2=0$. In particular a nonzero
square-zero ideal does not give a formally smooth quotient.

For $B=k[X]/(X^p)$, let $x$ denote the class of $X$.
Result~\ref{res:KD-21} gives $\Omega_{B/k}=B\,dx$ since $(X^p)'=0$.
But the conormal map kills the nonzero class of $X^p$ in
$(X^p)/(X^{2p})$. Theorem~\ref{res:FL-07} excludes formal smoothness.
\end{proof}

\begin{corollary}[A split sequence for a formally smooth stage]\label{res:FL-13}
For $A\to B\to D$ with $B\to D$ formally smooth, the sequence
\[
  0\to D\otimes_B\Omega_{B/A}\to\Omega_{D/A}
    \to\Omega_{D/B}\to 0
\]
is split exact as a sequence of $D$-modules.
\end{corollary}
\begin{proof}
For any $D$-module $M$ and any $A$-derivation on $B$ into $M$, use the same
square-zero test as in Theorem~\ref{res:FL-11}. Existence gives an extension
to an $A$-derivation on $D$. Take $M=D\otimes_B\Omega_{B/A}$ and start
with the derivation $b\mapsto1\otimes db$. An extension produces a $D$-linear
left inverse to the first arrow. The rest is the transitivity sequence
of Result~\ref{res:KD-09}.
\end{proof}

\section{Thickenings and descent}

\begin{proposition}[Corresponding ideal thickenings]\label{res:FL-14}
Let $A\to B$ be formally \'{e}tale. Let $J\subseteq B$ be an ideal such
that the induced $A\to B/J$ is surjective, and write $I$ for its kernel.
For every $n\ge1$ the canonical map
$A/I^n\to B/J^n$ is an isomorphism. These maps induce an isomorphism of
associated graded rings
$\bigoplus_{n\ge0} I^n/I^{n+1}\cong\bigoplus_{n\ge0}J^n/J^{n+1}$.
\end{proposition}
\begin{proof}
Identify $B/J$ with $A/I$. Apply Lemma~\ref{res:FL-02} to the nilpotent
kernel of $A/I^n\to A/I$. This gives a unique lift
$B\to A/I^n$ of the reduction map to $A/I$. The lift sends $J$ into
$I/I^n$, so it kills $J^n$ and gives a map $B/J^n\to A/I^n$.
Its composition in one order is the identity by $A$-linearity and
surjectivity of $A\to A/I^n$. In the other order, the lift
$B\to A/I^n\to B/J^n$ and the quotient $B\to B/J^n$ both reduce to the same
map to $B/J$. They agree by uniqueness along $B/J^n\to B/J$.
The canonical isomorphisms commute with the transition quotients.
Passing to their kernels identifies each positive graded piece, with
degree zero given by $A/I\cong B/J$. The product is preserved because all
maps are induced by $A\to B$.
\end{proof}

\begin{corollary}[Change of scalars for the diagonal thickenings]
\label{res:FL-15}
Let $A\to B\to D$ have $B\to D$ formally \'{e}tale and put
$J_B=\ker(B\otimes_A B\to B)$, $J_D=\ker(D\otimes_A D\to D)$.
Using the first copy of $B$ for scalar extension, the natural map
\[
 D\otimes_B \bigl((B\otimes_A B)/J_B^{n+1}\bigr)
    \longrightarrow (D\otimes_A D)/J_D^{n+1}
\]
is an isomorphism for $n\ge0$. For $n=1$ this induces the comparison of
Theorem~\ref{res:FL-11} on the square-zero diagonal ideals.
\end{corollary}
\begin{proof}
The map $D\otimes_A B\to D\otimes_A D$ applying $B\to D$ on the second
factor is a base change of that map, and so is formally \'{e}tale by
Proposition~\ref{res:FL-09}. Its composite with multiplication gives
the surjective evaluation $D\otimes_A B\to D$, $d\otimes b\mapsto d\,b_D$.
The kernel is the extension of $J_B$: quotienting by this extension
is the base change of $(B\otimes_A B)/J_B=B$ along $B\to D$ on the first
factor. Powers of extended ideals are extensions of powers. Apply
Proposition~\ref{res:FL-14} to $D\otimes_A B\to D\otimes_A D$ and the ideal
$J_D$. For $n=1$, Result~\ref{res:KD-03} identifies the first diagonal
ideals with differentials. The multiplication quotients onto $B$ and $D$
split linearly by their first factors, so identifying the kernels here
requires no flatness. The comparison takes the class
$[1\otimes b-b\otimes1]$ to $[1\otimes b_D-b_D\otimes1]$.
\end{proof}

\begin{proposition}[Filtered colimits and uniqueness]\label{res:FL-16}
A filtered colimit of formally \'{e}tale $A$-algebras is formally
\'{e}tale over $A$. A filtered colimit of formally unramified $A$-algebras
is formally unramified over $A$.
\end{proposition}
\begin{proof}
For a test map from the colimit to $C/N$, restrict to each algebra in the
diagram. In the formally \'{e}tale case the unique component lifts are
compatible with the transition maps, again by uniqueness, and so give
the required map out of the colimit. Any two lifts of the colimit map
agree whenever their component restrictions agree. This proves uniqueness
in both cases.
\end{proof}

\begin{proposition}[Faithfully flat descent in finite presentation]
\label{res:FL-17}
Let $A'$ be faithfully flat over $A$ and write $B'=A'\otimes_A B$.
Formal unramifiedness descends from $A'\to B'$ to $A\to B$. If $B$ is
finitely presented over $A$, formal smoothness and formal \'{e}taleness
also descend.
\end{proposition}
\begin{proof}
For formal unramifiedness, use Proposition~\ref{res:FL-04},
Result~\ref{res:KD-17}, and faithfulness of extension of scalars from
$B$ to $B'$. For smoothness choose a finite polynomial presentation with
$I=(f_1,\ldots,f_m)$ and conormal map $\delta:L\to F$, where $F=B^n$.
Flatness identifies the new kernel with $A'\otimes_A I$. It also identifies
its square with the extension of $I^2$, so the new conormal map is
$B'\otimes_B\delta$, by the comparison on $df_j$ in Result~\ref{res:KD-17}.

A linear map $F\to L$ is specified by its $n$ basis images. The map
\[
  L^n\longrightarrow L^m,\qquad
  h\longmapsto\bigl(h(\delta([f_j]))\bigr)_{j=1}^m
\]
is linear and has a preimage of $\bigl([f_j]\bigr)_j$ exactly when
$\delta$ has a retraction, since the classes $[f_j]$ generate $L$.
After extending to $B'$ the tuple lies in the image by
Theorem~\ref{res:FL-07}. Its class in the cokernel vanishes already over $B$:
if it were nonzero, the cyclic module it generates would remain nonzero by
faithful flatness, and would still inject into the base-changed cokernel by
flatness. Thus the required retraction exists over $B$. Applying the
splitting criterion proves formal smoothness, and the formally \'{e}tale
claim follows by combining existence and uniqueness.
\end{proof}

\begin{proposition}[Tests with nil ideals in finite presentation]
\label{res:FL-18}
If $B$ is finitely presented over $A$ and formally smooth, maps from $B$
lift even through ideals all of whose elements are nilpotent. If $B$ is
of finite type and formally unramified, lifts are unique through such ideals.
In particular a finitely presented formally \'{e}tale algebra has both
properties for nil ideals.
\end{proposition}
\begin{proof}
Use a presentation $A[X_1,\ldots,X_n]/(f_1,\ldots,f_m)$ for existence.
Given $u:B\to C/N$ with $N$ a nil ideal, choose representatives $c_i$
of $u(\bar X_i)$. The finitely many elements $f_j(c_1,\ldots,c_n)$ lie in
$N$, so their generated ideal $N_0$ is nilpotent. Indeed, if the
generators have nilpotence exponents $e_j$, every monomial of degree
$1+\sum_j(e_j-1)$ contains a vanishing power. The representatives define
a map $B\to C/N_0$ reducing to $u$, and Lemma~\ref{res:FL-02} lifts it to $C$.
For uniqueness, two lifts agreeing modulo $N$ differ on finitely many
generators by elements of $N$. The ideal generated by those differences
is nilpotent by the same argument. The lifts agree modulo that ideal and
hence agree in $C$ by Lemma~\ref{res:FL-02}.
\end{proof}

The next chapter imposes finite presentation on formal smoothness. The
splitting of the conormal map then has finite data, which can be studied by
Jacobian matrices on principal open subsets.
\chapter{Smooth Algebras and the Jacobian Criterion}
\label{ch:smooth}

Let $A$ be a commutative ring. The lifting criterion of
Chapter~\ref{ch:lifting} has no finiteness hypothesis. In this chapter
finite presentation lets us turn that criterion into a finite
Jacobian calculation near each prime of an $A$-algebra.

\section{The finite-presentation condition}

\begin{definition}[Smoothness]
\label{res:SM-01}
An $A$-algebra $B$ is \emph{smooth} if it is of finite presentation
over $A$ and is formally smooth over $A$. Finite presentation means
that $B$ admits a presentation
\[
  B \simeq A[X_1,\ldots,X_n]/(f_1,\ldots,f_c)
\]
with $n$ and $c$ finite. The definition requires the square-zero
existence property of Chapter~\ref{ch:lifting} for every test
$A$-algebra and every square-zero ideal.
\end{definition}

\begin{theorem}[Conormal criterion]
\label{res:SM-02}
Suppose $B$ is of finite presentation over $A$ and
$P=A[X_1,\ldots,X_n]\twoheadrightarrow B$ has kernel $I$.
Then $B$ is smooth over $A$ if and only if the map
\[
  d_I\colon I/I^2\longrightarrow B\otimes_P\Omega_{P/A},
  \qquad [f]\longmapsto 1\otimes df ,
\]
admits a $B$-linear left inverse. In that case
\[
  0\longrightarrow I/I^2\xrightarrow{d_I}
  B\otimes_P\Omega_{P/A}\longrightarrow\Omega_{B/A}
  \longrightarrow0
\]
is split exact, and both $I/I^2$ and $\Omega_{B/A}$ are finite
projective over $B$. Equivalently, $d_I$ is injective and
$\Omega_{B/A}$ is projective.
\end{theorem}

\begin{proof}
The square-zero criterion of Theorem~\ref{res:FL-07} gives the
equivalence with a left inverse. The remaining maps are those of
the conormal sequence in \ref{res:KD-11}. By \ref{res:KD-07}, its
middle term is free of rank $n$. A retraction splits that free
module into two direct summands, which are finite projective. In
the other direction, injectivity gives short exactness and
projectivity of its quotient gives a section of the last map.
Projecting onto the first summand then gives a left inverse of
$d_I$. The argument applies to every finite polynomial
presentation, as it characterizes the same lifting property.
\end{proof}

\section{Jacobian presentations}

\begin{definition}[Standard smooth presentation]
\label{res:SM-03}
A presentation $B=P/I$ as in \ref{res:SM-01} with $0\le c\le n$
is \emph{standard smooth} if a $c\times c$ minor of the matrix
$(\partial f_j/\partial X_i)_{1\le i\le n,\,1\le j\le c}$
has unit image in $B$. The minor uses all $c$ columns. After
reordering the variables one can take its row indices to be
$1,\ldots,c$. For $c=0$ its determinant is $1$. An algebra is
standard smooth if it has such a presentation.
\end{definition}

\begin{proposition}[The invertible-minor criterion]
\label{res:SM-04}
A standard smooth presentation defines a smooth algebra.
With the row indices ordered as in \ref{res:SM-03}, $I/I^2$
has basis $[f_1],\ldots,[f_c]$, and $\Omega_{B/A}$ has basis
$d\overline X_{c+1},\ldots,d\overline X_n$.
\end{proposition}

\begin{proof}
Write $B^n=B^c\oplus B^{n-c}$ using the ordered differentials.
There is a surjection $u\colon B^c\to I/I^2$ sending the standard
basis to the relation classes. By \ref{res:KD-12}, the composite
\[
  B^c\xrightarrow{u} I/I^2\xrightarrow{d_I} B^n
       \xrightarrow{\operatorname{pr}_1} B^c
\]
is the invertible block. Hence $u$ is injective and therefore an
isomorphism, while the composition of $\operatorname{pr}_1$
with the inverse of $\operatorname{pr}_1 d_I$ retracts $d_I$.
The image of $d_I$ is the graph of a map $L\colon B^c\to B^{n-c}$.
The quotient by this graph is identified with $B^{n-c}$ by
$(v,w)\mapsto w-Lv$. This gives the claimed differential basis
and proves smoothness by \ref{res:SM-02}.
\end{proof}

\begin{theorem}[Local standard presentations]
\label{res:SM-05}
Let $B$ be smooth over $A$. For every prime $\mathfrak q\subset B$
there is a $g\in B\setminus\mathfrak q$ such that $B_g$ is standard
smooth over $A$. Finitely many such $g$ may be chosen to generate
the unit ideal.
\end{theorem}

\begin{proof}
Choose $P=A[X_1,\ldots,X_n]\twoheadrightarrow B$ with finitely
generated kernel $I$ and let $\mathfrak Q$ be the inverse image
of $\mathfrak q$ in $P$. Write $k=\kappa(\mathfrak q)$.
The split inclusion in \ref{res:SM-02} is still injective after
tensoring with $k$. Choose $f_1,\ldots,f_c\in I$ whose classes
give a $k$-basis of $(I/I^2)\otimes_B k$. This is possible because
the image of $I$ spans that vector space. The columns $df_j$ are
independent in $k^n$, so some minor $\Delta\in P$ avoids
$\mathfrak Q$.

By Nakayama's lemma \cite[Theorem ALG-05-C04-R01]{p103}, the
chosen classes generate $(I/I^2)_{\mathfrak q}$. Consequently
\[
  I_{\mathfrak Q}
    =(f_1,\ldots,f_c)P_{\mathfrak Q}+I_{\mathfrak Q}^2.
\]
Apply Nakayama again to
$I_{\mathfrak Q}/(f_1,\ldots,f_c)P_{\mathfrak Q}$, a finite
$P_{\mathfrak Q}$-module, using
$I_{\mathfrak Q}\subseteq\mathfrak QP_{\mathfrak Q}$. We get
$I_{\mathfrak Q}=(f_1,\ldots,f_c)P_{\mathfrak Q}$. Since $I$
has finitely many generators, clearing denominators in these
generator expressions produces $h\notin\mathfrak Q$ with
$I_h=(f_1,\ldots,f_c)P_h$.

Set $s=h\Delta$ and let $g=\overline s$ in $B$. Then $g$ avoids
$\mathfrak q$ and
\[
  B_g\simeq P[T]/(f_1,\ldots,f_c,sT-1).
\]
The new variable adjoins the inverse of $s$. In the Jacobian use
the old $c$ rows and the $T$ row, and include the column for
$sT-1$. The rows and columns so selected form a block triangular
matrix, with blocks having determinants $\Delta$ and $s$.
Their product is a unit in $B_g$ because $s=h\Delta$ is a unit
there. This is a standard smooth presentation, including $c=0$.
Finally the ideal generated by all chosen $g$ is contained in no
prime and hence is the unit ideal. An expression for $1$ uses
only finitely many of them.
\end{proof}

\section{Permanence and the smooth locus}

\begin{proposition}[Composition, base change and localization]
\label{res:SM-06}
Smooth algebras are preserved by composition and arbitrary base
change. If $B$ is smooth over $A$ and $g\in B$, then $B_g$
is smooth over $A$. If $a\in A$ is already a unit in $B$, the
induced map $A_a\to B$ is smooth as well.
\end{proposition}

\begin{proof}
For composition, take finite presentations for the two maps and
lift the coefficients of the second set of equations to the first
polynomial ring. The resulting combined variables and equations
present the composite. The lifting property is preserved because
one lifts through the middle algebra first and then uses that lift
as the middle-algebra structure on the test ring
(\ref{res:FL-09}). Base change transports a finite presentation
coefficientwise and preserves the lifting property by the
pushout argument in \ref{res:FL-09}. Principal localization
$B\to B_g$ has one variable and one equation $gT-1$. Its lifts
exist because an element invertible modulo a square-zero ideal
is already invertible, so composition applies. The last assertion
is base change along $A\to A_a$, since $A_a\otimes_A B\simeq B$.
\end{proof}

\begin{theorem}[Locality on the target]
\label{res:SM-07}
Let $B$ be any $A$-algebra. If $g_1,\ldots,g_m\in B$ generate the
unit ideal and all $B_{g_i}$ are smooth over $A$, then $B$ is smooth
over $A$. Define $B$ to be smooth over $A$ at a prime
$\mathfrak q\subset B$ if $B_g$ is smooth for some
$g\notin\mathfrak q$. If $B$ is of finite presentation, the
following conditions are equivalent:
\begin{enumerate}
\item $B$ is smooth over $A$ at $\mathfrak q$.
\item $B_{\mathfrak q}$ is formally smooth over $A$.
\item The localized conormal map of a finite polynomial
  presentation has a left inverse over $B_{\mathfrak q}$.
\end{enumerate}
In particular, the smooth locus is open. Smoothness is equivalent
to smoothness at every prime, and to admitting a finite cover by
standard smooth principal localizations.
\end{theorem}

\begin{proof}
We first descend finite presentation. Choose coefficients $b_i\in B$
with $\sum_i b_i g_i=1$. Let $C$ be the finite type $A$-subalgebra
of $B$ generated by the $g_i$, the $b_i$ and numerators of finite
algebra generating sets for $B_{g_i}$. Each map
$C_{g_i}\to B_{g_i}$ is injective by the localization equality
test and surjective by construction. If $v\in B$, for each $i$ some
power of $g_i$ times $v$ belongs to $C$. These powers generate
the unit ideal in $C$: a maximal ideal containing all the powers
would contain all the $g_i$. Thus $v\in C$, and $B=C$ has finite
algebra generators.

Take a finite polynomial presentation $P\twoheadrightarrow B$
with kernel $I$ for which we have yet to prove finite generation
of $I$. Lift $g_i$ to $h_i\in P$. The ideal $I_{h_i}$ is
finitely generated. Indeed, for any surjection from a finite
polynomial algebra onto a finitely presented algebra, the kernel
is finitely generated. To see this, use the known finite
presentation on variables $Y$, add the given variables $X$
with equations expressing $X$ in terms of $Y$, and adjoin
equations expressing $Y$ in terms of $X$. These last equations
already hold in the quotient. Eliminating $Y$ leaves finitely
many relations in $X$. Apply this argument to
$P[T]\twoheadrightarrow B_{g_i}$ with $T$ mapping to $g_i^{-1}$,
then pass to $P_{h_i}=P[T]/(h_iT-1)$.

Choose finitely many elements of $I$ whose images generate
each $I_{h_i}$ by clearing their denominators. Let $J$ be the
ideal they generate together with $1-\sum_i \widetilde b_i h_i$,
where $\widetilde b_i\in P$ lifts $b_i$. The primes of $P$ missing
an $h_i$ see $I=J$ after localization. At a prime containing all
$h_i$, the added generator is a unit after localization, so
again the two localized ideals agree. Thus $I/J=0$, since a
nonzero cyclic submodule survives at a maximal ideal containing
its annihilator. This proves finite presentation of $B$.

Now Proposition~\ref{res:FL-10} applies to the formally smooth
algebras $B_{g_i}$ and the finitely presented algebra $B$.
It gives formal smoothness over $A$, so $B$ is smooth.

For the pointwise statements, localization at $\mathfrak q$
preserves formal smoothness by the unit argument of
\ref{res:FL-10}. In the local conormal sequence use
$P_{\mathfrak Q}\twoheadrightarrow B_{\mathfrak q}$, with
$\mathfrak Q$ the inverse image of $\mathfrak q$. By
\ref{res:KD-16} the middle module and its map are obtained by
localizing those for $P$. Since $P_{\mathfrak Q}$ is formally
smooth over $A$, \ref{res:FL-07} gives the equivalence of the
formal lifting condition and a retraction there. A local
retraction has finitely many denominators on the free basis,
all outside $\mathfrak q$. Its retraction identity can be
cleared on finitely many generators of $I/I^2$. Inverting the
product $g$ of these elements gives an actual retraction over
$B_g$. This algebra is finitely presented over $A$ by adjoining
one inverse, so it is smooth by \ref{res:FL-07}.

The sets $D(g)$ in the definition are open and cover the smooth
locus. A cover of the whole spectrum by such $D(g)$ contains a
finite subcover, as an expression for $1$ in the ideal they
generate has finite support. This proves the remaining equivalences
using \ref{res:SM-05} and the locality already proved.
\end{proof}

\section{Flatness and relative dimension}

\begin{theorem}[Flatness of smooth algebras]
\label{res:SM-08}
Every smooth $A$-algebra is flat as an $A$-module.
\end{theorem}

\begin{proof}
Suppose first that $A$ is Noetherian. Write $B=P/I$ for a finite
polynomial algebra $P$ over $A$. For every $j\ge 1$, the kernel
of $P/I^{j+1}\to P/I^j$ is square-zero because $2j\ge j+1$.
Start with the identity $B\to P/I$. Successive applications of
formal smoothness give compatible $A$-algebra maps
$B\to P/I^j$. Their limit is a section of
$\widehat P\to P/I=B$. The ring $P$ is Noetherian by the
Hilbert basis theorem \cite[Theorem ALG-05-C02-R01]{p103}.
The completion $\widehat P$ is flat over $P$ by
\cite[Theorem ALG-05-C04-S01]{p103}. Since $P$ is free over $A$,
$\widehat P$ is flat over $A$: tensoring over $A$ is the
composition of the two exact scalar-extension functors.
The section exhibits $B$ as a direct summand of $\widehat P$
as an $A$-module. A direct summand of a flat module is flat,
as can be checked on the tensor of each injection.

For a standard smooth presentation over arbitrary $A$, lift the
unit condition on the chosen minor $\Delta$ to a polynomial
identity $u\Delta-1=\sum_j v_j f_j$ in $P$.
Let $A_0\subseteq A$ be the subring generated over the image
of $\mathbb Z$ by the coefficients of the $f_j$, $u$ and $v_j$.
The same identity holds over $A_0$, since its inclusion into
$A$ is injective. Thus the same equations give a standard
smooth $A_0$-algebra $B_0$ with $B\simeq A\otimes_{A_0}B_0$.
The ring $A_0$ is Noetherian, being a quotient of a finite
polynomial algebra over $\mathbb Z$. The first part applied
to $B_0$ and \ref{res:SM-04} give flatness. Arbitrary base
change of a flat algebra is flat because for each $A$-module
$M$ the tensor $M\otimes_A(A\otimes_{A_0}B_0)$ is naturally
$M\otimes_{A_0}B_0$.

For general smooth $B$, choose the cover of \ref{res:SM-05}.
To test an injection of $A$-modules after tensoring with $B$,
its kernel is a $B$-module and vanishes after localization
at each of the covering elements, by the flatness already
proved for $B_g$. If an element of this kernel is killed by
a power of every covering element, these powers generate
the unit ideal, so the element itself is zero.
\end{proof}

\begin{theorem}[Relative rank]
\label{res:SM-09}
For $B$ smooth over $A$, the function
\[
  r(\mathfrak q)=
  \dim_{\kappa(\mathfrak q)}
    \bigl(\kappa(\mathfrak q)\otimes_B\Omega_{B/A}\bigr)
\]
is locally constant on $\operatorname{Spec}B$. On a standard
smooth chart with $n$ variables and $c$ relations it is $n-c$.
It is preserved by base change at corresponding primes.
If a smooth $A$-algebra $B$ is followed by a smooth $B$-algebra
$C$, the relative ranks add at a prime of $C$, using its
contraction for the rank of $B$ over $A$.
\end{theorem}

\begin{proof}
The charts of \ref{res:SM-05} and the bases of \ref{res:SM-04}
give the formula and local constancy. For a prime upstairs
after base change, use \ref{res:KD-17} and identify the new
cotangent fiber with the old one tensored along the extension
of residue fields. For composition, apply the split
transitivity sequence of \ref{res:FL-13} to $A\to B\to C$.
Tensor the split sequence with the residue field at the prime
of $C$. The left vector space is obtained from the cotangent
fiber for $B/A$ by the extension of residue fields. Dimensions
of the two end terms sum to the dimension of the middle term.
\end{proof}

\section{The rank test at a specified prime}

\begin{proposition}[Jacobian criterion with local relations]
\label{res:SM-10}
Let $B=P/I$ be finitely presented over $A$ with
$P=A[X_1,\ldots,X_n]$. Let $\mathfrak Q$ be a prime of $P$
containing $I$, set $\mathfrak q=\mathfrak Q/I$, and choose
$f_1,\ldots,f_c\in I$ whose classes are a basis of
$I_{\mathfrak Q}/\mathfrak Q I_{\mathfrak Q}$
over $\kappa(\mathfrak q)$. Then $B$ is smooth at $\mathfrak q$
if and only if the matrix of the partial derivatives of these
$f_j$ has rank $c$ over $\kappa(\mathfrak q)$. In particular
the full column-rank test applies to a prescribed list of
equations at primes where the list is a minimal generating
set of $I_{\mathfrak Q}$.
\end{proposition}

\begin{proof}
The source of the conormal map after tensoring with
$\kappa(\mathfrak q)$ identifies with
$I_{\mathfrak Q}/\mathfrak Q I_{\mathfrak Q}$, since
$I_{\mathfrak Q}^2\subseteq\mathfrak Q I_{\mathfrak Q}$.
A splitting at the prime gives injectivity of this map over
the field and hence the stated rank. For the converse, full
rank supplies a minor avoiding $\mathfrak Q$. Nakayama gives
that the $f_j$ generate $I_{\mathfrak Q}$. The argument of
\ref{res:SM-05} that clears denominators, inverts the minor
and adds its inverse now gives a standard smooth presentation
on a neighborhood. Hence $B$ is smooth there.
\end{proof}

\begin{proposition}[Permanence of standard presentations]
\label{res:SM-11}
Standard smoothness is preserved by composition, arbitrary
base change and target principal localization. Under those
operations the number $n-c$ adds for composition and is
unchanged for base change and localization in the displayed
presentations.
\end{proposition}

\begin{proof}
The combined presentation for composition used in
\ref{res:SM-06} gives a block triangular Jacobian minor,
because the old equations are independent of the new variables.
Its diagonal blocks are the two chosen minors, which map to
units. For base change the coefficients and the determinant are
sent to their images. For localization, choose a polynomial
$s$ lifting the element to invert and adjoin $sT-1$, as in
\ref{res:SM-05}. The chosen determinant is multiplied by $s$,
which becomes a unit. The dimension counts follow from adding
the numbers of variables and equations, including exactly one
of each for principal localization.
\end{proof}

\begin{proposition}[Dimension of a standard smooth fiber]
\label{res:SM-12}
If $B$ has a standard smooth presentation over $A$ with
$n$ variables and $c$ relations, each nonzero field fiber
$F\otimes_A B$ has Krull dimension $n-c$, for any ring map
$A\to F$ to a field. After passing to an algebraic closure of
$F$, all local rings at closed points of that fiber are
regular local rings of dimension $n-c$.
\end{proposition}

\begin{proof}
By \ref{res:SM-11} the fiber is standard smooth over $F$,
and the same holds over an algebraic closure $L$. The map
$F\otimes_A B\to L\otimes_A B$ is injective and integral.
Indeed it is a scalar extension by a field, and every element
uses finitely many elements algebraic over $F$, hence lies
in an extension which is finite as a module. For integral
inclusions Krull dimensions agree by lying over, going up
and incomparability \cite[Theorems ALG-05-C01-S04 and
ALG-05-C01-R04]{p103}. Thus it suffices to study the
fiber over $L$.

At a maximal ideal of that fiber, the variables have values
$a_i\in L$ by the weak Nullstellensatz
\cite[Theorem ALG-05-C01-S05]{p103}. The local polynomial
ring $T=L[X_1,\ldots,X_n]_{(X_1-a_1,\ldots,X_n-a_n)}$ is
regular local of dimension $n$: the coordinate prime chain
gives the lower bound and the dimension of the polynomial
ring the upper bound, and the $n$ coordinate differences
generate the maximal ideal. Modulo the square of that ideal,
the equations $f_j$ have linear parts
$\sum_i(\partial f_j/\partial X_i)(a)(X_i-a_i)$.
The chosen invertible minor makes them independent. Complete
their classes to a basis and choose lifts of the extra basis
vectors. This is a minimal generating set of the maximal
ideal by Nakayama. After quotient by the first generator,
the remaining $n-1$ generators bound the dimension above
by $n-1$, and the system-of-parameters bound gives dimension
at least $n-1$. Hence the quotient is regular local and the
remaining elements are a minimal generating set. Repeat for
the first $c$ generators
\cite[Theorems ALG-05-C03-S01 and ALG-05-C03-S03]{p103}.
The resulting quotient is the local ring of the fiber at the
chosen closed point. Every prime chain is contained in a
maximal ideal, so taking the supremum of the local dimensions
over maximal ideals gives the dimension asserted.
\end{proof}

\section{Further consequences of locality}

\begin{proposition}[The smooth locus under flat base change]
\label{res:SM-13}
Let $B$ be finitely presented over $A$. If $A\to A'$ is flat
and $B'=A'\otimes_A B$, the smooth locus of $B'/A'$ is
the inverse image of the smooth locus of $B/A$
under $\operatorname{Spec}B'\to\operatorname{Spec}B$.
\end{proposition}

\begin{proof}
If $B$ is smooth on $D(g)$, its base change is smooth on the
inverse image by \ref{res:SM-06}. Conversely, take a prime
$\mathfrak q'$ of $B'$ at which $B'/A'$ is smooth, and let
$\mathfrak q$ be its contraction to $B$. The local map
$B_{\mathfrak q}\to B'_{\mathfrak q'}$ is flat by base change
and localization, and faithfully flat because it is local
\cite[Theorems ALG-03-C02-S04 and ALG-05-C01-R04]{p103}.
Choose a finite polynomial presentation $B=P/I$.
Flatness of $A'$ over $A$ identifies the extended ideal and
its square with $A'\otimes_A I$ and $A'\otimes_A I^2$ in
$A'\otimes_A P$: their images are the extension ideals, and
products of images extend the products in $I^2$. Therefore
the conormal sequence upstairs is obtained by tensoring the
one downstairs with $B'$. The middle term has the same
polynomial basis after base change and the last term agrees
by \ref{res:KD-17}.

Now put $R=B_{\mathfrak q}$, $S=B'_{\mathfrak q'}$ and write
the localized conormal sequence over $R$ as
$M\xrightarrow{d}F\xrightarrow{\pi}N\to0$.
Its base change to $S$ is split short exact by
\ref{res:SM-07}. Faithful flatness detects injectivity of $d$.
To descend a section of $\pi$, note that $N$ is finitely
presented by \ref{res:KD-13}. For $E=F,N$ the natural maps
\[
  S\otimes_R\operatorname{Hom}_R(N,E)
    \simeq \operatorname{Hom}_S(S\otimes_R N,S\otimes_R E)
\]
are isomorphisms: choose a finite free presentation of $N$,
write each Hom as a kernel of a map between finite powers,
and use flatness to preserve that kernel. Postcomposition with
$\pi$ is compatible with these identifications. The class
of $\operatorname{id}_N$ in its cokernel becomes zero over
$S$, so it was already zero by faithful flatness. Thus $\pi$
has a section, and \ref{res:SM-07} gives smoothness at
$\mathfrak q$.
\end{proof}

\begin{corollary}[Finite products]
\label{res:SM-14}
A finite product of $A$-algebras is smooth if and only if its
factors are smooth. The empty product is the zero algebra,
which is smooth.
\end{corollary}

\begin{proof}
For a nonempty product, the coordinate idempotents generate
the unit ideal. Inverting a coordinate idempotent gives the
corresponding factor, since that idempotent becomes $1$ and
annihilates the other coordinates. Apply \ref{res:SM-06} in
one direction and \ref{res:SM-07} in the other. The zero
algebra has the standard smooth presentation $A[T]/(0T-1)$,
where the derivative has unit image in the zero ring.
\end{proof}

\begin{proposition}[Finite coefficient models and lifting charts]
\label{res:SM-15}
If $B$ is smooth over $A$, there is a finitely generated
$\mathbb Z$-subalgebra $A_0\subseteq A$ and a smooth $A_0$-algebra
$B_0$ such that $B\simeq A\otimes_{A_0}B_0$.
If $J\subseteq A$ is an ideal and $C$ is smooth over $A/J$,
there is a finite standard smooth principal cover of $C$
whose charts are reductions modulo $J$ of standard smooth
$A$-algebras.
\end{proposition}

\begin{proof}
For the first assertion take a finite presentation
$B=A[X_1,\ldots,X_n]/I$ with $I=(f_1,\ldots,f_c)$ and a
section $B\to P/I^2$ from \ref{res:FL-07}. Choose polynomial
lifts $h_i$ of the images of the variables. The section
identity and the fact that it respects the relations give
$h_i-X_i\in I$ and $f_j(h_1,\ldots,h_n)\in I^2$.
Express these finitely many inclusions by polynomial
combinations of the $f_j$ and products $f_jf_\ell$.
Take $A_0$ to be generated by the coefficients of the equations,
the $h_i$ and all those combinations. The identities remain
true over $A_0$, so the same $h_i$ give a section modulo the
square of the relation ideal for the resulting quotient $B_0$.
Apply \ref{res:FL-07} and finite presentation over $A_0$.
Base change of the quotient presentation returns $B$.

For the second assertion choose the cover of \ref{res:SM-05}.
For one chart lift its defining polynomials from $A/J$ to $A$
and let $\Delta$ be the chosen minor of the lifted presentation.
Adjoin a variable $T$ and the equation $\Delta T-1$. The
enlarged minor is $\Delta^2$, hence a unit in the quotient,
which is standard smooth over $A$. Reducing modulo $J$
recovers the original chart because its minor was already a
unit there. This works independently on each chart.
\end{proof}

\section{Examples and failure tests}

\begin{example}[Free variables and a hyperbola]
\label{res:SM-16}
Polynomial algebras in finitely many variables are standard
smooth with no equations. The algebra $A[T]/(aT-1)$ is
standard smooth of rank zero, since the derivative $a$ is a
unit in the quotient. For $b\in A$ put
$B=A[X,Y]/(XY-b)$. On $D(\overline X)$ and $D(\overline Y)$
one derivative of $XY-b$ is a unit, so the presentation after
inverting that element is standard smooth of rank one
by the localization construction in \ref{res:SM-05}.
For $A=k$ a field and $b=0$, smoothness fails at the origin.
There the single relation $XY$ minimally generates the
localized nonzero ideal, and both partial derivatives
vanish modulo the maximal ideal, so \ref{res:SM-10} applies.
\end{example}

\begin{example}[A free differential module on a nonsmooth algebra]
\label{res:SM-17}
Let $k$ be a field of characteristic $p>0$ and
$B=k[X]/(X^p)$. Formula \ref{res:KD-21} gives
$\Omega_{B/k}\simeq B\,d\overline X$. Nevertheless the
identity $B\to k[X]/(X^p)$ has no $k$-algebra lift through
$k[X]/(X^{p+1})\to k[X]/(X^p)$, whose kernel is square-zero.
Every candidate image of $\overline X$ is
$X+\lambda X^p$ modulo $X^{p+1}$ with $\lambda\in k$.
Its $p$th power is $X^p+\lambda^p X^{p^2}=X^p\ne0$ in
the lifting ring. This contradicts the defining relation
of $B$. The same failure was detected by the conormal
sequence in \ref{res:FL-12}.
\end{example}

\begin{remark}
The local Jacobian test uses equations which generate the
localized ideal and form a basis modulo its maximal ideal.
Redundant equations can make the full matrix lose rank even
for a smooth algebra. For example, $(X,X^2)=(X)$ defines
$A[X]/(X)\simeq A$. The two columns from this redundant
presentation are no test for the number of independent
relations. The pointwise criterion uses one relation here.
\end{remark}

The general form of the smooth fiber criterion and its
application to geometrically regular fibers belong to the
scheme-theoretic continuation. For the next chapter, we
will use the conormal criterion, flatness and the rank-zero
case of these Jacobian presentations.

For the comparison with the two-term cotangent complex
and for the original finite-presentation formulations,
see \cite[Sections 10.137 and 10.138]{stacks}, especially
\href{https://stacks.math.columbia.edu/tag/00TN}{Tag 00TN}.
Eisenbud's Jacobian matrix for computing differentials is
in \cite[Section 16.1]{eisenbud}.
\chapter{\'Etale Algebras}\label{ch:etale}

Formal \'{e}taleness asks for one lift through each square-zero extension.
Finite presentation gives a polynomial description of this lifting property
and allows one to test it on a cover by principal localizations. We first
derive a square Jacobian criterion, then study local one-variable
presentations and the fibers over fields.

\section{Finite presentation and the square Jacobian}

\begin{definition}[\'Etale and unramified algebras]\label{res:ET-01}
An $A$-algebra $B$ is \emph{\'etale} if it is finitely presented and
formally \'{e}tale over $A$. It is \emph{unramified} if it is of finite
type and formally unramified. We say \emph{finitely presented unramified}
when the stronger finiteness condition is intended.
\end{definition}

\begin{proposition}[Differential characterization]\label{res:ET-02}
The following are equivalent:
\begin{enumerate}
\item $B$ is \'{e}tale over $A$;
\item $B$ is smooth over $A$ and $\Omega_{B/A}=0$;
\item $B$ is finitely presented and, for one (equivalently every)
      polynomial presentation $P\twoheadrightarrow B$ with kernel $I$,
      the conormal map
      $I/I^2\to B\otimes_P\Omega_{P/A}$ is an isomorphism.
\end{enumerate}
\end{proposition}
\begin{proof}
Unique existence of lifts is existence together with uniqueness.
By \ref{res:FL-04}, uniqueness is equivalent to $\Omega_{B/A}=0$.
With finite presentation, existence is smoothness. For a polynomial
presentation, \ref{res:FL-07} expresses existence as split injectivity
of the conormal map. Its cokernel is $\Omega_{B/A}$ by \ref{res:KD-11}.
Thus split injectivity together with trivial cokernel is precisely
invertibility. These properties were derived from the lifting condition
and consequently hold with any polynomial presentation.
\end{proof}

\begin{theorem}[Square Jacobian criterion]\label{res:ET-03}
$B$ is \'{e}tale over $A$ if and only if there is an $A$-algebra
presentation
\[
  B\simeq A[X_1,\ldots,X_n]/(f_1,\ldots,f_n)
\]
in which the image of the determinant of
$(\partial f_j/\partial X_i)_{i,j}$ is a unit of $B$.
The determinant of the empty matrix is $1$.
\end{theorem}
\begin{proof}
For the forward implication start with a finite polynomial presentation
$P\twoheadrightarrow B$ using $n$ variables and with finitely generated
kernel $I$. By \ref{res:ET-02}, the conormal map identifies $I/I^2$
with $B^n$. Lift the preimages of the standard basis to $f_1,\ldots,f_n$
in $I$. Put $J=(f_1,\ldots,f_n)$. Then $I=J+I^2$, so for the
finitely generated $P$-module $N=I/J$ we have $N=IN$.
If $N$ is generated by $u_1,\ldots,u_r$, choose $a_{ij}\in I$ with
$u_i=\sum_j a_{ij}u_j$. Multiplication by the adjugate of
$(\delta_{ij}-a_{ij})$ gives $sN=0$ for
$s=\det(\delta_{ij}-a_{ij})\in1+I$. If $N=0$ one may take $s=1$.
In particular $I P_s=J P_s$, while $s$ has image $1$ in $B$.
Consequently
\[
 B\simeq P_s/JP_s
   \simeq A[X_1,\ldots,X_n,T]/(f_1,\ldots,f_n,sT-1).
\]
The upper left block of its Jacobian is invertible in $B$ by the choice
of the $f_j$. The last row is zero in the first $n$ entries, with last
entry $s$. Hence the full determinant is a unit.

Conversely suppose the displayed Jacobian in the statement is invertible.
The map $B^n\twoheadrightarrow I/I^2$ sending basis vectors to the
relation classes is surjective. Its composite with the conormal map
is the stated Jacobian by \ref{res:KD-12}. Since the composite is
injective, the surjection is an isomorphism. The conormal map itself
is then an isomorphism. Apply \ref{res:ET-02}.
\end{proof}

\begin{proposition}[Permanence and the \'etale locus]\label{res:ET-04}
\'Etaleness is preserved by arbitrary base change, composition, and
principal localization in the target. For any $A$-algebra $B$, if
$b_1,\ldots,b_r$ generate the unit ideal and each $B_{b_i}$ is
\'etale over $A$, then $B$ is \'{e}tale over $A$. If $B$ is finitely
presented, the set of primes $\mathfrak q$ admitting such an
\'etale principal neighborhood is open and equals
\[
 \{\text{smooth primes of }B/A\}\setminus
 \operatorname{Supp}_B(\Omega_{B/A}).
\]
Here the support of a module $M$ is
$\{\mathfrak q:M_{\mathfrak q}\ne0\}$.
\end{proposition}
\begin{proof}
Finite presentation is preserved under base change and composition.
For base change, apply the structural map to the coefficients of a
finite list of defining relations. For composition, take generators of
the two presentations together, and lift each coefficient from the
intermediate algebra to its polynomial presentation. Quotienting by
the first list of relations and then the second gives the composite
algebra. The lifting conditions under these operations follow from
\ref{res:FL-09}. The algebra $B_b$ is finitely presented over $B$
by the equation $bT-1=0$, and $B\to B_b$ is formally \'{e}tale by
\ref{res:FL-10}. For the principal-cover assertion, the $B_{b_i}$
are smooth. By \ref{res:SM-07}, $B$ is smooth. Localization
identifies $(\Omega_{B/A})_{b_i}$ with $\Omega_{B_{b_i}/A}$
by \ref{res:KD-16}. A module that is zero on this cover is zero:
the annihilator of any element contains a power of every $b_i$,
and these powers still generate the unit ideal. This proves the
assertion using \ref{res:ET-02}.

At a smooth prime, shrink first to a smooth principal neighborhood.
The module $\Omega_{B/A}$ is finitely generated by \ref{res:KD-13}.
If its localization at the prime is zero, multiply the finitely many
denominators killing a generating set to get a principal neighborhood
on which it is zero. The converse follows by localization.
\end{proof}

\begin{proposition}[Cancellation in an \'{e}tale tower]\label{res:ET-05}
Suppose $A\to B\to C$ are compatible. If $B/A$ is formally
unramified and $C/A$ formally \'{e}tale, then $C/B$ is formally
\'etale. In particular, if $B/A$ and $C/A$ are both \'{e}tale, then
$C/B$ is \'{e}tale.
\end{proposition}
\begin{proof}
Give a $B$-algebra $D$, a square-zero ideal $J$, and a $B$-algebra map
$C\to D/J$. Its unique $A$-algebra lift $C\to D$ restricts on $B$ to
the specified structural map. Indeed, both restrictions lift the
same map $B\to D/J$, and formal unramifiedness gives equality.
This proves the formal claim. For the finite presentation claim, take
finite generators $b_i$ of $B/A$, and write
$C=A[Y_1,\ldots,Y_m]/(g_1,\ldots,g_\ell)$.
Represent the images of $b_i$ in $C$ by polynomials $h_i(Y)$.
Over $B$ use the relations $g_j(Y)$ and $b_i-h_i(Y)$. The map from
this quotient to $C$ has inverse induced by the $A$-presentation of
$C$. On the $Y_j$ the two composites are identities, and on the
coefficients from $B$ this follows from the relations for the $b_i$.
\end{proof}

\section{Standard local forms and quotients}

\begin{definition}[Standard \'etale algebra]\label{res:ET-06}
A standard \'{e}tale $A$-algebra has an $A$-algebra presentation
$(A[X]/(f))_g$ with $f$ monic, $g\in A[X]$, and $f'$ invertible in
the localized quotient. An $A$-algebra that is \'{e}tale on some
principal neighborhood of a prime $\mathfrak q$ is \'{e}tale at
$\mathfrak q$.
\end{definition}

\begin{proposition}[Lifting a simple root]\label{res:ET-07}
If $f'$ is a unit in $(A[X]/(f))_g$, then this algebra is
\'etale over $A$. In particular every standard \'{e}tale algebra
is \'{e}tale. Principal localization and arbitrary base change
preserve the standard form.
\end{proposition}
\begin{proof}
The algebra has finite presentation by adjoining a variable $Y$
with $g(X)Y-1=0$. Let $D$ be an $A$-algebra and $J^2=0$ in $D$.
A map to $D/J$ is given by a root $\bar x$ of $f$ at which
$g$ is invertible. For any lift $x\in D$, $g(x)$ and $f'(x)$ are units:
an inverse modulo $J$ can be adjusted because $1+j$ has inverse $1-j$
for $j\in J$. Also $f(x)\in J$. For $u\in J$, expansion of each
monomial gives $f(x+u)=f(x)+u f'(x)$. Thus exactly one choice,
$u=-f(x)/f'(x)$, yields a root. Its $g$-value is still a unit,
and the image of $Y$ is forced. This proves the lifting claim,
without using monicity in the calculation.

If $C=(A[X]/(f))_g$ and $c=h/g^m\in C$, then $C_c$ is
$(A[X]/(f))_{gh}$. The original $f$ remains monic and $f'$
remains invertible.
Under any base change, use the images of $f,g$ in $A'[X]$.
Monicity is preserved, and invertibility of the image of the
derivative follows from its inverse in $C$.
\end{proof}

\begin{proposition}[\'Etale quotients and idempotents]\label{res:ET-08}
For an ideal $I$ in $A$, $A\to A/I$ is formally \'{e}tale if and
only if $I^2=I$. It is \'{e}tale if and only if $I=Ae$ for an
idempotent $e\in A$. In this case $A/I\simeq A_{1-e}$ over $A$.
\end{proposition}
\begin{proof}
If the map is formally \'{e}tale, lift the identity on $A/I$
through $A/I^2\to A/I$. Its $A$-linearity forces $I$ to map to
zero in $A/I^2$, so $I=I^2$. Conversely, for any lifting problem
through $D\to D/J$ with $J^2=0$, the structural map sends $I$
into $J$, hence sends $I=I^2$ to zero. There is exactly one
$A$-map $A/I\to D$ factoring that structural map.

If the quotient is finitely presented, take a finite polynomial
presentation, and choose $a_i\in A$ lifting the images of its
variables. All elements of $I$ lie in the ideal of $A$ generated by
the defining relations evaluated at the $a_i$: a constant in the
kernel is in the relation ideal in the polynomial ring, and evaluation
at the $a_i$ gives the claim. Conversely, finite generators of $I$
present the quotient with no variables. Now assume $I$ is finitely
generated and $I=I^2$. Apply the determinant argument in the proof
of \ref{res:ET-03} to $I$ to get $s\in1+I$ with $sI=0$.
Put $e=1-s\in I$. Then $ex=x$ for all $x\in I$, so $I=Ae$
and $e^2=e$. Conversely an idempotent generates a finitely generated
idempotent ideal. Inverting $1-e$ kills $e$, and in $A/Ae$ the
image of $1-e$ is $1$. The universal properties give the isomorphism.
\end{proof}

\begin{theorem}[Local one-variable presentation]\label{res:ET-09}
If $B$ is finitely presented over $A$ and \'{e}tale at
$\mathfrak q$, then some $B_b$ with $b\notin\mathfrak q$ is
standard \'{e}tale over $A$. For an \'{e}tale algebra the
resulting principal opens admit a finite subcover.
\end{theorem}
\begin{proof}
We use the proved local unramified presentation theorem
\texttt{Algebra.IsUnramifiedAt.}\allowbreak
\texttt{exists\_hasStandardEtaleSurjectionOn} in fixed Mathlib.
It says that a finite-type $A$-algebra $B$ whose localization
$B_{\mathfrak q}$ is formally unramified admits a surjection
$C\twoheadrightarrow B_f$ from a standard \'{e}tale $A$-algebra
for some $f\notin\mathfrak q$. Our finite presentation gives finite
type. By hypothesis there is $t\notin\mathfrak q$ with $B_t$ \'{e}tale;
further localization gives formal unramifiedness of $B_{\mathfrak q}$.

Choose a preimage of $t$ in $C$ and invert it. This gives a
surjection $C'\twoheadrightarrow B_{ft}$ with $C'$ still standard
\'etale by \ref{res:ET-07}. The target is \'{e}tale over $A$ by
\ref{res:ET-04}. By \ref{res:ET-05}, the surjection is \'{e}tale as
a map from $C'$, so by \ref{res:ET-08} the quotient is a principal
localization of $C'$. It is standard \'{e}tale again by
\ref{res:ET-07}. The finite subcover follows because a family of
principal opens that covers the spectrum generates the unit ideal,
and a representation of $1$ uses finitely many generators.
\end{proof}

\section{Fields, flatness, and fibers}

\begin{proposition}[Extensions of fields]\label{res:ET-10}
Every separable algebraic field extension $L/K$ is formally
\'etale. A finite extension is \'{e}tale exactly when it is
separable. More generally, for an extension $L/K$ finitely
generated as a field, formal unramifiedness is equivalent to
being finite separable.
\end{proposition}
\begin{proof}
Each finite separable subextension is simple by the primitive element
theorem \cite[Chapter 12]{d07}. Its minimal polynomial has invertible
derivative at the generator, so \ref{res:ET-07} applies. For a general
separable algebraic extension, use the directed union of its finite
subextensions. Given any square-zero lifting problem, each restriction
has a unique lift, and those lifts agree under restriction by uniqueness.
They therefore define a unique lift on $L$.

For the converse in the finitely generated field case, use the
\href{https://github.com/leanprover-community/mathlib4/blob/56a005855ea17e0b0958f73ec801ae140b0a1679/Mathlib/RingTheory/Unramified/Field.lean\#L217}
{fixed-Mathlib separability criterion for formally unramified field extensions}.
A finitely generated
extension field is a localization of the algebra generated by the
chosen field generators, so it satisfies the essentially finite type
hypothesis. The separability in that criterion includes algebraicity,
and finitely many algebraic field generators imply finite degree.
This also proves the claimed converse for a finite \'{e}tale
extension, since \'{e}taleness implies formal unramifiedness.
\end{proof}

\begin{theorem}[The field case]\label{res:ET-11}
For a finite-type algebra $B$ over a field $K$, the following
conditions are equivalent:
\begin{enumerate}
\item $\Omega_{B/K}=0$;
\item $B$ is \'{e}tale over $K$;
\item $B$ is a finite product of finite separable extensions of $K$.
\end{enumerate}
The empty product gives the zero ring.
\end{theorem}
\begin{proof}
For the forward direction use two fixed-Mathlib results for a
formally unramified, essentially finite-type algebra over a field,
\texttt{Algebra.FormallyUnramified.finite\_of\_free} and
\texttt{Algebra.FormallyUnramified.isReduced\_of\_field}.
The first applies because any vector space over a field is free.
Thus (1), by \ref{res:FL-04}, makes $B$ finite-dimensional
and reduced. A finite-dimensional $K$-algebra is Artinian by
the descending chain condition on subspaces. A commutative
Artinian ring is a finite product of local Artinian rings
\cite[Chapter 3, ALG-03-C03-S02]{p103}, with nilpotent maximal
ideal in each factor. A reduced factor therefore has zero maximal
ideal and is a field. This gives the product of residue fields.
Each factor is a quotient of $B$; its differentials over $K$ vanish
by the surjectivity in \ref{res:KD-10}, or by \ref{res:KD-20}.
By \ref{res:ET-10}, each finite field factor is separable.

For (3) implies (2), a finite product of finite extensions is a
finite-dimensional algebra. It is of finite type, and it is
finitely presented by Hilbert's basis theorem
\cite[Chapter 5]{p103}. By \ref{res:ET-10} each factor is formally
\'etale. In a square-zero lifting problem the orthogonal idempotents
of the product lift uniquely. To see this, for an idempotent
$\bar e$ and a lift $e$, the defect $e^2-e$ lies in the square-zero
ideal and $2e-1$ is a unit modulo it. Correcting $e$ by
$-(e^2-e)/(2e-1)$ gives the unique idempotent lift, by the
calculation in \ref{res:ET-07}. Products of distinct lifts are
idempotent lifts of zero, and their sum is an idempotent lift of
one. They therefore split the target into the corresponding
factors. Lift the map from each field factor separately. This
gives one lift of the product map, uniquely, and proves (2).
The implication (2) implies (1) is \ref{res:ET-02}.
\end{proof}

\begin{theorem}[The flat unramified criterion]\label{res:ET-12}
An $A$-algebra $B$ is \'{e}tale if and only if it is finitely
presented, flat over $A$, and $\Omega_{B/A}=0$.
\end{theorem}
\begin{proof}
Necessity follows from \ref{res:ET-02} and \ref{res:SM-08}.
For sufficiency, write $B=A[X_1,\ldots,X_n]/I$, with $I$ finitely
generated. Fix a prime $\mathfrak q$ in $B$ and let
$\mathfrak p=\mathfrak q\cap A$. Because the cokernel of the
conormal map vanishes, at $\mathfrak q$ its image spans
$\kappa(\mathfrak q)^n$. Choose $n$ of the original relations
$f_1,\ldots,f_n$ whose differential vectors are a basis in
that fiber. Let $\Delta$ be their Jacobian determinant, and form
\[
 C=\big(A[X_1,\ldots,X_n]/(f_1,\ldots,f_n)\big)_\Delta.
\]
There is a surjection $C\to B_\delta$, where $\delta$ is the
image of $\Delta$ in $B$. Present $C$ by adding a variable $T$
and the relation $\Delta T-1$. The square Jacobian is block
triangular with determinant $\Delta^2$, a unit in $C$, so
$C$ is \'{e}tale by \ref{res:ET-03}. The kernel $J$ of the
surjection is generated by the remaining relations.
For $n=0$ the empty selection and $\Delta=1$ give the same
construction.

Write $\mathfrak r$ for the inverse image in $C$ of the prime of
$B_\delta$ corresponding to $\mathfrak q$. Localize the exact
sequence at $\mathfrak r$ to obtain
\[
 0\longrightarrow J_{\mathfrak r}\longrightarrow C_{\mathfrak r}
   \longrightarrow (B_\delta)_{\mathfrak q B_\delta}
   \longrightarrow0 .
\]
The last module is flat over $A$. In fact localization of a flat
$A$-algebra is flat over $A$, as follows by composing the two exact
tensor functors. Tensoring this sequence over $A$ with
$\kappa(\mathfrak p)$ is left exact here: the $\operatorname{Tor}_1$
term of the flat quotient is zero \cite[Chapter 11]{p103}.
By \ref{res:ET-04} and \ref{res:ET-11}, the fiber
$C\otimes_A\kappa(\mathfrak p)$ is a finite product of fields.
Its localization at the prime induced by $\mathfrak r$ is a
field and is precisely $C_{\mathfrak r}\otimes_A\kappa(\mathfrak p)$.
The localized quotient fiber is nonzero since it has a map to
$\kappa(\mathfrak q)$. A field surjecting onto a nonzero ring
has zero kernel. Thus $J_{\mathfrak r}\otimes_A\kappa(\mathfrak p)=0$.
The elements of $A\setminus\mathfrak p$ act invertibly on
$J_{\mathfrak r}$ since their images lie outside $\mathfrak r$.
The tensor product is therefore $J_{\mathfrak r}/\mathfrak pJ_{\mathfrak r}$.
Since $\mathfrak pC_{\mathfrak r}$ is contained in the maximal
ideal, Nakayama's lemma gives $J_{\mathfrak r}=0$
\cite[Chapter 7]{p103}.

Clear a denominator outside $\mathfrak r$ for each of the finitely
many generators of $J$ and multiply them to obtain
$s\in C\setminus\mathfrak r$ with $J_s=0$. Then $C_s$ is
isomorphic to the corresponding localization of $B_\delta$.
This shows that $\mathfrak q$ has an \'{e}tale principal
neighborhood. These neighborhoods cover $\operatorname{Spec}(B)$.
Choose finitely many whose defining elements generate the unit ideal,
and apply \ref{res:ET-04}.
\end{proof}

\begin{proposition}[Local residue fields and the diagonal]\label{res:ET-13}
Suppose $B$ is \'{e}tale at $\mathfrak q$, and let
$\mathfrak p=\mathfrak q\cap A$. Then
$\mathfrak qB_{\mathfrak q}=\mathfrak pB_{\mathfrak q}$ and
$\kappa(\mathfrak q)/\kappa(\mathfrak p)$ is finite separable.
If $B/A$ is \'{e}tale, its multiplication map
$B\otimes_A B\to B$ is the localization at an idempotent element
up to a canonical isomorphism of algebras over $B\otimes_A B$.
\end{proposition}
\begin{proof}
Take an \'{e}tale principal neighborhood and base change to
$\kappa(\mathfrak p)$ by \ref{res:ET-04}. Its fiber is a
finite product of finite separable extensions by \ref{res:ET-11}.
The local ring of the fiber at the point in question is therefore
the residue field there. In terms of the original rings, that local
fiber is $B_{\mathfrak q}/\mathfrak p B_{\mathfrak q}$, which
proves both assertions about the local rings and residue fields.

Put $D=B\otimes_A B$ and give it its $B$-algebra structure via
the first factor. Then $B\to D$ is \'{e}tale by base change, and
the composite $B\to D\to B$ is the identity. Hence $D\to B$ is
\'etale by \ref{res:ET-05}. It is surjective, so \ref{res:ET-08}
identifies its kernel with $De$ for an idempotent $e$, and
$B$ with $D_{1-e}$. Thus in affine scheme language the diagonal
is both open and closed.
\end{proof}

\section{Change of base and affine use}

\begin{proposition}[Descent and flat pullback of the locus]\label{res:ET-14}
Let $A\to A'$ and $B'=A'\otimes_A B$.
If $A'$ is faithfully flat over $A$, then $B'$ being \'{e}tale
over $A'$ implies $B$ is \'{e}tale over $A$.
If $B$ is finitely presented and $A'$ is flat over $A$, the
\'etale locus of $B'/A'$ is the inverse image of that of $B/A$.
\end{proposition}
\begin{proof}
Under faithful flatness, finite presentation descends by
\href{https://stacks.math.columbia.edu/tag/08XE}{Stacks Tag 08XE},
also proved in the
\href{https://github.com/leanprover-community/mathlib4/blob/56a005855ea17e0b0958f73ec801ae140b0a1679/Mathlib/RingTheory/Finiteness/Descent.lean\#L98}
{fixed-Mathlib finite-presentation descent theorem}.
Formal \'{e}taleness then descends by \ref{res:FL-17}.
For the pullback statement, a prime $\mathfrak q'$ of $B'$
over $\mathfrak q$ of $B$ gives a flat local map
$B_{\mathfrak q}\to B'_{\mathfrak q'}$. Such a map is faithfully
flat, since its closed fiber is nonzero and faithful flatness
is detected on maximal fibers \cite[Chapter 2]{p103}.
Under flat base change the smooth locus pulls back by
\ref{res:SM-13}. The vanishing condition for the localized
differentials also pulls back by \ref{res:KD-17} and faithful
flatness of this local map. Use the locus description in
\ref{res:ET-04}.
\end{proof}

\begin{proposition}[Finite coefficient models]\label{res:ET-15}
Every \'{e}tale $A$-algebra is obtained by base change from an
\'etale algebra over a finitely generated $\mathbb Z$-algebra.
More generally, if $A$ is the filtered colimit of rings $A_i$,
an \'{e}tale $A$-algebra descends to an \'{e}tale algebra over
some $A_i$. If $A$ is a localization $S^{-1}A_0$, the descended
algebra may be taken over $(A_0)_s$ for some $s\in S$.
\end{proposition}
\begin{proof}
Use the square presentation of \ref{res:ET-03}. In its polynomial
ring, choose polynomials $u,h_1,\ldots,h_n$ with
$u\Delta-1=\sum_j h_jf_j$, where $\Delta$ is the determinant in
that presentation. This is a finite certificate for invertibility
of $\Delta$ in the quotient. The finitely many coefficients of
the polynomials generate a finitely generated subring of $A$.
The same identity already holds in its polynomial ring. This
gives the first assertion by \ref{res:ET-03}.
For a filtered colimit, lift the coefficients to a common stage
and pass to a later stage at which the finitely many coefficient
equalities in the certificate hold. The presentation at that
stage is square with unit determinant. A localization is the
filtered colimit of its principal localizations indexed by
finite products from its multiplicative set, giving the last claim.
\end{proof}

\begin{example}[Computations and finiteness distinctions]\label{res:ET-16}
For any $a\in A$, $A\to A_a$ is \'{e}tale. Its presentation
$A[X]/(aX-1)$ has Jacobian $a$, a unit in the quotient.
An arbitrary localization is formally \'{e}tale by \ref{res:FL-10}
and may fail to be of finite type. For
example, a finite list of rational generators has denominators
involving only finitely many primes, so cannot generate $\mathbb Q$
over $\mathbb Z$. If $2,a\in A^\times$, then $A[X]/(X^2-a)$
is \'{e}tale because $2X$ is a unit in the quotient.

For a prime integer $p$, the quotient
$\mathbb Z\to\mathbb F_p$ is finitely presented unramified
by \ref{res:KD-06} and \ref{res:FL-04}. Tensoring
$\mathbb Z\xrightarrow{p}\mathbb Z$ with $\mathbb F_p$ produces
the zero map on a nonzero module, so it is not flat or \'{e}tale.
Over a field $k$, $k[\epsilon]/(\epsilon^2)$ has nonzero
differentials over $k$ by \ref{res:KD-21}, including in
characteristic two. The same quotient regarded over $k[X]$ via
$X\mapsto\epsilon$ has zero relative differentials, and its
kernel $(X^2)$ is not idempotent. The two base maps illustrate
the dependence of unramifiedness and \'{e}taleness on the
specified base algebra.
\end{example}

\begin{corollary}[Finite products]\label{res:ET-17}
A finite product of $A$-algebras is \'{e}tale exactly when each
factor is \'{e}tale.
\end{corollary}
\begin{proof}
Apply \ref{res:ET-02}, the smooth product criterion
\ref{res:SM-14}, and the differential product isomorphism
\ref{res:KD-20}. The empty product is the zero algebra, which is
smooth by \ref{res:SM-14} and has zero differential module
by \ref{res:KD-06}.
\end{proof}

\begin{proposition}[Pointwise flat and separable criterion]\label{res:ET-18}
Let $B$ be finitely presented over $A$, let $\mathfrak q$ be a
prime of $B$, and put $\mathfrak p=\mathfrak q\cap A$. Then $B/A$ is \'{e}tale at
$\mathfrak q$ if and only if
\begin{enumerate}
\item $B_{\mathfrak q}$ is flat over $A_{\mathfrak p}$;
\item $\mathfrak p B_{\mathfrak q}=\mathfrak q B_{\mathfrak q}$; and
\item $\kappa(\mathfrak q)$ is finite separable over $\kappa(\mathfrak p)$.
\end{enumerate}
In particular the fibers of an \'{e}tale algebra are finite products
of finite separable field extensions of the base residue fields.
\end{proposition}
\begin{proof}
For necessity, use \ref{res:ET-02}, \ref{res:SM-08} and
\ref{res:ET-13}. Flatness after localization over $A$ is also
flatness over $A_{\mathfrak p}$, because tensor on
$A_{\mathfrak p}$-modules agrees with the tensor over $A$.
For sufficiency, conditions (2) and (3) identify the local
fiber with a finite separable field over $\kappa(\mathfrak p)$.
By \ref{res:KD-16}, \ref{res:KD-17} and \ref{res:ET-10}, the
differentials of $B$ tensored with $\kappa(\mathfrak q)$ vanish.
The construction in the proof of \ref{res:ET-12} at a single
prime uses only this vanishing to choose the determinant, and
flatness of $B_{\mathfrak q}$ over $A$ to preserve the injection
of the localized kernel when tensoring with $\kappa(\mathfrak p)$.
The latter follows from (1) and flatness of $A_{\mathfrak p}$
over $A$. Thus it produces an \'{e}tale principal neighborhood
of $\mathfrak q$. The final assertion follows from arbitrary
base change and \ref{res:ET-11}; it also gives the finite-fiber
property of the associated affine morphism.
\end{proof}

\begin{proposition}[Lifting from a quotient and from a residue field]
\label{res:ET-19}
For any ideal $I\subseteq A$, every \'{e}tale $A/I$-algebra
$\overline B$ is $B/IB$ for some \'{e}tale $A$-algebra $B$.
Also, if $L/\kappa(\mathfrak p)$ is a finite separable field
extension at a prime $\mathfrak p$ of $A$, some \'{e}tale
$A$-algebra has a prime over $\mathfrak p$ with residue field $L$.
\end{proposition}
\begin{proof}
Take a square presentation of $\overline B$ by \ref{res:ET-03}.
Lift the equations to $f_1,\ldots,f_n$ over $A$ and let $\Delta$
be their Jacobian determinant. Define $B$ by these equations and
one extra equation $\Delta T-1$. Its Jacobian determinant is
$\Delta^2$, which is a unit, so $B$ is \'{e}tale. Reducing modulo
$I$ gives $\overline B$ because the inverse of the old determinant
was already present there.

For the residue-field assertion, choose a primitive element of
$L/\kappa(\mathfrak p)$ and its monic minimal polynomial $\bar f$
\cite[Chapter 12]{d07}. The map
$A_{\mathfrak p}\to\kappa(\mathfrak p)$ is surjective. Lift its
finitely many coefficients to $A_{\mathfrak p}$ and hence to
$A_s$ for a common $s\notin\mathfrak p$, giving a monic polynomial
$f$. Let $B=(A_s[X]/(f))_{f'}$. It is \'{e}tale over $A_s$
by \ref{res:ET-07}, and over $A$ by \ref{res:ET-04} and
\ref{res:ET-16}. Its fiber over $\mathfrak p$ is
$\kappa(\mathfrak p)[X]/(\bar f)=L$, since the derivative
has nonzero image in $L$. The prime of that fiber gives the
required prime and residue field.
\end{proof}

\begin{example}[A composite of standard algebras]\label{res:ET-20}
The maps $\mathbb F_2\to\mathbb F_4$ and
$\mathbb F_4\to\mathbb F_4^4$ are standard \'{e}tale.
The first uses $X^2+X+1$, and the second uses
$\prod_{a\in\mathbb F_4}(X-a)$, whose distinct linear factors
give a product by the Chinese remainder theorem and whose
derivative is a unit at each factor. Their composite is
\'etale by \ref{res:ET-04}. It has no single standard
presentation over $\mathbb F_2$. In any localized one-variable
quotient over $\mathbb F_2$, primes with residue field
$\mathbb F_4$ correspond to surviving irreducible polynomials
of degree two in $\mathbb F_2[X]$. There is only one such
polynomial, since the constant term must be $1$ and absence
of a root at $1$ forces the linear coefficient to be $1$.
But the product $\mathbb F_4^4$ has four such primes.
\end{example}

\section*{Passage to affine schemes}

The local comparison in \ref{res:KD-16} identifies the restriction
of the differential module on a principal open with the module
for its localized algebra. Its compatibility with iterated
localization follows from the formula $db\mapsto d(b/1)$ and
generation by such differentials after extending scalars.
The base-change map of \ref{res:KD-17} is compatible with these
restriction maps for the same reason. These are the inputs for
constructing the relative differential sheaf on affine charts.
The local criteria in \ref{res:SM-05} and \ref{res:ET-09} give
the corresponding smooth and \'{e}tale charts. The pointwise
comparison in \ref{res:KD-24} and the residue-field result in
\ref{res:ET-13} identify the fibers that enter the geometric
formulation.

\part{Local Algebra, Depth, and Normality}
\label{part:local}
\chapter{Regular Sequences and the Koszul Complex}\label{ch:regular-koszul}

A regular sequence gives two kinds of control over a quotient. Its Koszul
complex supplies a finite free resolution when the ambient module is the
ring itself. Its associated graded module describes the successive powers
of the generated ideal. We develop both descriptions before using Koszul
homology to measure depth.

\section{Successive nonzerodivisors}

\begin{definition}[Weak and proper regular sequences]\label{res:RK-01}
Let $R$ be a ring, let $M$ be an $R$-module, and put $I_j=(x_1,\ldots,x_j)$,
with $I_0=0$. The ordered list $x=(x_1,\ldots,x_r)$ is
\emph{weakly $M$-regular} if multiplication by $x_j$ on
$M/I_{j-1}M$ is injective for every $j$. It is \emph{$M$-regular} if,
in addition, $M/I_rM\ne0$. A regular sequence in a ring means an
$R$-regular sequence.
\end{definition}

The empty list is weakly regular on every module and regular exactly when
$M\ne0$. Every list is weakly regular on the zero module, and none is
regular on it. A weakly regular sequence can start with a unit, but such
a sequence is not regular because the first quotient is already zero.
The properness condition also implies that each of the intermediate
quotients in a regular sequence is nonzero.

\begin{proposition}[Concatenation, scalar extension, and localization]
\label{res:RK-02}
An ordered concatenation $(x,y)$ is weakly $M$-regular if and only if $x$
is weakly $M$-regular and $y$ is weakly regular on $M/(x)M$. The corresponding
equivalence holds for regularity. If $N$ is an $R$-flat $S$-module for a ring
map $R\to S$, a weakly $M$-regular sequence remains weakly regular on
$M\otimes_R N$, both as an $R$-sequence and through its image in $S$. It is
regular on that module if the final quotient is nonzero. If $N$ is faithfully
flat over $R$, both kinds of regularity are reflected. In particular, for a
finite module $M$ and a prime $\mathfrak p$ containing $(x)$ with
$M_{\mathfrak p}\ne0$, an $M$-regular sequence remains regular on
$M_{\mathfrak p}$.
\end{proposition}
\begin{proof}
The successive quotients for the concatenation agree with the quotients
for $x$ followed by those for $y$ on $M/(x)M$.
For scalar extension, right exactness identifies the intermediate
quotient with $(M/I_jM)\otimes_R N$, and multiplication by $x_{j+1}$
is the tensor of the original multiplication map. Flatness preserves
its injectivity. Under faithful flatness, the kernel before tensoring
vanishes if it vanishes after tensoring, and so does the final quotient.
For the prime assertion, $R_{\mathfrak p}$ is flat and the final
quotient is nonzero by Nakayama's lemma, since the images of all $x_i$
lie in the maximal ideal of $R_{\mathfrak p}$.
\end{proof}

The order cannot in general be changed. In $R=k[X,Y,Z]$ over a field, the
list $X,Y(1-X),Z(1-X)$ is regular by reduction modulo $X$. For the order
$Y(1-X),Z(1-X),X$, the class of $Y$ modulo $Y(1-X)$ is nonzero and is
killed by $Z(1-X)$. Nonvanishing follows by evaluation at $X=1$.

\section{Construction and comparison maps}

\begin{definition}[The Koszul complex]\label{res:RK-03}
For a linear map $u:E\to R$ from an $R$-module, the Koszul complex
$K(u)$ has degree $q$ term $\bigwedge^q_R E$, in degrees $q\ge0$.
The differential is the contraction
\[
 d(e_1\wedge\cdots\wedge e_q)
 =\sum_{i=1}^q(-1)^{i-1}u(e_i)\,
     e_1\wedge\cdots\wedge\widehat{e_i}\wedge\cdots\wedge e_q.
\]
For a list $x=(x_1,\ldots,x_r)$, write $K(x)$ for the case
$E=R^r$ and $u(e_i)=x_i$, where $e_i$ is the $i$th standard basis
vector. Write $K(x;M)=K(x)\otimes_R M$ and $H_q(x;M)$ for its
homology. For the empty list the complex consists of $M$ in degree zero.
\end{definition}

The formula is multilinear and alternating. If two arguments agree,
the two summands that remove one of the equal arguments cancel, while
the others contain both. The alternating exterior identities are thus
valid even in characteristic two. In $d^2$, the two ways to delete
two positions have opposite signs and the same scalar coefficient,
so $d^2=0$. Splitting the deletion sum at the boundary between a
$p$-fold wedge and the next wedge gives
$d(a\wedge b)=da\wedge b+(-1)^p a\wedge db$. In the finite free case
$K_q(x)$ has the basis of increasing $q$-fold wedges, with
$\binom rq$ basis elements, and $H_0(x;M)=M/(x)M$.

\begin{proposition}[Functoriality and change of coordinates]
\label{res:RK-04}
If $u'\colon E'\to R$ and $\phi\colon E\to E'$ satisfy $u'\phi=u$,
the exterior powers of $\phi$ give a chain map $K(u)\to K(u')$.
These maps respect identities and composition. For every ring map $R\to S$,
there is a natural comparison $S\otimes_R K(u)\simeq K(S\otimes_R u)$.
In particular permutations and invertible linear changes of a
generating list give isomorphic Koszul complexes. Flat base change
also commutes with Koszul homology with coefficients.
\end{proposition}
\begin{proof}
Apply the relation $u'\phi=u$ term by term in the contraction sum to
verify the chain-map identity. Exterior powers are defined by the
universal alternating quotient of tensor powers. After scalar extension
that quotient satisfies the alternating universal property over $S$,
so the comparison is an isomorphism and respects the displayed
differentials. An invertible change of coordinates is an isomorphism
of the two linear maps after choosing the corresponding transformation
of the bases. For a flat scalar extension, cycles and boundaries
commute with tensoring by exactness, hence so do their quotients.
\end{proof}

\begin{proposition}[Direct sums, the cone, and its exact sequence]
\label{res:RK-05}
The complex for $u\oplus v:E\oplus F\to R$ is the tensor product of
the complexes for $u$ and $v$, with the total differential
$d(a\otimes b)=da\otimes b+(-1)^{\deg a}a\otimes db$. If $y$ is
appended to a list $x$, $K(x,y;M)$ is the mapping cone of multiplication
by $y$ on $C=K(x;M)$. In particular there is a natural long exact
sequence with segment
\[
 H_q(C)\xrightarrow{y}H_q(C)\longrightarrow H_q(x,y;M)
 \longrightarrow H_{q-1}(C)\xrightarrow{y}H_{q-1}(C).
\]
\end{proposition}
\begin{proof}
Exterior multiplication identifies the tensor product of the two
exterior algebras, with the graded product, with the exterior algebra
of the direct sum. To see invertibility, send $e\in E$ and $f\in F$
to their degree-one images in the tensor product. Their sum for
$(e,f)$ has square zero, by the exterior relations and the sign in the
graded product. The universal extension to the exterior algebra is
inverse on both sets of generators. The contraction derivation gives
the asserted tensor differential.

For one appended vector $e$ use $a+e\wedge b$ to identify degree $q$
with $C_q\oplus C_{q-1}$. Its boundary is
$(da+yb,-db)$. The first summand is the subcomplex $C$, and the
quotient is the degree shift with sign $-d$. The connecting map
sends the class of a cycle $b$ to the class of $yb$, by taking
the boundary of the lift $(0,b)$. Thus the homology sequence of this
short exact sequence gives the stated signs.
\end{proof}

\begin{proposition}[Homotopies and redundant generators]\label{res:RK-06}
Multiplication by every element of $I=(x)$ is homotopic to zero on
$K(x;M)$. The ideal $I$ annihilates its homology and its Koszul
cohomology with coefficients. If $I=R$ the complex is contractible.
Appending $s$ redundant elements of $I$ gives an isomorphism of
complexes with $K(x;M)\otimes_R\bigwedge^\bullet R^s$, where the
second differential is zero. In particular, for $M/IM\ne0$ the
quantity $r-\max\{q:H_q(x;M)\ne0\}$ depends only on $I$ and $M$.
\end{proposition}
\begin{proof}
If $a=\sum_i b_ix_i$, wedge on the left by $e=\sum_i b_i e_i$ provides
$h$ with $dh+hd=a\,\mathrm{id}$. The same identity precomposed on the
dual cochain groups gives a cochain homotopy. Taking $a=1$ gives a
contraction. For one extra coordinate $a$, replace its basis vector
by $e_{r+1}-\sum_i b_ie_i$. This is an invertible change of basis
and the new coordinate has value zero. Do this for each extra
coordinate and apply the tensor description. The zero-differential
exterior complex is finite free in each degree, so adding $s$
coordinates increases the highest nonzero degree by exactly $s$.
Given two generating lists, append the second to the first and appeal
to the coordinate permutation isomorphism to compare in either order.
\end{proof}

\section{Acyclicity and its converse}

\begin{theorem}[The Koszul resolution]\label{res:RK-07}
If $x$ is weakly $M$-regular, then $H_q(x;M)=0$ for $q>0$.
If $x$ is $R$-regular, $K(x)$ augmented to $R/(x)$ is a finite free
resolution of length $r$. It follows that for every $M$,
\[
 H_q(x;M)\simeq \operatorname{Tor}_q^R(R/(x),M).
\]
In particular a list regular on $R$ and weakly regular on $M$
has zero higher Tor on this pair.
\end{theorem}
\begin{proof}
Induct on the length. For the preceding list, the only remaining
homology is its degree-zero quotient. In the long exact sequence of
\ref{res:RK-05}, the sole possible positive homology is degree one,
the kernel of the next multiplication map on that quotient. It is
zero by weak regularity. The augmentation onto the final quotient
then gives the claimed finite free resolution for $M=R$, using the
basis in \ref{res:RK-03}. Tensor this resolution with $M$ to compute Tor.
\end{proof}

\begin{theorem}[The degree-one converse]\label{res:RK-08}
Suppose $R$ is Noetherian, $M$ is finite, and the entries of $x$
lie in the Jacobson radical of $R$. If $H_1(x;M)=0$, then $x$ is
weakly $M$-regular. If $M\ne0$, it is $M$-regular. Hence for a
nonzero finite module over a Noetherian local ring and a list in
the maximal ideal, regularity, acyclicity in all positive degrees,
and $H_1(x;M)=0$ are equivalent.
\end{theorem}
\begin{proof}
Let $x=(x',y)$ and $C=K(x';M)$. The long exact sequence shows that
multiplication by $y$ is surjective on $H_1(C)$ and injective on
$H_0(C)$. The group $H_1(C)$ is finite, as a subquotient of a
finite module over a Noetherian ring. Nakayama gives $H_1(C)=0$.
Induction supplies weak regularity of $x'$, and the injection on
$H_0(C)=M/(x')M$ is the last condition. If $M\ne0$, Nakayama
also gives $M/(x)M\ne0$. The other implications follow from
\ref{res:RK-07} and the positive-degree condition itself.
\end{proof}

\begin{proposition}[Permutations, positive powers, and local generators]
\label{res:RK-09}
For a nonzero finite module over a Noetherian local ring, every
permutation of a regular sequence in the maximal ideal is regular.
If $x$ is weakly $M$-regular over an arbitrary ring, replacing its
entries by positive powers preserves weak regularity, and preserves
regularity when the original list was regular. More generally,
if a list remains weakly regular after replacing one coordinate
$a$ by $b$, then the list with that coordinate $ab$ is weakly
regular. If both given lists are regular, the product list is regular.
In the Noetherian
local case, every regular list in the maximal ideal is a minimal
generating list of its ideal, and every other minimal generating
list of that ideal is regular on the same module.
\end{proposition}
\begin{proof}
Permutation preserves $H_1$ by \ref{res:RK-04}, and the local
converse applies. For the product statement let $L$ be the quotient
of $M$ by the fixed prefix. Multiplication by $a$ and by $b$ on $L$
are injective, as is multiplication by $ab$. The exact sequence
\[
 0\longrightarrow L/aL\xrightarrow{\ b\ }L/abL
   \longrightarrow L/bL\longrightarrow0
\]
is injective on the left because $b$ acts injectively on $L$.
An element of the middle term mapping to zero on the right is
represented by $b\ell$, so it is in the image. The remaining suffix
is weakly regular on the outer terms. It is therefore weakly regular
on the middle by the extension argument in \ref{res:RK-12}.
The new ideal is contained in both old ideals, so a nonzero final
quotient remains nonzero. Iterating products with $a$ gives the
positive powers, and the procedure applies to each coordinate.

In the local case, a redundant generator could be put last by
permutation. Its multiplication would be zero on the nonzero
preceding quotient, which is impossible. Two minimal generating
lists of the same ideal both reduce to bases of
$I/\mathfrak m I$, by Nakayama. Write one list as a matrix
times the other; reduction of that matrix is invertible, so its
determinant is a unit over the local ring. Apply \ref{res:RK-04}
and \ref{res:RK-08} to transport regularity.
\end{proof}

\section{Duality and the ideal filtration}

\begin{proposition}[Self-duality and coefficient cohomology]
\label{res:RK-10}
Let $E=R^r$, let $D=\bigwedge^r E$, and give
$\operatorname{Hom}_R(K(x),M)$ the cochain differential
$\delta(f)=f\circ d$. There is an identification of cochain
complexes
\[
 \operatorname{Hom}_R(K(x),M)^q
 \simeq K_{r-q}(x;M\otimes_R D^\vee),\qquad
 D^\vee=\operatorname{Hom}_R(D,R),
\]
where the differential on the reversed indexing is $d$.
After orienting $E$ by its ordered basis this identifies Koszul
cohomology $H^q(x;M)$ with $H_{r-q}(x;M)$. For a regular
sequence on $R$ the cochain complex computes
$\operatorname{Ext}_R^q(R/I,M)$. In particular for $M=R$ the
Ext groups vanish except in degree $r$, where the orientation
identifies the Ext group with $R/I$.
\end{proposition}
\begin{proof}
Exterior multiplication into $D$ pairs every increasing basis
wedge with its unique complementary wedge, with coefficient $1$
or $-1$. Thus it is perfect, giving the required module identifications
after applying $D^\vee$ and $M$. For $w$ of degree $q+1$ and
$v$ of degree $r-q$, the wedge $w\wedge v$ has degree $r+1$
and is zero. Consequently
$(dw)\wedge v=(-1)^q w\wedge dv$. The pairing taking $v$ to
the map $w\mapsto (-1)^{q(q-1)/2}w\wedge v$ in degree $q$,
with evaluation in $M$ after applying $D^\vee$, therefore
commutes with the differentials. The ordered basis trivializes
$D$. The Ext assertions follow by applying $\operatorname{Hom}_R(-,M)$
to the resolution in \ref{res:RK-07} and, for $M=R$, using its
homology.
\end{proof}

\begin{theorem}[Rees's associated graded description]\label{res:RK-11}
Put $I=(x_1,\ldots,x_r)$ and define
$\operatorname{gr}_I M=\bigoplus_{n\ge0} I^nM/I^{n+1}M$.
Evaluation of the variables at the classes of the $x_i$ gives
a graded surjection
\[
 (M/IM)[X_1,\ldots,X_r]\longrightarrow\operatorname{gr}_I M
\]
of $(R/I)[X_1,\ldots,X_r]$-modules. If $x$ is weakly $M$-regular
this map is an isomorphism. In the case $M=R$ it is an isomorphism
of graded algebras. In particular $I/I^2$ is free over $R/I$ with
basis the classes of the $x_i$ for an $R$-regular sequence.

Call the list \emph{$M$-quasi-regular} if the graded map is an
isomorphism and $M/IM\ne0$. If each of
$M,M/x_1M,\ldots,M/(x_1,\ldots,x_{r-1})M$ is $I$-adically
separated, quasi-regularity implies regularity. In particular
when $R$ is Noetherian local, $M\ne0$ is finite, and $I$ lies
in the maximal ideal, it is equivalent to the conditions in
\ref{res:RK-08}.
\end{theorem}
\begin{proof}
Interpret module-valued polynomials as finite sums of
monomials with coefficients in the given module. The
degree-$n$ map sends $\sum m_\alpha X^\alpha$ to the class of
$\sum x^\alpha m_\alpha$ in $I^nM/I^{n+1}M$. It is well
defined, surjective by the description of $I^n$, and respects
the indicated graded action. It respects products in the
ring case.

For injectivity, induct on $r$. The case $r=0$ is the trivial
filtration. Set $a=x_r$ and $J=(x_1,\ldots,x_{r-1})$. The induction
hypothesis identifies $\operatorname{gr}_J M$ with a polynomial
module over $M/JM$. Multiplication by $a$ is injective on $M/JM$
by weak regularity and hence on every homogeneous component of
$\operatorname{gr}_J M$. Thus it is injective on $M/J^nM$ for every
$n\ge1$. To see this, if $am\in J^nM$, then injectivity modulo $JM$
first gives $m\in JM$. For $i<n$, injectivity on
$J^iM/J^{i+1}M$ takes $m\in J^iM$ to $m\in J^{i+1}M$.

Now induct on the degree $n$ of a homogeneous relation $F$ with
$F(x)\in I^{n+1}M$. When $n=0$ this just says its single coefficient
is in $IM$. In any degree, express $F(x)$ by a homogeneous degree
$n+1$ polynomial evaluated at $x$. Write that polynomial as
$\sum_i X_iG_i$ and subtract $\sum_i x_iG_i$ from $F$.
This changes each coefficient by an element of $IM$ and reduces
to $F(x)=0$. Write $F=G+X_rH$, with $G$ in the first $r-1$
variables. Then $aH(x)=-G(x)\in J^nM$. The preceding injection
gives $H(x)\in J^nM\subseteq I^nM$, so induction on degree
shows that all coefficients of $H$ lie in $IM$. Choose a
homogeneous polynomial $H'$ of degree $n$ in the first $r-1$
variables evaluating to $H(x)$, using $H(x)\in J^nM$.
Now $G+aH'$ evaluates to zero on the prefix. Induction on
length puts its coefficients in $JM$, so the coefficients of
$G$ lie in $IM$. This proves the graded isomorphism.
Its degree-one component gives the conormal basis assertion.

Conversely, suppose $x$ is quasi-regular. Multiplication
by $x_1$ on the graded module is multiplication by the indeterminate
$X_1$ under the graded identification. It is injective. If
$x_1m=0$ in $M$, then degree-by-degree injectivity implies
$m\in I^nM$ for every $n$, and separatedness gives $m=0$.
For $n\ge1$, the same argument stopped in degree $n-1$ gives
$x_1M\cap I^nM=x_1I^{n-1}M$: if $x_1m\in I^nM$, apply
injectivity in degrees $0$ through $n-2$ to conclude
$m\in I^{n-1}M$. The degree-$n$ kernel of the map to
$\operatorname{gr}_I(M/x_1M)$ is consequently the image of
$I^{n+1}M+x_1 I^{n-1}M$ in $I^nM/I^{n+1}M$, precisely the
degree-$n$ image of $X_1$. In degree zero the kernel is zero.
Thus the graded module of the quotient is the polynomial
module with $X_1$ removed. The remaining list is quasi-regular
there, with the same nonzero final quotient. Induct on length
using the stated separatedness of the successive modules.
In the Noetherian local case this separatedness follows from
Krull intersection \cite[Theorem ALG-05-C04-R04]{p103}.
\end{proof}

For a polynomial presentation $P\twoheadrightarrow B=P/I$ as in
\ref{res:KD-11} whose kernel is generated by a $P$-regular sequence,
\ref{res:RK-11} proves that the displayed relation classes are
a $B$-basis of the conormal term. It does not assert injectivity of the
conormal map. For example the regular element $X^2$ of $k[X]$ gives
the nonzero conormal class of $X^2$ for $k\to k[X]/(X^2)$, whose
image is $2X\,dX$. Even when $2$ is invertible, multiplying that
class by the nonzero element $\bar X$ kills its image without
killing the class. Thus regular relations alone do not give the
split criterion of \ref{res:FL-07}.

\begin{proposition}[Exactness with coefficients and reduction]
\label{res:RK-12}
A short exact sequence of coefficient modules induces a short exact
sequence of Koszul complexes and hence long exact Koszul homology
and cohomology sequences. If a list is weakly regular on every
term of an exact complex, quotienting every term by the generated
ideal gives an exact complex. For a short exact sequence
$0\to L\to M\to N\to0$, weak regularity on the two outer
modules implies weak regularity on the middle.
\end{proposition}
\begin{proof}
Each $K_q(x)$ is finite free, so tensoring and applying
$\operatorname{Hom}_R(K_q(x),-)$ are exact on coefficient sequences.
This yields the first assertion by the homology sequence.
If one element acts injectively in every degree of an exact complex,
multiplication by it gives a short exact sequence of complexes with
quotient the reduced complex. The first two complexes have zero
homology, so the third does as well. Repeat at each stage of a
list, using injectivity on the reduced modules.
For the last assertion, in the short exact sequence an element of
the middle module killed by the first entry has image zero on the
right and then preimage zero on the left. Quotienting by the first
entry keeps the sequence exact on the left since the right module
has no kernel for that multiplication. The same argument
iterates along the list.
\end{proof}

\begin{example}[Calculations and boundary cases]\label{res:RK-13}
For a nonzero ring $A$, the variables $X_1,\ldots,X_r$ form
a regular sequence in $R=A[X_1,\ldots,X_r]$. The successive
quotients are polynomial rings in the remaining variables,
in which the next variable is injective on coefficients. The Koszul
complex resolves $A$ as an $R$-module. If one also uses $A$ as
the coefficient module, all variables act as zero and
$H_q(x;A)=\bigwedge_A^q A^r$.

For $R=k[\epsilon]/(\epsilon^2)$ and the single element
$x=(\epsilon)$, the complex has one boundary map, multiplication
by $\epsilon$. Its degree-one homology is $(\epsilon)$ and its
degree-zero homology is $k$. Thus a nonzero nilpotent can give
positive Koszul homology. For a sequence containing a unit the
whole complex is contractible by \ref{res:RK-06}. A sequence with
$I$ proper whose Koszul complex is acyclic in positive degrees
need not be regular in its given order outside the local
hypotheses. Use the two orderings of the polynomial-ring example
following \ref{res:RK-02} and apply \ref{res:RK-04} and
\ref{res:RK-07}.
\end{example}

\begin{proposition}[Top homology and a regular element in the ideal]
\label{res:RK-14}
The ordered basis identifies $H_r(x;M)$ with $0:_M I$ and
$H^0(x;M)$ with $\operatorname{Hom}_R(R/I,M)$.
If $y\in I$ acts injectively on $M$, there are short exact
sequences
\[
 0\longrightarrow H_q(x;M)\longrightarrow H_q(x;M/yM)
   \longrightarrow H_{q-1}(x;M)\longrightarrow0.
\]
Here homology outside degrees $0,\ldots,r$ is zero. More generally,
if $M/IM\ne0$ and $y_1,\ldots,y_s\in I$ is an $M$-regular sequence,
then $s\le r$,
$H_q(x;M)=0$ for $q>r-s$, and
\[
 H_{r-s}(x;M)\simeq
 0:_{\,M/(y_1,\ldots,y_s)M}I.
\]
\end{proposition}
\begin{proof}
The differential in degree $r$ has components multiplication
by each $x_i$, with signs, and there is no incoming boundary.
The kernel is $0:_M I$. For $r=0$ this is $M$. Dually the
degree-zero cochain cycles are the maps from $R$ annihilating $I$,
which is the stated Hom group. For an injective multiplication
by $y$, apply \ref{res:RK-12} to
$0\to M\xrightarrow{y}M\to M/yM\to0$. Both homology maps
labelled $y$ vanish by \ref{res:RK-06}, giving the short
exact sequences.

For the iterated assertion $H_0(x;M)$
is nonzero, and quotienting by each $y_i$ leaves $H_0$
unchanged. If the highest nonzero degree at one stage is $h$,
the short exact sequences imply that it becomes $h+1$ and
that the old group in degree $h$ maps isomorphically to the
new group in degree $h+1$. Since no group exists above degree $r$,
$h+s\le r$. Applying the short exact sequences successively at
the highest potentially nonzero degrees identifies
$H_{r-s}(x;M)$ with the degree-$r$ group of the final quotient,
allowing both groups to be zero. The top-group formula completes
the identification.
\end{proof}

\begin{proposition}[Degree-one acyclicity and the graded filtration]
\label{res:RK-15}
Over an arbitrary ring, if $H_1(x;M)=0$ then the graded map in
\ref{res:RK-11} is an isomorphism. No separation or finiteness
hypothesis is needed for this direction. Both degree-one acyclicity
and acyclicity in all positive degrees are preserved under replacing
a coordinate by a positive power. More generally they are preserved
under multiplying two possible values of a fixed coordinate, provided
both resulting lists before multiplication satisfy the same acyclicity.
\end{proposition}
\begin{proof}
After a common permutation we may put the coordinate last. Let
$C$ be the coefficient Koszul complex of the fixed prefix. Use the
cone signs in \ref{res:RK-05}. In any degree $q$, suppose that the
cones of multiplication by $a$ and $b$ have zero $q$th homology.
A cycle $(u,v)$ in the cone of $ab$ gives the cycle $(u,av)$ in the
cone of $b$. Write $(u,av)=(dw+bz,-dz)$.
Subtract the boundary $(dw,0)$ from $(u,v)$ to get $(bz,v)$.
The pair $(z,v)$ is a cycle in the cone of $a$, since
$dz+av=0$ and $dv=0$. Its expression as a boundary
$(dW+aZ,-dZ)$ makes $(bz,v)$ the boundary of $(bW,Z)$
in the cone of $ab$. This proves the product assertion in each
degree. Iteration and permutation give the assertion about powers.

Now take a homogeneous degree-$n$ polynomial $F$ in the variables
$X_i$ with coefficients in $M$ and $F(x)\in I^{n+1}M$.
We may assume $r,n\ge1$ since the empty list and degree zero
are immediate. Set $a=x_r$ and consider the coefficient $m$ of
$X_r^n$. Modulo $\sum_{i<r} x_iM$, the evaluation $F(x)$ is
$a^nm$, while every element of $I^{n+1}M$ is represented by
$a^{n+1}m'$ modulo that same submodule. Their difference gives a
relation on the list $(x_1,\ldots,x_{r-1},a^n)$ with last
coefficient $m-am'$. This list is acyclic in degree one by the
power assertion. The relation is therefore a Koszul boundary.
The last coordinate of such a boundary belongs to
$\sum_{i<r}x_iM$, by the formula for $d(e_i\wedge e_r)$.
Consequently $m\in IM$.

To detect every coefficient, extend scalars to
$S=R[t_1,\ldots,t_r,t_r^{-1}]$. This ring is free as an $R$-module
with its Laurent monomials as basis, so the base-change comparison
in \ref{res:RK-04} preserves degree-one acyclicity.
Use the invertible change of coordinates
\[
 g_r=x_r/t_r,\qquad g_i=x_i-t_i g_r\quad(1\le i<r).
\]
In terms of the $g_i$, the coefficient of the pure power $g_r^n$
in the expression for $F(x)$ is $\sum_{|\alpha|=n}m_\alpha t^\alpha$.
The ideal generated by $g$ is $IS$. The pure-power argument just
proved places this coefficient in $I(M\otimes_R S)$. Independence
of the Laurent monomials after passing to $M/IM$ gives
$m_\alpha\in IM$ for each $\alpha$. This is injectivity of the
graded map in \ref{res:RK-11}.
\end{proof}
\chapter{Depth and Homological Comparisons}\label{ch:depth}

Regular sequences measure how long successive nonzerodivisors can be
chosen on a module. Their common maximal length also appears as a
vanishing threshold in Ext and in Koszul cohomology. The relative
version with respect to an ideal is needed to compare depths at
different primes. The finite projective dimension formula at the end
relates these thresholds to the length of a minimal free resolution.
The principal sources are \cite[Chapters 17--18]{eisenbud},
\cite[Sections 1.2--1.3 and 1.6]{brunsherzog}, and
\cite[Sections 16--17]{matsumura}.

\section{Grade and first nonvanishing Ext}

\begin{definition}[Grade, depth, and the zero-depth test]\label{res:DP-01}
Let $R$ be Noetherian, let $M$ be a finite $R$-module, and let $I$ be an
ideal of $R$. Set $\operatorname{grade}_R(I,M)=+\infty$ if $IM=M$.
Otherwise define it as the supremum of the lengths of $M$-regular
sequences in $I$, using the regular-sequence convention of
Chapter~\ref{ch:regular-koszul}. In a Noetherian local ring
$(R,\mathfrak m,k)$ put
$\operatorname{depth}_R M=\operatorname{grade}_R(\mathfrak m,M)$.
In particular the zero module has infinite depth. For $IM\ne M$ the
following are equivalent:
\[
 \operatorname{grade}_R(I,M)=0,\qquad
 I\subseteq\mathfrak p \text{ for some }\mathfrak p\in\operatorname{Ass}_R M,
 \qquad (0:_M I)\ne0.
\]
For a nonzero finite module over a local ring, depth zero is therefore
equivalent to a nonzero socle $(0:_M\mathfrak m)$.
\end{definition}
\begin{proof}
When $IM\ne M$, any injective multiplication by $x\in I$ gives a
regular sequence of length one, since $M/xM$ surjects onto $M/IM$.
The zero-divisor theorem and the finiteness of the set of associated
primes \cite[Theorems ALG-05-C02-S03 and ALG-05-C02-R03]{p103}
identify the set of failed choices with a finite union of primes.
Prime avoidance then gives the middle condition if every $x\in I$
is a zero divisor on $M$. Conversely an element of $M$ whose
annihilator is an associated prime containing $I$ is a nonzero
element of $(0:_M I)$. Such an element prevents every $x\in I$ from
acting injectively. In a local ring $M\ne0$ ensures $\mathfrak mM\ne M$
by Nakayama's lemma.
\end{proof}

\begin{proposition}[Dimension and bounded lengths]\label{res:DP-02}
Over a Noetherian local ring $(R,\mathfrak m)$, if $M\ne0$ is finite
and $x\in\mathfrak m$ is a nonzerodivisor on $M$, then
\[
 \dim\operatorname{Supp}(M/xM)=\dim\operatorname{Supp}(M)-1.
\]
In particular $\operatorname{depth}_R M\le\dim\operatorname{Supp}M$.
Over an arbitrary Noetherian ring the lengths of $M$-regular sequences
in $I$ are bounded whenever $IM\ne M$.
\end{proposition}
\begin{proof}
The support of $M/xM$ is $\operatorname{Supp}M\cap V(x)$.
Indeed localization and Nakayama give nonzero fibers exactly there.
Every minimal prime of $\operatorname{Supp}M$ is associated
\cite[Theorem ALG-05-C02-R03]{p103}, so $x$ avoids all of them.
A chain in $\operatorname{Supp}(M/xM)$ can therefore be extended
strictly downward by adjoining a minimal prime of
$\operatorname{Supp}M$ below its first member. This gives the upper
bound on dimension. The opposite bound follows by prepending $x$
to a system of parameters of $M/xM$
\cite[Theorem ALG-05-C03-S01]{p103}. Repeating the dimension drop proves
the depth bound. If $IM\ne M$ over a general Noetherian ring, choose
$\mathfrak p\in\operatorname{Supp}(M/IM)$. A regular sequence in $I$
remains regular after localization at $\mathfrak p$: localization
preserves each injection and the terminal quotient still surjects
onto $(M/IM)_{\mathfrak p}\ne0$. Its length is bounded by
$\dim\operatorname{Supp}_{R_{\mathfrak p}}M_{\mathfrak p}$, which is
finite by the parameter theorem for the Noetherian local ring.
\end{proof}

\begin{lemma}[Hom detection and the two scalar actions on Ext]
\label{res:DP-03}
For finite modules $N,M$ over a Noetherian ring,
$\operatorname{Hom}_R(N,M)\ne0$ if and only if
$\operatorname{Supp}N\cap\operatorname{Ass}M\ne\varnothing$.
For modules over a commutative ring the scalar action on
$\operatorname{Ext}^i_R(N,M)$ can be computed by precomposition
with multiplication on $N$ or by postcomposition with multiplication
on $M$. In particular $\operatorname{Ann}_R N$ acts by zero.
\end{lemma}
\begin{proof}
A nonzero image of a homomorphism is a quotient of $N$ and a submodule
of $M$. It has an associated prime, which belongs to both sets in the
intersection. Conversely, if $\mathfrak p$ lies in the intersection,
$N_{\mathfrak p}$ has a quotient isomorphic to $\kappa(\mathfrak p)$
by Nakayama. The associated-prime witness embeds
$\kappa(\mathfrak p)$ into $M_{\mathfrak p}$. The composite is a
nonzero localized Hom and hence comes from a nonzero Hom group,
by localization of Hom for finitely presented $N$
\cite[Theorems ALG-05-C02-R01 and ALG-10-C03-S03]{p103}.

Take a projective resolution $P_\bullet\to N$. Multiplication by
$a\in R$ on $P_\bullet$ lifts multiplication on $N$. Precomposition
with this chain map and postcomposition by multiplication on $M$
are literally the same on $\operatorname{Hom}_R(P_\bullet,M)$.
If $aN=0$, the chain map lifts the zero map and is homotopic to zero
by the projective comparison theorem
\cite[Theorem ALG-10-C02-R02]{p103}. Thus the action in cohomology is
zero.
\end{proof}

\begin{theorem}[Regular sequences and Ext detect the same grade]
\label{res:DP-04}
Let $R$ be Noetherian and $N,M$ finite. Then
\[
 \operatorname{grade}_R(\operatorname{Ann}_RN,M)
   =\inf\{i\ge0:\operatorname{Ext}^i_R(N,M)\ne0\},
 \qquad \inf\varnothing=+\infty.
\]
Taking $N=R/I$ gives the Ext criterion for grade. In the local case,
taking $N=k$ gives the Ext criterion for depth. When $IM\ne M$,
every maximal $M$-regular sequence in $I$ has length
$\operatorname{grade}_R(I,M)$.
\end{theorem}
\begin{proof}
Write $I=\operatorname{Ann}_R N$. If $IM=M$, then
$\operatorname{Supp}M$ and $V(I)$ are disjoint, since a prime in the
intersection would contradict local Nakayama. As $\operatorname{Supp}N=V(I)$,
at least one of $M_{\mathfrak p},N_{\mathfrak p}$ is zero at every prime.
Ext commutes with localization by
\cite[Theorem ALG-10-C03-S03]{p103}, so every Ext group is zero.

Otherwise, choose a maximal $M$-regular sequence in $I$. The length
bound in \ref{res:DP-02} ensures the successive choices terminate.
On the terminal quotient $L$ there is no regular element of $I$.
By \ref{res:DP-01} and \ref{res:DP-03},
$\operatorname{Hom}_R(N,L)\ne0$. At each preceding step let $U$ be
the current module, $x\in I$ its chosen nonzerodivisor, and $V=U/xU$.
The Ext sequence for $0\to U\xrightarrow{x}U\to V\to0$
has all its multiplication-by-$x$ arrows zero by \ref{res:DP-03}.
Also $\operatorname{Hom}_R(N,U)=0$, because a map from $N$ has image
killed by $x$ whereas $x$ is injective on $U$. Thus for every $i\ge0$
there is an exact sequence
\[
 0\longrightarrow \operatorname{Ext}_R^i(N,U)
 \longrightarrow \operatorname{Ext}_R^i(N,V)
 \longrightarrow \operatorname{Ext}_R^{i+1}(N,U)
 \longrightarrow0 .
\]
If the first nonzero term for $V$ is in degree $j$, induction along
these short exact sequences shows that the first nonzero term
for $U$ is in degree $j+1$. Starting from $L$, we obtain the claimed
threshold equal to the chosen length. Applying this argument with
$N=R/I$ to every maximal sequence gives the common length. Every
regular sequence in $I$ extends to a maximal one, by the finite
length bound, so this length equals the supremum in the definition.
\end{proof}

For later use, the same connecting maps give
$\operatorname{Ext}_R^i(N,M)=0$ for $i<r$ and a natural isomorphism
$\operatorname{Ext}_R^r(N,M)\simeq\operatorname{Hom}_R(N,M/JM)$
whenever $x_1,\ldots,x_r$ is $M$-regular in $\operatorname{Ann}_R N$
and $J=(x_1,\ldots,x_r)$. Indeed, start with the vanishing of Hom
before each regular step, then successively use the short exact
cohomology sequences above to shift by one. This assertion does
not require the given sequence to be maximal.

\begin{corollary}[Regular quotient and the type of a module]\label{res:DP-05}
If $x_1,\ldots,x_r\in I$ is $M$-regular and $IM\ne M$, then with
$J=(x_1,\ldots,x_r)$,
\[
 \operatorname{grade}_{R/J}(I/J,M/JM)
 =\operatorname{grade}_R(I,M/JM)
 =\operatorname{grade}_R(I,M)-r .
\]
In a local ring this gives
$\operatorname{depth}_{R/J}(M/JM)=\operatorname{depth}_R M-r$
when the sequence lies in $\mathfrak m$. More generally, a quotient
of the base ring by an ideal that annihilates the module preserves
grade with respect to the image of the given ideal.

For $M\ne0$ over $(R,\mathfrak m,k)$, put
$d=\operatorname{depth}_R M$. Its \emph{type}
$\dim_k\operatorname{Ext}_R^d(k,M)$ is the dimension over $k$ of the
socle of the quotient by any maximal $M$-regular sequence in
$\mathfrak m$.
\end{corollary}
\begin{proof}
Actions by sequences on a module killed by $J$ depend only on their
images in $R/J$, and each sequence in the image ideal admits lifts to
$I$. The successive quotients and the final nonzero condition are
identical. Equality of grade over the two base rings follows, including
the infinite convention. Extend the given sequence to a maximal one.
In the quotient the possible extensions are precisely the regular
sequences in the image of $I$, so \ref{res:DP-04} gives the length drop.
For the type, the exact sequence in the proof of \ref{res:DP-04} for
$N=k$ and a regular $x\in\mathfrak m$ identifies
$\operatorname{Ext}_R^{d-1}(k,M/xM)$ with
$\operatorname{Ext}_R^d(k,M)$, since the preceding term is zero.
Iterating over the sequence identifies the latter group with
$\operatorname{Hom}_R(k,M/JM)$. Each comparison is $k$-linear
because $\mathfrak m$ annihilates the Ext groups by \ref{res:DP-03}.
These groups are finite-dimensional: use a resolution of $k$ by
degreewise finite free modules and Noetherianity to see that their
cohomology modules are finite.
\end{proof}

\begin{proposition}[Local detection of grade]\label{res:DP-06}
For $R$ Noetherian and $M$ finite,
\[
 \operatorname{grade}_R(I,M)=
 \inf_{\mathfrak p\in V(I)\cap\operatorname{Supp}M}
 \operatorname{depth}_{R_{\mathfrak p}}M_{\mathfrak p}.
\]
Consequently grade depends only on $\sqrt I$. For ideals $I,J$,
the grades of $IJ$ and of $I\cap J$ both equal the minimum of the
grades of $I$ and $J$ on $M$. For a proper ideal $I$,
$\operatorname{grade}_R(I,R)\le\operatorname{ht}I$.
\end{proposition}
\begin{proof}
If $IM=M$, the indexing set is empty by local Nakayama.
If $IM\ne M$, take a maximal $M$-regular sequence in $I$ of length
$r$. For every prime in the indexing set the localized sequence is
regular, so the local depth is at least $r$. Its terminal quotient
$L$ has $I$-grade zero. By \ref{res:DP-01} take
$\mathfrak p\in\operatorname{Ass}L$ containing $I$. At this prime the
terminal quotient has a nonzero socle, by localization of the
associated-prime witness. Thus the depth is exactly $r$ before
quotienting, by \ref{res:DP-05}. This proves the equality. The last
claims use $V(I)=V(\sqrt I)$ and $V(IJ)=V(I\cap J)=V(I)\cup V(J)$.
The height bound follows from \ref{res:DP-02}, since
$\dim R_{\mathfrak p}=\operatorname{ht}\mathfrak p$ and height of
$I$ is the infimum of these heights over $V(I)$.
\end{proof}

\section{Exact sequences and changes of rings}

\begin{proposition}[The depth lemma]\label{res:DP-07}
Fix an ideal $I$ in a Noetherian ring $R$. In an exact sequence
$0\to L\to M\to N\to0$ of finite modules write
$a=\operatorname{grade}_R(I,L)$, $b=\operatorname{grade}_R(I,M)$
and $c=\operatorname{grade}_R(I,N)$. Then
\[
 b\ge\min(a,c),\qquad
 a\ge\min(b,c+1),\qquad
 c\ge\min(b,a-1).
\]
Subtraction by one at zero is interpreted as zero, and
$+\infty\pm1=+\infty$. The same holds for local depth.
If $a<c$, then $b=a$. If $a>c+1$, then $b=c$.
If $a<b$, then $a\ge1$ and $c=a-1$.
\end{proposition}
\begin{proof}
Use the long exact sequence of Ext with first variable $R/I$
\cite[Theorem ALG-10-C03-R02]{p103}.
In degree $i$ each inequality follows from the vanishing of
the two displayed neighbors of the corresponding group:
\[
 \begin{aligned}
  &\operatorname{Ext}^i(R/I,L) \to
      \operatorname{Ext}^i(R/I,M) \to \operatorname{Ext}^i(R/I,N),\\
  &\operatorname{Ext}^{i-1}(R/I,N) \to
      \operatorname{Ext}^i(R/I,L) \to \operatorname{Ext}^i(R/I,M),\\
  &\operatorname{Ext}^i(R/I,M) \to
      \operatorname{Ext}^i(R/I,N) \to \operatorname{Ext}^{i+1}(R/I,L).
 \end{aligned}
\]
In the second line the preceding term is zero for $i=0$ because Hom
is left exact. If $a<c$, the first two inequalities force $b=a$:
the first gives $b\ge a$, and $c+1>a$ makes the second impossible
if $b>a$. If $a>c+1$, the first forces $b\ge c$, whereas the third
would fail for $b>c$. Finally, if $a<b$,
the second inequality forces $c+1\le a$ while the third gives
$c\ge a-1$. This also forces $a\ge1$ and proves the last claim.
\end{proof}

\begin{proposition}[Flat and finite changes of rings]\label{res:DP-08}
Let $R,S$ be Noetherian and $I\subseteq R$.
If $R\to S$ is flat and $M$ is finite over $R$, then
\[
 \operatorname{grade}_R(I,M)\le
 \operatorname{grade}_S(IS,S\otimes_R M).
\]
For faithfully flat maps this is an equality.
If $S$ is a finite $R$-algebra and $N$ is a finite $S$-module, then
$\operatorname{grade}_R(I,N)=\operatorname{grade}_S(IS,N)$.
In particular, when $R$ is local, $\operatorname{depth}_R N$ is
the minimum of the depths of $N_{\mathfrak n}$ over the local rings
$S_{\mathfrak n}$ at the maximal ideals of $S$.
\end{proposition}
\begin{proof}
For the flat assertion the Ext base-change theorem
\cite[Theorem ALG-10-C03-S03]{p103} applies to the finite first
argument $R/I$. It identifies the Ext groups on the right with
$S\otimes_R\operatorname{Ext}_R^i(R/I,M)$, which gives the
inequality by \ref{res:DP-04}. Faithful flatness detects zero
modules, giving equality of the thresholds in that case.

For the finite assertion, $IN=(IS)N$, and a list from $I$ has
exactly the same multiplication maps in either interpretation.
If $(IS)N=N$ both grades are infinite. Otherwise, whenever the
$IS$-grade of a successive quotient $U$ is positive, no prime in
$\operatorname{Ass}_S U$ contains $IS$. The contractions of these
finitely many primes to $R$ therefore do not contain $I$. Choose
$x\in I$ avoiding their union by prime avoidance. Its image is a
nonzerodivisor on $U$. Repeat until the $IS$-grade becomes zero,
decreasing it by one at each step using \ref{res:DP-05}. The chosen
list shows that the $R$-grade is at least the $S$-grade, and the
reverse inequality follows by viewing every $I$-sequence as an
$IS$-sequence. For local $R$ with maximal ideal $\mathfrak m$,
the primes of $S$ over $\mathfrak m$ are precisely its maximal ideals:
$S$ is integral over the image of $R$, and $S/\mathfrak m S$ is
finite over the residue field. Apply \ref{res:DP-06} to $IS=\mathfrak m S$.
\end{proof}

\begin{theorem}[The Koszul vanishing threshold]\label{res:DP-09}
Let $R$ be Noetherian, let $M$ be finite, and let
$I=(a_1,\ldots,a_n)$. Set
$K^i(\mathbf a;M)=\operatorname{Hom}_R(K_i(\mathbf a;R),M)$
with the differential induced by the Koszul differential.
Then
\[
 \operatorname{grade}_R(I,M)
  =\inf\{i:H^i(K^\bullet(\mathbf a;M))\ne0\}
  =n-\sup\{j:H_j(K_\bullet(\mathbf a;M))\ne0\}.
\]
When both complexes are acyclic the right-hand thresholds are
interpreted as $+\infty$. In particular, if $M/IM\ne0$ then
$\operatorname{grade}_R(I,M)\le n$. The cohomology groups depend on
the chosen Koszul complex, but their first nonvanishing degree
is independent of a generating list.
\end{theorem}
\begin{proof}
The wedge homotopies of \ref{res:RK-06} and \ref{res:RK-10}
show that $I$ acts by zero on cohomology. The kernel in degree zero
is $(0:_M I)$, by the formula for the first differential.
For $IM\ne M$ and grade zero this is nonzero by \ref{res:DP-01}.
If $IM=M$, choose a finite generating list of $M$ and write it as
a matrix with entries in $I$ times itself. Multiplying by the
adjugate of $1$ minus that matrix gives an element $1-a$ with
$a\in I$ killing $M$. The operator $a$ thus acts as the identity
on the entire Koszul complex with coefficients in $M$, while its
homological action is zero. This gives acyclicity. The zero-module
case is included.

If the grade is positive and finite, choose a nonzerodivisor $x\in I$.
The short exact sequence of coefficient modules gives the Koszul
cohomology sequence by \ref{res:RK-12}. Multiplication by $x$
induces zero on cohomology, so the sequence for successive degrees
breaks up into
\[
 0\to H^i(K^\bullet(\mathbf a;M))\to
 H^i(K^\bullet(\mathbf a;M/xM))\to
 H^{i+1}(K^\bullet(\mathbf a;M))\to0 .
\]
Also $H^0(K^\bullet(\mathbf a;M))=0$ by injectivity of $x$.
The same vanishing-threshold argument as in \ref{res:DP-04},
together with the grade drop \ref{res:DP-05}, proves the first
equality. Self-duality \ref{res:RK-10} identifies cohomology in
degree $i$ with homology in degree $n-i$ using the ordered basis
of the free module defining Koszul. The last bound follows because
the complex has no terms above degree $n$.
\end{proof}

\section{Localization and homological dimensions}

\begin{theorem}[Depth bounds at associated primes and after localization]
\label{res:DP-10}
Let $(R,\mathfrak m)$ be Noetherian local and $M\ne0$ finite. Then
\[
 \begin{aligned}
 \operatorname{depth}_R M &\le \dim(R/\mathfrak p)
           &&(\mathfrak p\in\operatorname{Ass}_R M),\\
 \operatorname{depth}_R M &\le \operatorname{depth}_{R_{\mathfrak p}}M_{\mathfrak p}
                                      +\dim(R/\mathfrak p)
           &&(\mathfrak p\in\operatorname{Spec}R).
 \end{aligned}
\]
Depth itself may either decrease or increase under localization.
\end{theorem}
\begin{proof}
Use induction on the depth for the first inequality. At depth zero
it holds. At positive depth choose a nonzerodivisor $x\in\mathfrak m$.
Fix $\mathfrak p\in\operatorname{Ass}M$, so $x\notin\mathfrak p$,
and regard $U=R/\mathfrak p$ as a submodule of $M$. Artin--Rees
\cite[Theorem ALG-05-C04-R02]{p103} gives
$U\cap x^s M\subseteq xU$ for some $s\ge1$. Its image
$\overline U$ in $M/x^sM$ is killed by $\mathfrak p+(x^s)$
and surjects onto $U/xU$. Thus
$\operatorname{Supp}\overline U=V(\mathfrak p+(x))$.
Any prime $\mathfrak q$ minimal in this set is associated to
$\overline U$, hence to $M/x^sM$. It strictly contains $\mathfrak p$,
so any chain from $\mathfrak q$ extends downward through
$\mathfrak p$ and $\dim(R/\mathfrak q)\le\dim(R/\mathfrak p)-1$.
The depth of $M/x^sM$ is one less by \ref{res:DP-05}.
Induction proves the first bound.

For the localization inequality the result is immediate if
$M_{\mathfrak p}=0$ or
$\operatorname{depth}_R M\le\dim(R/\mathfrak p)$.
Otherwise $\mathfrak p$ cannot be contained in an associated
prime $\mathfrak q$ of $M$, since that would give
$\dim(R/\mathfrak q)\le\dim(R/\mathfrak p)$ in the first bound.
Choose a nonzerodivisor $x\in\mathfrak p$ by prime avoidance.
Locally at $\mathfrak p$ it is still a nonzerodivisor and its
quotient is nonzero. The depths of $M$ and $M_{\mathfrak p}$
both drop by one on taking the quotient. Induction on
$\operatorname{depth}_R M$ completes the inequality.

For example, $R=k[x,y]_{(x,y)}$ has depth two, with $x,y$ a
regular sequence and $\dim R=2$. Its localization at $(x)$ has
depth one, since it is $k(y)[x]_{(x)}$. For the opposite direction,
$R\oplus k$ has depth zero at $(x,y)$ because of the $k$ summand.
At $(x)$ this summand vanishes and the localized module has depth one.
\end{proof}

\begin{corollary}[Two Ext-vanishing bounds]\label{res:DP-11}
If $R$ is Noetherian local and $N,M$ are nonzero finite modules, then
\[
 \operatorname{Ext}^i_R(N,M)=0
   \qquad (0\le i<\operatorname{depth}_R M-\dim\operatorname{Supp}N).
\]
Over any Noetherian ring, let $E$ be finite, put
$g=\operatorname{grade}_R(\operatorname{Ann}_R E,R)$, and suppose
$N$ is finite of finite projective dimension $\ell$. Then
$\operatorname{Ext}^i_R(E,N)=0$ for $0\le i<g-\ell$,
with the infinite convention when $g=+\infty$.
In particular
$\operatorname{grade}_R(\operatorname{Ann}_R E,R)\le
\operatorname{pd}_R E$ for $E\ne0$.
\end{corollary}
\begin{proof}
For the first assertion, \ref{res:DP-04} and \ref{res:DP-06}
give the threshold as the infimum of the local depths for
$\mathfrak p\in\operatorname{Supp}N\cap\operatorname{Supp}M$.
For such a prime $V(\mathfrak p)\subseteq\operatorname{Supp}N$,
so $\dim(R/\mathfrak p)\le\dim\operatorname{Supp}N$.
Apply \ref{res:DP-10} and take the infimum.

For the second assertion, a finite projective module is a direct
summand of a finite free module. In the projective case the claim
therefore follows from \ref{res:DP-04} applied to $E,R$.
For $\ell>0$ choose $0\to U\to F\to N\to0$ with $F$ finite free.
Then $U$ is finite by Noetherianity and has projective dimension at
most $\ell-1$: dimension shifting
\cite[Theorem ALG-10-C02-R04]{p103} reduces eventual vanishing
to the projectivity test proved below in \ref{res:DP-12}.
For $i<g-\ell$ the group $\operatorname{Ext}^i_R(E,F)$ vanishes by
the free case. The group $\operatorname{Ext}^{i+1}_R(E,U)$ vanishes by
the induction hypothesis. Exactness gives the result.
For the last bound, finite projective dimension of $E$ forces
$\operatorname{Ext}_R^i(E,R)=0$ above that dimension. Since $E\ne0$,
its annihilator is proper, so \ref{res:DP-02} and \ref{res:DP-04}
give a finite first nonzero degree for this Ext group. If
$\operatorname{pd}_R E$ is infinite the bound holds as well.
\end{proof}

\begin{proposition}[Minimal free resolutions and finite projective dimension]
\label{res:DP-12}
Let $(R,\mathfrak m,k)$ be Noetherian local, and let $M$ be finite.
There is a free resolution $F_\bullet\to M$ with every $F_i$ finite
free and the image of $F_{i+1}\to F_i$ contained in $\mathfrak m F_i$.
For such a resolution
\[
 k\otimes_R F_i\simeq \operatorname{Tor}^R_i(k,M).
\]
If $M\ne0$ has finite projective dimension $e$, then the last nonzero
term of this minimal free resolution is $F_e$.
\end{proposition}
\begin{proof}
Choose a basis of $M/\mathfrak mM$ and lift it to $M$.
Nakayama gives a surjection from a finite free module $F_0$,
and its kernel is contained in $\mathfrak mF_0$ because the reduction
of the covering map is an isomorphism. The kernel is finite, so the
construction iterates. The module maps after tensoring with $k$ are
zero, giving the displayed formula by the computation of Tor using
free resolutions and balance
\cite[Results ALG-10-C03-R01 and ALG-10-C03-S01]{p103}.

Here is the projectivity test used for finite bounds. Over any ring,
a module $T$ with $\operatorname{Ext}_R^1(T,V)=0$ for every module
$V$ is projective. For a free surjection $P\to T$ with kernel $U$,
the identity of $T$ lifts to a section by applying
$\operatorname{Hom}_R(T,-)$ and using
$\operatorname{Ext}_R^1(T,U)=0$. Thus $T$ is a direct summand of $P$.
It follows by dimension shifting that vanishing of
$\operatorname{Ext}_R^{b+1}(M,-)$ makes the $b$th syzygy projective,
and hence bounds the projective dimension of $M$ by $b$.

Over the local ring a finite projective module is free: take its
minimal free surjection as above. The split kernel has zero reduction
modulo $\mathfrak m$ because the surjection reduces to an isomorphism,
so that kernel is zero by Nakayama. If $M$ has finite projective
dimension $e$, the test and dimension shifting make its $e$th syzygy
finite projective and thus free. Its minimal free cover is an
isomorphism, so the resolution stops there. It cannot stop earlier,
by the definition of projective dimension.
\end{proof}

\begin{theorem}[Auslander--Buchsbaum]\label{res:DP-13}
Let $(R,\mathfrak m,k)$ be Noetherian local and let $M$ be a nonzero
finite module of finite projective dimension $e$. Then
\[
 \operatorname{pd}_R M+\operatorname{depth}_R M
   =\operatorname{depth}_R R.
\]
If $x\in\mathfrak m$ is a nonzerodivisor on $M$, then
$\operatorname{pd}_R(M/xM)=e+1$. If $x$ is also a nonzerodivisor
on $R$, then $\operatorname{pd}_{R/(x)}(M/xM)=e$.
\end{theorem}
\begin{proof}
Write $d=\operatorname{depth}_R R$. At $e=0$ the module is nonzero
finite free by \ref{res:DP-12}, so its first nonzero residue Ext is
in degree $d$. For $e>0$ take the first step
$0\to U\to F_0\to M\to0$ of a minimal resolution. The tail is
minimal for $U$, which is nonzero of projective dimension $e-1$
by \ref{res:DP-12} and its Tor detection. Induction gives
$t=\operatorname{depth}_R U=d-(e-1)$. If $e>1$, then $t<d$.
In this case $t=0$ is impossible, because the exact Hom sequence
injects $\operatorname{Hom}_R(k,U)$ into
$\operatorname{Hom}_R(k,F_0)=0$. The first nonzero residue Ext of
$U$ is therefore in degree $t\ge1$. In degrees $0\le i<t-1$,
both $\operatorname{Ext}_R^i(k,F_0)$ and
$\operatorname{Ext}_R^{i+1}(k,U)$ vanish. At $i=t-1$ the connecting
homomorphism identifies $\operatorname{Ext}_R^{t-1}(k,M)$ with
$\operatorname{Ext}_R^t(k,U)$, since the free terms in both
degrees are zero. Thus $\operatorname{depth}_R M=t-1=d-e$.

For $e=1$, $U$ is finite free and every matrix entry of $U\to F_0$
lies in $\mathfrak m$. The induced map
$\operatorname{Ext}^d_R(k,U)\to\operatorname{Ext}^d_R(k,F_0)$
is zero, since its entries act as zero by \ref{res:DP-03}.
If $d=0$, this contradicts injectivity on Hom and the nonzero
source. If $d\ge1$, the same exact-sequence calculation now gives
vanishing below $d-1$ and an isomorphism
$\operatorname{Ext}_R^{d-1}(k,M)\simeq\operatorname{Ext}_R^d(k,U)\ne0$.
This covers the boundary case.

For the first quotient assertion use the long exact Tor sequence
for $0\to M\xrightarrow{x}M\to M/xM\to0$ with $k$ fixed.
Multiplication by $x$ acts as zero on $\operatorname{Tor}_i^R(k,M)$:
on a free resolution of $M$ it is represented by multiplication by
$x$, which is zero after tensoring with $k$. Exactness therefore
gives zero Tor in degrees above $e+1$ and an isomorphism
$\operatorname{Tor}_{e+1}^R(k,M/xM)\simeq\operatorname{Tor}_e^R(k,M)\ne0$.
The minimal free resolution of $M/xM$ has last nonzero term in
degree $e+1$ by \ref{res:DP-12}. Its length is the projective dimension.

For the other quotient assertion, a two-term free resolution of $R/(x)$
is $R\xrightarrow{x}R$ because $x$ is a nonzerodivisor on $R$.
Tensoring with $M$ gives zero positive Tor because $x$ is also
injective on $M$. Thus the free resolution of $M$ stays exact in
positive degrees on reduction modulo $x$, and resolves $M/xM$.
It is still minimal and its last nonzero term is still nonzero
modulo the maximal ideal. The Tor test over $R/(x)$ gives the claimed
projective dimension.
\end{proof}

\section{Two interfaces for quotient and restriction arguments}

\begin{proposition}[Change of rings by a nonzerodivisor]\label{res:DP-14}
Let $R$ be commutative, let $x$ be a nonzerodivisor on both $R$ and
an $R$-module $M$, put $S=R/(x)$, and let $N$ be an $S$-module.
Then $\operatorname{Hom}_R(N,M)=0$ and for all $j\ge0$ there are
natural isomorphisms
\[
 \operatorname{Ext}^{j+1}_R(N,M)
    \simeq \operatorname{Ext}_S^j(N,M/xM).
\]
For a sequence regular on both $R$ and $M$, the result iterates
with one degree shift per element for a module $N$ killed by the
whole ideal of the sequence.
\end{proposition}
\begin{proof}
The Hom group vanishes because $x$ kills $N$ but is injective on $M$.
Take an injective resolution $M\to E^\bullet$ over $R$ and set
$C^\bullet=\operatorname{Hom}_R(S,E^\bullet)$.
Each $C^i$ is injective over $S$: for an $S$-module $T$, evaluation
identifies $\operatorname{Hom}_S(T,C^i)$ with
$\operatorname{Hom}_R(T,E^i)$, an exact functor of $T$ because
restriction of scalars is exact. By the two-term free resolution
$0\to R\xrightarrow{x}R\to S\to0$, the cohomology of $C^\bullet$
is $M/xM$ in degree one and zero elsewhere, using balanced
computation of Ext \cite[Theorem ALG-10-C03-S01]{p103}.

The differential $C^0\to C^1$ is injective and splits because
$C^0$ is injective over $S$. Hence
$C^1/\operatorname{im}(C^0)$, followed by $C^2,C^3,\ldots$, is an
injective resolution of $M/xM$ when started in degree zero.
The original complex is the sum of this shifted resolution and
a two-term contractible complex. Applying $\operatorname{Hom}_S(N,-)$
and taking cohomology therefore gives the displayed isomorphisms
via the adjunction in each degree. Equivalently, use the canonical
quotient by the contractible two-term subcomplex, which shows that
the comparison is natural. Iteration applies the same argument
to the successive rings and modules.
\end{proof}

\begin{corollary}[Depth two and extension of linear maps]\label{res:DP-15}
For $R$ Noetherian, $I\subset R$ and $M$ a finite module,
$\operatorname{grade}_R(I,M)\ge2$ if and only if
the restriction map $\operatorname{Hom}_R(R,M)\to\operatorname{Hom}_R(I,M)$
is an isomorphism. Here $\operatorname{Hom}_R(R,M)$ is identified
with $M$ by evaluation at one.
\end{corollary}
\begin{proof}
The Ext exact sequence in the first variable for
$0\to I\to R\to R/I\to0$ identifies the kernel and cokernel of
the restriction with $\operatorname{Hom}_R(R/I,M)$ and
$\operatorname{Ext}_R^1(R/I,M)$, respectively, since $R$ is free.
Both groups vanish precisely when the grade is at least two,
by \ref{res:DP-04}.
\end{proof}
\chapter{Cohen--Macaulay Modules and Rings}
\label{ch:cohen-macaulay}

In a Cohen--Macaulay module, dimension measures the length of a
regular sequence. This equality makes parameters, associated primes,
and localizations accessible through the depth theory of Chapter~\ref{ch:depth}.
We use the local dimension and completion results of \cite{p103} and
compare the formulations in \cite[Chapter 18]{eisenbud},
\cite[Sections 2.1--2.3]{brunsherzog}, and \cite[Section 17]{matsumura}.

\section{Parameters and regular quotients}

\begin{definition}[Cohen--Macaulay modules and rings]\label{res:CM-01}
Let $(R,\mathfrak m)$ be a Noetherian local ring. For a nonzero
finite $R$-module $M$, set
$d_R(M)=\dim\operatorname{Supp}_R M=\dim(R/\operatorname{Ann}_R M)$.
The module is \emph{Cohen--Macaulay} if
$\operatorname{depth}_R M=d_R(M)$. It is \emph{maximal
Cohen--Macaulay} if $\operatorname{depth}_R M=\dim R$. A ring is
Cohen--Macaulay at a prime if its local ring there is Cohen--Macaulay
as a module over itself. Over a Noetherian ring, a finite module
is Cohen--Macaulay if it has that property at every prime in its
support. The zero module satisfies the local and global
Cohen--Macaulay conditions by convention, and is admitted along
with the nonzero maximal Cohen--Macaulay modules when discussing
syzygies. Numerical assertions about the dimension of a nonzero
module refer to the part of its support where it is nonzero.
\end{definition}

\begin{proposition}[Support and parameter sequences]\label{res:CM-02}
Let $M\ne0$ be a finite module over a Noetherian local ring
$(R,\mathfrak m)$ with $d=d_R(M)$. If $J\subseteq\mathfrak m$, then
\[
 \operatorname{Supp}(M/JM)=\operatorname{Supp}(M)\cap V(J).
\]
The module $M/JM$ has finite length if and only if this support
consists of $\mathfrak m$. A \emph{system of parameters} for $M$ is
a sequence in $\mathfrak m$ of length $d$ with finite-length
quotient. Such sequences exist and $d$ is the minimum length
of any sequence in $\mathfrak m$ with finite-length quotient.
If $x\in\mathfrak m$ acts injectively on $M$, then
$d_R(M/xM)=d-1$.
\end{proposition}
\begin{proof}
At $\mathfrak p$, localization makes the quotient zero if either
$M_{\mathfrak p}=0$ or $J\not\subseteq\mathfrak p$.
In the remaining case local Nakayama gives a nonzero quotient.
The finite-length and minimum-length statements are
\cite[Theorem ALG-05-C03-S01]{p103}; therein the minimum length
is proved by the height theorem and prime avoidance in
$R/\operatorname{Ann} M$. In particular
$d_R(M/xM)\ge d-1$. A nonzerodivisor on $M$ avoids its
associated primes, including those minimal in its support
\cite[Theorems ALG-05-C02-S03 and ALG-05-C02-R03]{p103}. Thus the first prime in
any chain in $\operatorname{Supp}(M/xM)$ is not minimal in
$\operatorname{Supp}(M)$. Prepending a minimal prime of the support
increases the length by one, proving $d_R(M/xM)\le d-1$.
\end{proof}

\begin{proposition}[Associated primes and equidimension]\label{res:CM-03}
For a nonzero finite Cohen--Macaulay module $M$ over a Noetherian
local ring, every associated prime $\mathfrak p$ satisfies
$\dim R/\mathfrak p=d_R(M)$ and is minimal in the support. All
minimal primes of the support are associated, so $M$ has
equidimensional support and no embedded associated primes.
\end{proposition}
\begin{proof}
The associated-prime depth bound in \ref{res:DP-10} gives
\[
 d_R(M)=\operatorname{depth}_R M\leq \dim R/\mathfrak p
 \leq d_R(M).
\]
If $\mathfrak p$ were not minimal, a longest chain from $\mathfrak p$
to the maximal ideal could be extended downward in the support.
This contradicts the equality. The assertion about the minimal
primes is \cite[Theorem ALG-05-C02-R03]{p103}.
\end{proof}

\begin{theorem}[Initial segments of systems of parameters]\label{res:CM-04}
Let $M\ne0$ be finite and Cohen--Macaulay over $(R,\mathfrak m)$,
and write $d=d_R(M)$. For $x_1,\ldots,x_c\in\mathfrak m$,
the sequence is $M$-regular if and only if its quotient has
dimension $d-c$. In either case it extends to a system of
parameters and every intermediate quotient is Cohen--Macaulay.
For a nonzero finite module, being Cohen--Macaulay is equivalent
to some system of parameters being regular, and equally to every
system of parameters being regular.
\end{theorem}
\begin{proof}
For a regular sequence each step drops dimension by one by
\ref{res:CM-02} and depth by one by \ref{res:DP-05}. Thus each
quotient is again Cohen--Macaulay. Conversely the dimension
can fall by at most one at each step of an arbitrary sequence
in $\mathfrak m$, by \ref{res:CM-02}. All intermediate quotients
are nonzero by Nakayama. If the final dimension is $d-c$, each
step must drop it by one. An element such that the dimension
drops avoids the associated primes of a Cohen--Macaulay module:
if it belonged to one of them, the quotient support would contain
that prime and would still have dimension $d$ by \ref{res:CM-03}.
Avoiding them makes multiplication injective by
\cite[Theorem ALG-05-C02-S03]{p103}. Apply this argument successively
to the intermediate quotients, using the regular case to retain
the Cohen--Macaulay property.

To extend the given sequence, lift a system of parameters of the
final quotient, which exists by \ref{res:CM-02}. The result has
length $d$ and finite-length quotient. A regular sequence of
length $d$ forces $\operatorname{depth} M\ge d$, and
\ref{res:DP-02} gives the reverse inequality. This proves
the last two characterizations, including $d=0$ and the
empty system of parameters.
\end{proof}

\begin{corollary}[Regular quotients and quotient scalars]\label{res:CM-05}
If $M$ is nonzero and finite over a Noetherian local ring $R$
and $x$ in its maximal ideal is $M$-regular, then $M$ is
Cohen--Macaulay if and only if $M/xM$ is. This equivalence holds
over $R$ or, for the quotient module, over $R/(x)$.
More generally a finite module over a surjective local map
$R\twoheadrightarrow S$ has the same dimension and depth
when viewed over either ring.
\end{corollary}
\begin{proof}
For a finite $S$-module the supports correspond under the
prime correspondence for $S=R/\ker(R\to S)$, hence have the
same dimension. A sequence in the maximal ideal of $R$ acts
on each successive quotient through its image in $S$, and
every sequence in the maximal ideal of $S$ lifts. Thus the
regular-sequence depth agrees as well. For the regular quotient,
combine the dimension and depth drops in
\ref{res:CM-02} and \ref{res:DP-05}.
\end{proof}

\section{Localization, heights, and regular rings}

\begin{theorem}[Localization and the local dimension equality]
\label{res:CM-06}
If a nonzero finite module $M$ over $(R,\mathfrak m)$ is
Cohen--Macaulay, then $M_{\mathfrak p}$ is Cohen--Macaulay for
each $\mathfrak p\in\operatorname{Supp}(M)$, and
\[
 d_R(M)=d_{R_{\mathfrak p}}(M_{\mathfrak p})+\dim R/\mathfrak p.
\]
Over a Noetherian ring it therefore suffices to test
Cohen--Macaulayness at maximal ideals, and the property persists
under arbitrary localization.
\end{theorem}
\begin{proof}
A chain in the localized support of maximum possible length,
ending at $\mathfrak p$, can be joined to a chain from
$\mathfrak p$ to $\mathfrak m$ of length $\dim R/\mathfrak p$.
The entire chain lies in the support because support is closed
under specialization. Thus the sum on the right is at most
$d_R(M)$. The depth localization bound of \ref{res:DP-10}
and the depth bound of \ref{res:DP-02} give
\[
 d_R(M)=\operatorname{depth}_R M
 \le \operatorname{depth}_{R_{\mathfrak p}} M_{\mathfrak p}
       +\dim R/\mathfrak p
 \le d_{R_{\mathfrak p}}(M_{\mathfrak p})+\dim R/\mathfrak p .
\]
All bounds are equalities. For the global statements, every
prime is contained in a maximal ideal, and repeated localization
identifies the corresponding local modules. Local rings of a
localization are local rings at primes disjoint from the
multiplicative set.
\end{proof}

\begin{corollary}[Heights and prime chains]\label{res:CM-07}
Let $R$ be Cohen--Macaulay and Noetherian. For every proper
ideal $I$, its grade on $R$ equals its height. If $R$ is local,
\[
 \operatorname{ht} I+\dim R/I=\dim R .
\]
Every saturated chain between two fixed primes
$\mathfrak p\subseteq\mathfrak q$ has length
$\operatorname{ht}\mathfrak q-\operatorname{ht}\mathfrak p$.
In particular $R$ is catenary, and all maximal chains in a
local Cohen--Macaulay ring have length $\dim R$.
\end{corollary}
\begin{proof}
Grade is the minimum of the localized depths at primes
containing $I$ by \ref{res:DP-06}. Their depths equal their heights by
the local Cohen--Macaulay condition. For $R$ local,
\ref{res:CM-06} gives
$\dim R/\mathfrak p=\dim R-\operatorname{ht}\mathfrak p$.
Compute $\dim R/I$ by taking the maximum over primes
containing $I$, or just the finitely many minimal primes there.
If $\mathfrak p\subsetneq\mathfrak q$ is saturated, the quotient
$R_{\mathfrak q}/\mathfrak pR_{\mathfrak q}$ has dimension one.
Apply the local dimension equality in $R_{\mathfrak q}$ to
$\mathfrak pR_{\mathfrak q}$. Its height there equals
$\operatorname{ht}\mathfrak p$, since localizing further gives
$R_{\mathfrak p}$. Therefore each saturated step increases
height by one. Summing over the chain proves the length assertion.
A maximal chain in a local ring starts at a minimal prime of
height zero and ends at the maximal ideal.
\end{proof}

\begin{proposition}[Regular local rings]\label{res:CM-08}
A regular local ring is Cohen--Macaulay. If its dimension is
$d$ and $x_1,\ldots,x_d$ minimally generate its maximal ideal,
then the quotient by the first $c$ elements is regular local
of dimension $d-c$, for $0\le c\le d$. In contrast, the
quotient by an arbitrary regular sequence is guaranteed
Cohen--Macaulay and can fail to be regular.
\end{proposition}
\begin{proof}
The basic structure theorem of \cite[ALG-05-C03-S03]{p103}
proves that the minimal generators form a regular sequence
and that a regular local ring of dimension zero is a field.
The first assertion follows by \ref{res:CM-04}.
The quotient by the first $c$ elements has dimension $d-c$
by \ref{res:CM-04}. The remaining $d-c$ generators generate
its maximal ideal. The height bound forces that number to
be minimal, so this quotient is regular local by definition.
For an arbitrary regular sequence use \ref{res:CM-05}
at each step. In a discrete valuation ring with parameter
$t$, the quotient by $t^2$ is zero-dimensional, has nonzero
maximal ideal, and is not regular.
\end{proof}

\begin{theorem}[Height, complete intersections, and unmixedness]
\label{res:CM-09}
Suppose a proper ideal $I$ of height $c$ in a Noetherian
Cohen--Macaulay ring is generated by $c$ elements. At every
maximal ideal containing $I$, the displayed generators
form a regular sequence. The quotient $R/I$ is
Cohen--Macaulay, and all its associated primes, viewed as
primes of $R$, are minimal over $I$ and have height $c$.
In particular a quotient of a regular local ring by
a regular sequence in its maximal ideal is Cohen--Macaulay.

Conversely, a Noetherian ring is Cohen--Macaulay if each
proper ideal generated by as many elements as its height
has all associated primes of that height.
\end{theorem}
\begin{proof}
The height theorem \cite[ALG-05-C03-R01]{p103} gives
$\operatorname{ht}\mathfrak p\le c$ for each prime minimal over
$I$. By the definition of height all have height $c$.
Localize at a maximal ideal containing $I$. Heights of its
subprimes are unchanged, and every prime of that localization
minimal over the localized ideal comes from a prime minimal
over $I$ contained in the maximal ideal. Thus the localized
ideal still has height $c$. By \ref{res:CM-07} its quotient
has dimension $\dim R_{\mathfrak m}-c$, so
\ref{res:CM-04} gives regularity and Cohen--Macaulayness of
the quotient. Then \ref{res:CM-06} gives the global result.
If an associated prime of $R/I$ were not minimal over $I$,
localization at that prime would yield an embedded associated
prime, contrary to \ref{res:CM-03}. Indeed an element with that
annihilator stays nonzero with maximal annihilator after
localizing. The assertion about regular local rings follows
from \ref{res:CM-08} and \ref{res:CM-05}.

For the converse fix a maximal ideal $\mathfrak m$.
Construct a regular sequence in $\mathfrak m$ on $R$ by
choosing, at each stage, an element of $\mathfrak m$
avoiding all associated primes of the current quotient.
This is possible by prime avoidance while $\mathfrak m$
is not associated, since it is not contained in any other
proper prime. The process terminates because the sequences
stay regular at $\mathfrak m$, where lengths are bounded
by $\dim R_{\mathfrak m}$. Write $x_1,\ldots,x_e$ for the
result and put $J=(x_1,\ldots,x_e)$. Then $\mathfrak m$
is associated to $R/J$.
For every prime over $J$ the sequence stays regular on the
localized ring, because localization preserves injectivity
and the entries remain in the maximal ideal. The bound of
\ref{res:DP-02} shows that every such prime has height at least
$e$, while the height theorem gives equality at each minimal
prime over $J$. The assumed unmixedness therefore gives
$\operatorname{ht}\mathfrak m=e$. The regular sequence
of this length at $\mathfrak m$ makes $R_{\mathfrak m}$
Cohen--Macaulay by \ref{res:DP-02}. As $\mathfrak m$ was arbitrary,
\ref{res:CM-06} proves the claim.
\end{proof}

\section{Flat local maps and Serre's conditions}

\begin{theorem}[Depth and dimension over a flat local family]
\label{res:CM-10}
Let $(R,\mathfrak m,k)\to(S,\mathfrak n)$ be a local map of
Noetherian local rings, $M\ne0$ a finite $R$-module, and $N\ne0$
a finite $S$-module flat over $R$. Put $N_0=N/\mathfrak mN$,
viewed over $S_0=S/\mathfrak mS$. Then
\begin{align*}
 d_S(M\otimes_R N)&=d_R(M)+d_{S_0}(N_0),\\
 \operatorname{depth}_S(M\otimes_R N)
   &=\operatorname{depth}_R M+\operatorname{depth}_{S_0}N_0 .
\end{align*}
Consequently $M\otimes_R N$ is Cohen--Macaulay if and only if
both $M$ and $N_0$ are. For a flat local map $R\to S$, the
ring $S$ is Cohen--Macaulay exactly when $R$ and $S/\mathfrak mS$ are.
Completion over a Noetherian local ring preserves the
dimension, depth, and Cohen--Macaulay property of finite modules.
\end{theorem}
\begin{proof}
The support of $M\otimes_R N$ consists of the primes
$\mathfrak q\in\operatorname{Supp}_S N$ whose contractions
$\mathfrak p$ belong to $\operatorname{Supp}_R M$.
Indeed, localize the tensor at $\mathfrak q$. If either
localized factor is zero, the tensor vanishes. If both are
nonzero, then $M_{\mathfrak p}/\mathfrak pM_{\mathfrak p}$
and $N_{\mathfrak q}/\mathfrak pN_{\mathfrak q}$ are nonzero
by Nakayama. Their tensor over $\kappa(\mathfrak p)$ is
nonzero and is the reduction of the localized tensor.
In particular $M\otimes_R N$ and $N_0$ are nonzero.

For the upper dimension bound choose a system of parameters
of $M$, then lift a system of parameters of $N_0$ to $S$.
Quotienting the tensor product by the first list replaces
$M$ by its finite-length quotient. The support formula
identifies the resulting support with that of $N_0$ in
the closed fiber. The second list leaves support consisting
of $\mathfrak n$, so apply \ref{res:CM-02}.

For the lower bound, choose a maximal-length chain in the
support of $N_0$, with bottom $\mathfrak q_0$ viewed in $\operatorname{Spec} S$.
If $\mathfrak q\in\operatorname{Supp}_S N$ contracts to
$\mathfrak p$, then $N_{\mathfrak q}$ is flat over
$R_{\mathfrak p}$. To see this, tensor over $R$ with $N$,
then localize, and use that tensor over $R_{\mathfrak p}$
agrees with tensor over $R$ on localized modules. Its closed
fiber is nonzero by local Nakayama, so it is faithfully
flat \cite[Theorem ALG-03-C02-S04]{p103}. For any
$\mathfrak p'\subseteq\mathfrak p$, the tensor of
$N_{\mathfrak q}$ with $\kappa(\mathfrak p')$ is nonzero.
A prime in its support over the corresponding fiber of
$S_{\mathfrak q}$ gives a prime below $\mathfrak q$ over
$\mathfrak p'$ where $N$ is nonzero. Hence a longest chain
in $\operatorname{Supp}_R M$ ending at $\mathfrak m$ can
be lifted downward from $\mathfrak q_0$. Appending the
fiber chain and using the support formula gives equality.

For depth, first suppose $y\in\mathfrak n$ is a
nonzerodivisor on $N_0$. Flatness identifies the graded
pieces of the $\mathfrak m$-adic filtration on $N$ with
$(\mathfrak m^i/\mathfrak m^{i+1})\otimes_k N_0$. Thus
$y$ is injective on the pieces. From the short exact sequence
between successive truncations it is injective on every
$N/\mathfrak m^iN$, starting with $i=1$.
Krull intersection over $S$ gives $\bigcap_i\mathfrak m^iN=0$
\cite[Theorem ALG-05-C04-R04]{p103}, hence $y$ is injective
on $N$. From $0\to N\xrightarrow{y}N\to N/yN\to0$
the long Tor sequence identifies $\operatorname{Tor}_1^R(k,N/yN)$
with the kernel of $y$ on $N_0$. It vanishes, so the local
criterion \cite[Theorem ALG-05-C04-S03]{p103} proves that
$N/yN$ is flat over $R$.

Induct on $\operatorname{depth}_R M+\operatorname{depth}_{S_0}N_0$.
When both depths are zero, the depth-zero criterion
\ref{res:DP-01} gives injections $k\to M$ and
$S/\mathfrak n\to N_0$. Flatness tensors the first to
$N_0\to M\otimes_R N$, giving depth zero for the tensor.
If $M$ has positive depth, choose an $M$-regular
$x\in\mathfrak m$. It remains regular on $M\otimes_R N$
by flatness, and the quotient is $(M/xM)\otimes_R N$.
Apply induction and \ref{res:DP-05}. If the fiber has
positive depth, lift a fiber-regular element to $y\in\mathfrak n$.
By the lifting argument, $N/yN$ is flat, so tensoring
$0\to N\xrightarrow{y}N\to N/yN\to0$ with $M$ is exact
also on the left by Tor vanishing. Apply induction to
$M\otimes_R(N/yN)$, lowering the fiber depth by one, and
use \ref{res:DP-05} for the change of depth on the tensor.
All successive quotients are nonzero by Nakayama.

Subtract the two formulas and use \ref{res:DP-02}:
both differences on the right are nonnegative. For the
ring case take $M=R$ and $N=S$. Finally use the local flat
map $R\to\widehat R$, with closed fiber $k$, and
$\widehat M\simeq M\otimes_R\widehat R$
\cite[Theorems ALG-05-C04-R03 and ALG-05-C04-S01]{p103}.
The zero module is preserved by completion as well.
\end{proof}

\begin{definition}[Serre's depth conditions]\label{res:CM-11}
A finite module $M$ over a Noetherian ring satisfies $(S_k)$
if for all $\mathfrak p\in\operatorname{Supp}M$,
\[
 \operatorname{depth}_{R_{\mathfrak p}}M_{\mathfrak p}
  \ge \min\{k,d_{R_{\mathfrak p}}(M_{\mathfrak p})\}.
\]
For a ring use $M=R$. A Cohen--Macaulay module satisfies
every $(S_k)$, and satisfying all $(S_k)$ implies
Cohen--Macaulayness. The condition $(S_1)$ holds exactly when
$M$ has no embedded associated primes.

For a flat map $R\to S$ of Noetherian rings and a prime
$\mathfrak q$ over $\mathfrak p$, if $S_{\mathfrak q}$
satisfies $(S_k)$ at all its primes, so does $R_{\mathfrak p}$.
If the base ring and every fiber satisfy $(S_k)$, then so
does $S$.
\end{definition}
\begin{proof}
The characterization by all $k$ follows from the depth
bound \ref{res:DP-02}, because each local dimension is finite.
A prime in the support has localized depth zero exactly if
it belongs to $\operatorname{Ass}_R M$. In one direction an
element with that annihilator stays nonzero with maximal
annihilator after localizing. Conversely, multiply a
representative of a nonzero element killed by
$\mathfrak pR_{\mathfrak p}$ by a common denominator
outside $\mathfrak p$, so that all the finitely many
generators of $\mathfrak p$ kill it in $M$. Its localization
remains nonzero, so no element outside $\mathfrak p$ kills
it. This gives annihilator $\mathfrak p$. Since the local
support dimension is zero exactly at its minimal primes,
the assertion about $(S_1)$ follows.

For descent localize at $\mathfrak p,\mathfrak q$. For every
$\mathfrak p'$ of the local base $R_{\mathfrak p}$ there
is a prime of $S_{\mathfrak q}$ over it by flat going down
\cite[Theorem ALG-05-C01-R04]{p103}. Replace the prime by
a minimal prime of the fiber below it. Its local fiber has
depth and dimension zero, so \ref{res:CM-10} transports the
$(S_k)$ inequality to $\mathfrak p'$. For ascent write
$a,b$ for the local base and fiber dimensions at the chosen
prime. The sum of the two depth bounds is at least
$\min(k,a)+\min(k,b)\ge\min(k,a+b)$, giving the result again
by \ref{res:CM-10}.
\end{proof}

\section{Syzygies, flatness, and examples}

\begin{proposition}[Syzygies and finite projective dimension]
\label{res:CM-12}
Let $R$ be a Cohen--Macaulay Noetherian local ring of
dimension $d$ and $M\ne0$ a finite $R$-module of depth $e$.
A $(d-e)$th syzygy obtained from successive finite free
surjections is maximal Cohen--Macaulay or zero. If $M$ has
finite projective dimension, then $M$ is Cohen--Macaulay
if and only if
\[
 \operatorname{pd}_R M=\operatorname{grade}_R(\operatorname{Ann}_R M,R).
\]
Over a regular local ring the nonzero maximal
Cohen--Macaulay modules are precisely the nonzero finite
free modules. In particular, all finite modules over a
regular local ring of dimension $d$ have projective
dimension at most $d$, and the residue field has projective
dimension exactly $d$.
\end{proposition}
\begin{proof}
In $0\to K\to F\to M\to0$ with $F$ nonzero finite free,
the depth of $F$ is $d$: regularity of a sequence is tested
componentwise. If $K\ne0$, \ref{res:DP-07} gives
$\operatorname{depth}K\ge\min(d,\operatorname{depth}M+1)$.
Iterate and use \ref{res:DP-02}. If a syzygy is zero, its
subsequent free modules and syzygies can be taken to be zero.

For finite projective dimension use the Auslander--Buchsbaum
formula \ref{res:DP-13} to express $\operatorname{pd}_R M$
as $d-\operatorname{depth}M$. By \ref{res:CM-07},
the grade of the annihilator is $d-d_R(M)$. Their equality
is exactly the Cohen--Macaulay condition.

Now take $R$ regular local and choose minimal generators
$x_1,\ldots,x_d$ of its maximal ideal. They are regular
on $R$ by the structure theorem in \ref{res:CM-08}.
For a nonzero maximal Cohen--Macaulay module $M$ one has
$d_R(M)=d$ by \ref{res:DP-02}. The list of generators is a
system of parameters on $M$ and hence is $M$-regular by
\ref{res:CM-04}. By \ref{res:RK-07} the Koszul resolution
of $k$ computes $\operatorname{Tor}_1^R(k,M)=0$. Choose a
free surjection $R^r\to M$ lifting a basis of $M/\mathfrak mM$.
Tensor the sequence of its kernel $K$ with $k$. The vanishing
just proved identifies $K/\mathfrak mK$ with the kernel of
$k^r\to M/\mathfrak mM$. The latter is zero, and Nakayama
gives $K=0$. Conversely regular sequences on $R$ act
regularly on every nonzero finite free module.
For a nonzero finite module, the earlier syzygy bound
now gives a free resolution of length at most $d-e$.
The residue field has depth zero, and the
Auslander--Buchsbaum formula \ref{res:DP-13} gives its
projective dimension $d$.
\end{proof}

\begin{theorem}[Miracle flatness with a dimension equality]
\label{res:CM-13}
Let $(R,\mathfrak m)\to(S,\mathfrak n)$ be a local map of
Noetherian local rings with $R$ regular local. Let $N\ne0$
be a finite Cohen--Macaulay $S$-module. If
\[
 d_S(N)=\dim R+d_{S/\mathfrak mS}(N/\mathfrak mN),
\]
then $N$ is flat over $R$. In particular a Cohen--Macaulay
$S$ is flat over a regular local $R$ whenever
$\dim S=\dim R+\dim(S/\mathfrak mS)$.
\end{theorem}
\begin{proof}
Choose a minimal sequence of generators $x_1,\ldots,x_d$
of $\mathfrak m$, regular on $R$ by \ref{res:CM-08}.
Lift a system of parameters of the nonzero finite module
$N/\mathfrak mN$ to elements of $\mathfrak n$. Concatenating
it with the images of the $x_i$ first gives a sequence with
finite-length quotient. The dimension equality says that
it is a system of parameters for $N$. Therefore its
first $d$ elements are $N$-regular by \ref{res:CM-04}.
The Koszul resolution of $R/\mathfrak m$ in
\ref{res:RK-07} remains exact in positive degrees after
tensoring with $N$. Thus $\operatorname{Tor}_1^R(R/\mathfrak m,N)=0$,
and \cite[Theorem ALG-05-C04-S03]{p103} proves flatness.
The empty parameter list is allowed if the base is a field.
\end{proof}

\begin{example}[One-dimensional and embedded cases]\label{res:CM-14}
The local ring at $(x,y)$ of $k[x,y]/(y^2-x^3)$ is
Cohen--Macaulay of dimension one and is not regular.
The local polynomial ring has dimension two by the
coordinate chain and \cite[Lemma polynomial-dimension]{p103},
and its maximal ideal has two generators. It is regular
local. The nonzero relation is regular there, so use
\ref{res:CM-02}, \ref{res:CM-05}, and \ref{res:CM-08}.
As the relation is in the square of the maximal ideal,
the coordinate classes still give a two-dimensional
cotangent space, proving nonregularity.

The localization of $k[x,y]/(x^2,xy)$ at $(x,y)$ has
dimension one and depth zero. Its primes correspond to
those of $k[y]$, because the class of $x$ is nilpotent.
The class of $x$ has annihilator exactly $(x,y)$, by the
two defining monomials. It stays nonzero after localization.
This exhibits an embedded prime above $(x)$ and a failure
of $(S_1)$.

Every one-dimensional reduced Noetherian ring is
Cohen--Macaulay. In a one-dimensional local ring of this
kind, the direct product of the minimal-prime domain
quotients receives an injective map from the ring. An
element avoiding the minimal primes is regular by that
injection, and prime avoidance finds one in the maximal
ideal. This gives depth one. The zero-dimensional local
cases are fields.
\end{example}

\section{Permanence and connectedness}

\begin{proposition}[Polynomial rings, formal series, and field extension]
\label{res:CM-15}
Let $R$ be Noetherian and $M$ a finite $R$-module.
Its scalar extension to a polynomial or formal power
series ring over $R$ in finitely many variables is
Cohen--Macaulay if and only if $M$ is. In particular a
Noetherian Cohen--Macaulay ring is universally catenary.

If $R$ is a Noetherian $k$-algebra over a field and
$K/k$ is a field extension, assume $R$ is of finite
type over $k$ or $K$ is finitely generated as a field
over $k$. Then $R\otimes_k K$ is Noetherian, and it is
Cohen--Macaulay if and only if $R$ is.
\end{proposition}
\begin{proof}
For one polynomial variable apply \ref{res:CM-10} at
each prime of $R[T]$ and its contraction in $R$. The
local closed fiber is a localization of $\kappa(\mathfrak p)[T]$.
It is either a field or a one-dimensional local principal
ideal domain. In the latter case a generator of the
maximal ideal is regular, so the fiber is Cohen--Macaulay.
When the localized base module is zero the scalar extension
is zero, and in the other cases the formula applies.
For descent take the prime $(\mathfrak p,T)$ over each
$\mathfrak p$. Iterate for finitely many variables.

For one formal variable, $R[[T]]$ is the $(T)$-adic
completion of $R[T]$: the compatible truncations are the
coefficients of a unique series. It is Noetherian and
flat over $R$ by the Hilbert basis and completion results
\cite[Theorems ALG-05-C02-R01 and ALG-05-C04-S01]{p103}.
Every maximal ideal contains $T$, since $1-Tu$ has an
inverse given by its geometric series for any $u$.
The maximal ideals thus correspond to those of $R$. If
$\mathfrak p$ is maximal, the fiber over it is
$(R/\mathfrak p)[[T]]$: using a fixed finite set of
generators of $\mathfrak p$, each series with coefficients
in $\mathfrak p$ is a sum of those generators times series.
Over the residue field, every nonzero series is a power
of $T$ times a unit, giving dimension and depth one.
Now apply \ref{res:CM-10} at the maximal ideals and use
\ref{res:CM-06} to pass from those to all primes.
Iteration gives the assertion for several variables.
Finite polynomial extensions of $R$ are Cohen--Macaulay,
hence catenary by \ref{res:CM-07}. Quotienting a catenary
ring preserves catenarity because intervals between primes
in the quotient are the corresponding intervals upstream.
All finite-type $R$-algebras are such quotients, proving
universal catenarity.

If $K/k$ is finitely generated as a field, take a transcendence
basis and then adjoin finitely many algebraic generators one
at a time. Tensoring $R$ first with the rational function field
gives a localization of a polynomial extension of $R$.
At each algebraic step the new ring is a polynomial extension
of the preceding ring modulo the monic minimal polynomial
of the added element. Multiplication by a monic polynomial
is injective, by the leading coefficient of a nonzero
polynomial. Thus the new ring is Noetherian and retains
Cohen--Macaulayness by the polynomial case and \ref{res:CM-05}
at each prime of the quotient. Descent follows from the
faithfully flat map $R\to R\otimes_k K$ and the local
Cohen--Macaulay criterion in \ref{res:CM-10}.

If instead $R$ is of finite type over $k$, its scalar extension
is of finite type over $K$ and is Noetherian. For each $\mathfrak p$,
the field $\kappa(\mathfrak p)$ is finitely generated over $k$.
The fiber $\kappa(\mathfrak p)\otimes_k K$ is Cohen--Macaulay
by the case just proved, with $K$ as the Noetherian algebra
and $\kappa(\mathfrak p)$ as the finitely generated field
extension. Apply \ref{res:CM-10} to the localized map at
each prime, using faithfulness for descent again.
\end{proof}

\begin{proposition}[Connectedness in codimension one]\label{res:CM-16}
Let $R$ be Noetherian local and let $I,J$ be proper ideals
with incomparable radicals. Suppose $I\cap J$ is nilpotent.
Then $\operatorname{grade}_R(I+J,R)\le1$. If $R$ satisfies
$(S_2)$, in particular if $R$ is Cohen--Macaulay, then
$\operatorname{ht}(I+J)\le1$. Thus a partition of the
minimal components into two disjoint nonempty families
cannot have their unions meet only in codimension at
least two.
\end{proposition}
\begin{proof}
Replacing the ideals by large powers leaves their radicals
unchanged and makes their product zero: $(IJ)^n$ is zero
for all large $n$ because $IJ\subseteq I\cap J$. The radical
of their sum is also unchanged, so grade and height of the
sum are unchanged by \ref{res:DP-06} and the prime description
of height. If the new intersection is nonzero, its nonzero
elements are killed by the sum since $IJ=0$. This gives
grade zero. Otherwise use the exact sequence
\[
 0\longrightarrow R\longrightarrow R/I\oplus R/J
   \longrightarrow R/(I+J)\longrightarrow0 ,
\]
where the first map sends $r$ to $(r,-r)$ modulo the ideals,
and the second adds the two coordinates. Exactness in the
middle follows because if $a+b=i+j$ with $i\in I,j\in J$,
then $r=a-i$ satisfies $(a,b)\equiv(r,-r)$.
If the grade of the sum were at least two,
$\operatorname{Ext}^1_R(R/(I+J),R)=0$ by \ref{res:DP-04}.
Then the extension would split
\cite[Theorem ALG-10-C03-S02]{p103}.
A retraction would express $1$ as a difference of an element
of $\operatorname{Ann} I$ and one of $\operatorname{Ann} J$.
Both $I$ and $J$ are nonzero because their radicals are
incomparable. Thus both annihilators are contained in the
maximal ideal, an impossibility.

If every prime containing $I+J$ had height at least two,
$(S_2)$ and the grade localization formula \ref{res:DP-06}
would give grade at least two. This proves the height
bound. For the final assertion take $I,J$ to be the
intersections of the minimal primes in the two families.
Their combined intersection is the nilradical, which is
nilpotent in a Noetherian ring. Their radicals are
incomparable: for any minimal prime in one family choose
an element from each prime in the other family outside
that first prime and multiply them. The product is in
the other family's intersection but outside the first prime.
The sum defines the intersection of the two unions,
as required.
\end{proof}

\begin{corollary}[Smooth algebras over a Cohen--Macaulay base]
\label{res:CM-17}
Suppose $A$ is Noetherian and Cohen--Macaulay and $B$ is
smooth over $A$. Then $B$ is Cohen--Macaulay.
\end{corollary}
\begin{proof}
By \ref{res:SM-08}, the localized base-to-target maps are
flat. It remains by \ref{res:CM-10} to check all local
rings of each field fiber. Smoothness is preserved by
base change and can be tested on standard smooth charts
by \ref{res:SM-05} and \ref{res:SM-06}. For a standard chart
over a field $F$, extend scalars to an algebraic closure.
The local rings at the closed points of that chart become
regular local by \ref{res:SM-12}, hence Cohen--Macaulay
by \ref{res:CM-08}. Every prime is under a maximal ideal,
so their localizations are Cohen--Macaulay by \ref{res:CM-06}.
The scalar extension of the chart is faithfully flat.
For each prime of the original chart, choose a prime over
it in the extended chart and use \ref{res:CM-10} to descend
Cohen--Macaulayness. All fiber rings at points follow by
restricting to these charts.
\end{proof}

\begin{proposition}[Type under flat local scalar extension]\label{res:CM-18}
Under the hypotheses and notation of \ref{res:CM-10}, the
types of the three nonzero modules satisfy
\[
 \operatorname{type}_S(M\otimes_R N)
  =\operatorname{type}_R M\cdot\operatorname{type}_{S_0} N_0 .
\]
\end{proposition}
\begin{proof}
Choose a maximal regular sequence on $M$ in $\mathfrak m$.
Flatness carries it to a regular sequence on the tensor.
Quotient $M$ by the sequence. Next choose a maximal regular
sequence on $N_0$ and lift it to $\mathfrak n$. The lifting
argument in \ref{res:CM-10}, applied successively, keeps the
quotients of $N$ flat over $R$ and the multiplication maps
injective on the tensor. Thus we have replaced $M,N$ by
quotients $M',N'$ for which both $M'$ and $N'/\mathfrak m N'$
have depth zero. By \ref{res:CM-10} the combined sequence is
maximal on the tensor, so in all three cases it remains to
compute the socles of the resulting quotients by
\ref{res:DP-05}.

For finite generators $a_1,\ldots,a_r$ of $\mathfrak m$,
the socle of $M'$ is the kernel of
$M'\to (M')^r$, $u\mapsto(a_1u,\ldots,a_ru)$.
Tensor this kernel sequence with the flat module $N'$.
The submodule of $M'\otimes_R N'$ killed by $\mathfrak m$
becomes
\[
 \operatorname{Soc}_R(M')\otimes_k (N'/\mathfrak mN').
\]
The socle of the full $S$-module is obtained by taking
the elements of this submodule killed by $\mathfrak n$.
Choose a $k$-basis of $\operatorname{Soc}_R(M')$.
By inspecting coordinates in that basis, the resulting
socle is
$\operatorname{Soc}_R(M')\otimes_k
 \operatorname{Soc}_{S_0}(N'/\mathfrak mN')$.
The dimension over $S/\mathfrak n$ is the asserted product.
\end{proof}

\begin{theorem}[The residue-field test for regularity]
\label{res:CM-19}
A Noetherian local ring $R$ is regular local if and only
if its residue field $k$ has finite projective dimension.
If $R$ is regular local of dimension $d$, its global
dimension is $d$, and its localizations at primes are
regular local.
\end{theorem}
\begin{proof}
The forward direction is \ref{res:CM-12}. For the converse
let $n=\operatorname{pd}_R k<\infty$ and induct on $n$.
If $n=0$, $k$ is a finite free module over the local ring.
Since $k$ is nonzero and annihilated by $\mathfrak m$, this
forces $\mathfrak m=0$. Suppose $n>0$.
Then $\operatorname{depth}_R R=n$ by \ref{res:DP-13},
so $\mathfrak m$ is not associated. The finite prime
avoidance lemma \cite[ALG-05-C02-S01]{p103}, applied
to the associated primes together with $\mathfrak m^2$,
gives $x\in\mathfrak m\setminus\mathfrak m^2$ avoiding
all associated primes. Here $\mathfrak m\ne\mathfrak m^2$
by Nakayama. Thus $x$ is regular and extends to a minimal
set of generators of $\mathfrak m$.

Put $B=R/(x)$. Take a minimal finite free resolution
$F_\bullet\to k$ as in \ref{res:DP-12}, with $F_0=R$
and with the first vector $e_1$ of a basis of $F_1$
mapping to $x$. This is possible by the construction in
\ref{res:DP-12} from the minimal generators of $\mathfrak m$.
Let $C_\bullet=F_\bullet\otimes_R B$. The short exact
sequence of complexes given by multiplication by $x$
on $F_\bullet$ shows that $H_i(C)=0$ for $i\ge2$ and
that its connecting map $H_1(C)\to k$ is an isomorphism.
The class of $e_1$ maps to $1$, because its differential
in $F$ is $x\cdot1\in F_0$.

Put $V=C_1/\operatorname{im}(C_2\to C_1)$. It has finite
projective dimension over $B$ since the exact free tail of
$C_\bullet$ resolves it. Taking the first coefficient and
reducing it modulo the maximal ideal of $B$ defines
a map $V\to k$: the image of $C_2\to C_1$ has entries
in that maximal ideal by minimality. It sends the class
of $e_1$ to $1$, so retracts the inclusion $H_1(C)\simeq
k\subseteq V$. A direct summand of a module of finite
projective dimension has finite projective dimension:
its higher Ext groups are summands of those of the module,
and dimension shifting followed by the projectivity test
from \ref{res:DP-12} converts their vanishing to a finite
projective resolution. Thus $k$ has finite projective
dimension over $B$. By \ref{res:DP-05} and \ref{res:DP-13}
this dimension is $n-1$. Induction makes $B$ regular
local. The image of $x$ in $\mathfrak m/\mathfrak m^2$
spans a one-dimensional subspace, so the embedding
dimension of $R$ is that of $B$ plus one. By \ref{res:CM-02}
the same equality holds for the Krull dimensions.
This proves regularity of $R$.

Now assume $R$ is regular. For any prime $\mathfrak p$,
a finite projective resolution of $R/\mathfrak p$ localizes
to one for the residue field of $R_{\mathfrak p}$. The
just-proved converse makes $R_{\mathfrak p}$ regular.
For the global dimension statement let $T$ be any
$R$-module, not necessarily finite. By \ref{res:CM-12},
every quotient $R/I$ has projective dimension at most
$d$. Apply dimension shifting to an injective resolution
of $T$. If $W$ is its $d$th cosyzygy, then
$\operatorname{Ext}^1_R(R/I,W)=0$ for every ideal $I$.
This is Baer's criterion for injectivity, since the
Hom sequence for $0\to I\to R\to R/I\to0$ extends each
map $I\to W$ to $R$ \cite[Chapter 2]{p103}. Hence every
module has injective dimension at most $d$. All Ext
groups in degrees greater than $d$ vanish, and dimension
shifting with the projectivity test in \ref{res:DP-12}
gives projective dimension at most $d$ for every module.
Equality is attained by $k$ by \ref{res:CM-12}.
\end{proof}

\paragraph{Further topics.}
The complete-intersection consequence of \ref{res:CM-09}
uses an explicit regular local presentation. Intrinsic
characterizations through completions and higher Tor, the
Cohen--Macaulayness of determinantal rings via determinantal
resolutions, and the structure of canonical modules require
their own theorems. See \cite[Section 2.3]{brunsherzog} for the
complete-intersection characterizations and
\cite[Theorem 18.18]{eisenbud} for the determinantal case.
\chapter{Serre's Criterion for Normality}\label{ch:normality}

Throughout this chapter the rings are Noetherian. The depth of a
local ring means its depth as a module over itself. We use the
associated-prime and localization results of P1-03
\cite[Chapter 5]{p103}. Every ring-theoretic Serre condition is
tested at the primes of the ring.

\section{Associated primes and total quotients}

\begin{definition}[Serre conditions and normality]\label{res:SN-01}
For an integer $j\ge 0$, $(R_j)$ means that $A_{\mathfrak p}$
is regular local whenever $\operatorname{ht}\mathfrak p\le j$.
For a finite $A$-module $M$, $(S_j)$ means
\[
 \operatorname{depth}_{A_{\mathfrak p}}M_{\mathfrak p}
 \ge \min\{j,\dim\operatorname{Supp}_{A_{\mathfrak p}}M_{\mathfrak p}\}
 \qquad(\mathfrak p\in\operatorname{Supp}M).
\]
For the ring, use $M=A$. The ring is normal if every prime
localization is an integrally closed domain. Write $Q(A)$ for the
localization of $A$ at its nonzerodivisors.
\end{definition}

\begin{lemma}[Detection by associated primes]\label{res:SN-02}
If $M$ is finite and $0\ne x\in M$, then some
$\mathfrak p\in\operatorname{Ass}_A(M)$ has $x/1\ne0$ in
$M_{\mathfrak p}$. Moreover $M$ satisfies $(S_1)$ exactly when
its associated primes are all minimal in its support.
\end{lemma}
\begin{proof}
The nonzero cyclic submodule $Ax$ has an associated prime
$\mathfrak p$, which is also associated to $M$. Its annihilator
contains $\operatorname{Ann}(x)$, so no element outside
$\mathfrak p$ annihilates $x$. For the second assertion, the
maximal ideal of $A_{\mathfrak p}$ is associated to
$M_{\mathfrak p}$ exactly when $\mathfrak p$ is associated
to $M$, by localization of associated primes. This is the
depth-zero condition of \ref{res:DP-01}. The local support has
dimension zero precisely when $\mathfrak p$ is minimal in
the support of $M$, since its primes correspond to support
primes below $\mathfrak p$. Combine the two statements.
\end{proof}

\begin{proposition}[The reducedness criterion]\label{res:SN-03}
The ring $A$ is reduced if and only if it satisfies $(R_0)$
and $(S_1)$.
\end{proposition}
\begin{proof}
For a reduced ring, each localization at a minimal prime is
a reduced zero-dimensional local ring. Its maximal ideal is
its nilradical, hence is zero, so the localization is a field.
Write $\mathfrak p_1,\ldots,\mathfrak p_r$ for the finitely many
minimal primes. Their intersection is zero. If $a$ avoids all
these primes and $ab=0$, primeness gives $b\in\mathfrak p_i$
for every $i$ and hence $b=0$. A nonminimal prime is contained
in none of the minimal primes, so prime avoidance supplies
in that prime an element avoiding all of them. After localizing
it gives a nonzerodivisor in the local maximal ideal. This
proves $(S_1)$.

Conversely, $(S_1)$ and \ref{res:SN-02} put all associated
primes among the minimal ones. A nilpotent element vanishes in
every such localization, which is a field by $(R_0)$. By the
detection statement it vanishes in $A$.
\end{proof}

\begin{proposition}[Total quotients of a reduced ring]\label{res:SN-04}
For reduced $A$, the total quotient ring is the product
\[
 Q(A)\ \simeq\ \prod_{\mathfrak p\ {\rm minimal}}\operatorname{Frac}(A/\mathfrak p).
\]
Moreover $A$ is normal if and only if it is integrally closed
in $Q(A)$, and then $A$ is a finite product of integrally closed
domains.
\end{proposition}
\begin{proof}
The nonzerodivisors are exactly the elements avoiding every
minimal prime. One inclusion was proved in \ref{res:SN-03}.
For the other, if $a\in\mathfrak p_i$, choose for each $j\ne i$
$b_j\in\mathfrak p_j\setminus\mathfrak p_i$. Their product
$b$ avoids $\mathfrak p_i$ and $ab$ belongs to every minimal
prime, so $b\ne0$ and $ab=0$. With a single minimal prime the
empty product is $1$.
Any prime disjoint from the nonzerodivisors is contained in
the union of the minimal primes and hence in one of them
by prime avoidance. It is therefore minimal. Thus $Q(A)$
is reduced, Noetherian, and zero-dimensional, so its
Artinian decomposition gives a product of fields
\cite[Chapter 3, ALG-03-C03-S02]{p103}. The factors are
$A_{\mathfrak p}$ for the minimal primes, hence the indicated
fraction fields.

Suppose $A$ is normal and $z\in Q(A)$ is integral over $A$.
For each maximal ideal $\mathfrak m$ there is an injection
$Q(A)_{\mathfrak m}\hookrightarrow Q(A_{\mathfrak m})$:
the former localizes $A_{\mathfrak m}$ at images of
nonzerodivisors, and those remain nonzerodivisors. Integrality
puts the image of $z$ in the integrally closed domain
$A_{\mathfrak m}$. Hence the class of $z$ in $Q(A)/A$ vanishes
locally at every maximal ideal, and it is zero. In fact a
nonzero element has proper annihilator, contained in a maximal
ideal where the element survives.

Conversely, assume $A$ is integrally closed in $Q(A)$. The
coordinate idempotents of $Q(A)$ are integral, so belong to
$A$ and decompose it into subrings $A_i$ of the corresponding
fields. Each $A_i$ has that field as fraction field: a
fraction of elements of $A$ projects to a fraction in $A_i$,
and its denominator projects to a unit in the field.
An element of that field integral over $A_i$ gives an integral
element of $Q(A)$ by placing it in that coordinate and zero
elsewhere, using its monic equation in the chosen coordinate
and the corresponding monomial elsewhere. Thus $A_i$ is
integrally closed. Localizations of an integrally closed domain
are integrally closed. Indeed, from a monic equation with
coefficients in $S^{-1}A_i$, choose $s\in S$ so that
multiplication of the coefficient of degree $n-j$ by $s^j$
clears its denominator. Then $sz$ is integral over $A_i$,
and $z\in S^{-1}A_i$. Prime localizations of $A$ are therefore
integrally closed domains. For the zero ring the product is
empty and all local assertions are vacuous.
\end{proof}

\section{Codimension one and the depth bound}

\begin{lemma}[The determinant test for fractional multiplication]\label{res:SN-05}
Let $A$ be a domain and let $L$ be a nonzero finitely generated
$A$-submodule of its fraction field. If $uL\subseteq L$ for
$u$ in the fraction field, then $u$ is integral over $A$.
\end{lemma}
\begin{proof}
Take generators $v_i$ of $L$ and write $uv_i=\sum_j a_{ij}v_j$
with $a_{ij}\in A$. The adjugate of $(u\delta_{ij}-a_{ij})$
shows that its determinant annihilates each $v_i$. One generator
is nonzero, so the monic determinant polynomial vanishes at $u$.
\end{proof}

\begin{theorem}[One-dimensional normal local domains]\label{res:SN-06}
For a one-dimensional Noetherian local domain $(A,\mathfrak m)$,
the following are equivalent: $A$ is integrally closed, $A$ is
regular local, and $A$ is a discrete valuation ring. In this
case, if $\pi$ generates $\mathfrak m$, every nonzero fractional
element is a unit times a unique integral power of $\pi$.
\end{theorem}
\begin{proof}
Assume $A$ is integrally closed. Choose $0\ne a\in\mathfrak m$.
The radical of $(a)$ is $\mathfrak m$, since any prime containing
$a$ is nonzero and must be maximal in a one-dimensional local
domain. Finite generation of $\mathfrak m$ implies
$\mathfrak m^n\subseteq(a)$ for some $n$: take powers of each
generator lying in $(a)$ and use the pigeonhole bound on
monomials. For the least such positive $n$, choose
$b\in \mathfrak m^{n-1}\setminus(a)$. Then $u=b/a$ is outside
$A$ and $u\mathfrak m\subseteq A$. If
$u\mathfrak m\subseteq\mathfrak m$, \ref{res:SN-05} gives
integrality, a contradiction. There is therefore
$\pi\in\mathfrak m$ with $u\pi$ a unit. For every
$x\in\mathfrak m$,
\[
 x=\pi\,(ux)/(u\pi)\ \in\ \pi A.
\]
Thus $\mathfrak m=(\pi)$, giving regularity.

If $A$ is regular local of dimension one, its maximal ideal
is principal, by the definition of regularity and Nakayama.
The Krull intersection theorem \cite[ALG-05-C04-R04]{p103}
says $\bigcap_{n\ge0}(\pi^n)=0$. Every nonzero element of $A$
is therefore a unit times a unique \emph{nonnegative} power of
$\pi$, by choosing its largest such power. Dividing two such
elements gives the integral exponent on the fraction field.
The exponent is additive under multiplication. Every nonzero
ideal of $A$ is generated by an element of least exponent, so
this is a discrete valuation ring.

Finally a fractional element with negative exponent cannot
solve a monic equation over $A$. Dividing that equation by
its leading power would express $1$ as an element of
$\mathfrak m$, since each negative power of the fractional
element lies in $\mathfrak m$. Every integral element thus has
nonnegative exponent, so is in $A$.
\end{proof}

\begin{proposition}[Normality gives depth at least two in dimension at least two]\label{res:SN-07}
If $A$ is a Noetherian integrally closed local domain, then
$\operatorname{depth}A\ge\min\{2,\dim A\}$.
\end{proposition}
\begin{proof}
The assertion in dimension zero is automatic, and a nonzero
element of the maximal ideal gives the depth-one lower bound
in positive dimension. Assume $\dim A\ge2$ and choose such an
element $a$. If the depth were one, \ref{res:DP-05} would
give depth zero for $A/aA$. By \ref{res:DP-01} choose
$b\notin aA$ whose class is killed by $\mathfrak m$.
Then $u=b/a\notin A$ and $u\mathfrak m\subseteq A$.
If $u\mathfrak m\subseteq\mathfrak m$ we contradict integral
closedness using \ref{res:SN-05}. Otherwise choose
$\pi\in\mathfrak m$ with $u\pi$ a unit and argue as in
\ref{res:SN-06} to get $\mathfrak m=(\pi)$. The principal
ideal theorem \cite[Chapter 6]{p103} gives
$\operatorname{ht}\mathfrak m\le1$, again a contradiction.
\end{proof}

\begin{theorem}[Codimension-one membership under $(S_2)$]\label{res:SN-08}
Suppose $A$ satisfies $(S_2)$. For $z\in Q(A)$, if its image
in $Q(A_{\mathfrak p})$ lies in $A_{\mathfrak p}$ at every
prime $\mathfrak p$ of height at most one, then $z\in A$.
For a positive-dimensional Noetherian normal domain this
gives an intersection in its fraction field,
\[
 A=\bigcap_{\operatorname{ht}\mathfrak p=1} A_{\mathfrak p}.
\]
\end{theorem}
\begin{proof}
Write $z=a/b$ with $b$ a nonzerodivisor. If $a\notin bA$,
\ref{res:SN-02} gives a prime
$\mathfrak p\in\operatorname{Ass}(A/bA)$ at which the class
of $a$ is nonzero. At this prime $b$ is a nonzerodivisor
in the maximal ideal and the quotient has depth zero, so
\ref{res:DP-05} gives $\operatorname{depth}A_{\mathfrak p}=1$.
Condition $(S_2)$ forces $\operatorname{ht}\mathfrak p\le1$.
But $a/b$ in the image of $A_{\mathfrak p}$ means
$a\in bA_{\mathfrak p}$, a contradiction. For a normal
domain, \ref{res:SN-07} applied to its prime localizations
gives $(S_2)$, and the one height-zero prime imposes no
additional restriction on the fraction.
\end{proof}

\begin{theorem}[Serre's criterion]\label{res:SN-09}
A Noetherian ring is normal if and only if it satisfies
$(R_1)$ and $(S_2)$.
\end{theorem}
\begin{proof}
If $A$ is normal, its localizations of height zero are fields.
In height one apply \ref{res:SN-06}, and at all primes apply
\ref{res:SN-07}. Conversely $(R_1)$ and $(S_2)$ imply
reducedness by \ref{res:SN-03}. If $z\in Q(A)$ is integral
over $A$, its image is integral in every localized total
quotient. At primes of height at most one the local ring
is a field or a regular local domain of dimension one
\cite[ALG-05-C03-S03]{p103}, so is integrally closed by
\ref{res:SN-06}. Use \ref{res:SN-08} to deduce $z\in A$ and
\ref{res:SN-04} to conclude normality.
\end{proof}

\begin{corollary}[Regular and Cohen--Macaulay rings]\label{res:SN-10}
A Noetherian ring regular at every prime is normal. A
Cohen--Macaulay Noetherian ring is normal exactly when it is
regular at the primes of height at most one.
\end{corollary}
\begin{proof}
At each prime a regular local ring is Cohen--Macaulay by
\ref{res:CM-08}. Both assertions now follow from
\ref{res:SN-09}, since the depth equals the local dimension
in a Cohen--Macaulay ring.
\end{proof}

\section{Codimension-two extension and change of rings}

\begin{corollary}[Extension of functions across codimension two]\label{res:SN-11}
Let $A$ be a Noetherian normal domain and let
$U\subseteq\operatorname{Spec}A$ be an open set containing all
primes of height at most one. Every regular function on $U$
is the restriction of a unique element of $A$.
\end{corollary}
\begin{proof}
The generic point lies in $U$. The value there is an element
$z$ of the fraction field. At every point of $U$, locally on
a principal affine open, a regular function is a fraction
with denominator invertible on that open. Its generic value
determines it there because the coordinate ring of the open
embeds in the fraction field. Thus $z$ belongs to every local
ring in $U$. By \ref{res:SN-08} it is in $A$, and the same
generic comparison proves equality and uniqueness.
\end{proof}

\begin{example}[The two Serre conditions are independent]\label{res:SN-12}
The local cusp ring $k[t^2,t^3]_{(t^2,t^3)}$ satisfies $(S_2)$
and fails $(R_1)$. Indeed $k[t]$ is integral and finite over
$k[t^2,t^3]$ and both have dimension one by integral dimension
invariance \cite[Chapter 6]{p103}; the displayed maximal prime
is nonzero. The local ring is a one-dimensional domain, so has
depth one. The element $t$ is integral and belongs to its
fraction field since $t=t^3/t^2$. It is outside the local ring:
multiplying by a polynomial in $t^2,t^3$ with nonzero constant
term leaves a nonzero linear coefficient. By \ref{res:SN-06}
the ring fails $(R_1)$.

For the other direction let $B=k[x,y]_{(x,y)}$,
$C=k[z,w]_{(z,w)}$, and $A=B\times_k C$ with both maps given
by residues. This is the localization at the origin of
$k[x,y,z,w]/(xz,xw,yz,yw)$: denominator pairs with equal
nonzero constants lift to denominators with that constant,
so localization gives the stated fiber product. Its two
projection kernels have product zero and sum equal to the
maximal ideal. Every other prime contains exactly one of
them, so its localization agrees with a corresponding
localization of $B$ or $C$. These localizations in height
at most one are fields or DVRs. In fact the two polynomial
rings are unique factorization domains by the
\href{https://github.com/leanprover-community/mathlib4/blob/56a005855ea17e0b0958f73ec801ae140b0a1679/Mathlib/RingTheory/Polynomial/UniqueFactorization.lean\#L128}{fixed Mathlib polynomial factorization theorem}.
A height-one prime in a unique
factorization domain is generated by a prime factor of any
nonzero element it contains. Localizing then gives a principal
maximal ideal in dimension one. The height comparison follows
from the local isomorphism. Each prime chain starts on at
least one component and stays on that component, so $\dim A=2$.
Thus $A$ satisfies $(R_1)$. The element $(x,z)$ is a
nonzerodivisor, since both coordinates are nonzero in
domains. The class of $(x,0)$ modulo $(x,z)$ is nonzero:
a multiplier would have to have residues $1$ and $0$ in
the two components. Every element $(a,b)$ of the maximal
ideal kills this class, since $(a,0)\in A$ and
$(a,b)(x,0)=(a,0)(x,z)$. The quotient has depth zero,
so by \ref{res:DP-05} the ring has depth one and fails $(S_2)$.
\end{example}

\begin{proposition}[Normality and \'{e}tale maps]\label{res:SN-13}
Let $A$ be Noetherian and $A\to B$ be \'{e}tale. If $A$ is
normal, then $B$ is normal. If the map is faithfully flat and
$B$ is normal, then $A$ is normal.
\end{proposition}
\begin{proof}
The ring $B$ is Noetherian by finite presentation and the
Hilbert basis theorem. For $\mathfrak q\in\operatorname{Spec}B$,
put $\mathfrak p=\mathfrak q\cap A$ and consider the flat local
map $R=A_{\mathfrak p}\to S=B_{\mathfrak q}$. The closed fiber
is a field and $\mathfrak m_S=\mathfrak m_R S$ by
\ref{res:ET-12} and \ref{res:ET-13}. By the flat local
depth and dimension formulas \ref{res:CM-10}, the two local
rings have the same depth and the same dimension. Further,
$\mathfrak m_S/\mathfrak m_S^2$ is the scalar extension of
$\mathfrak m_R/\mathfrak m_R^2$ to the residue field of $S$:
tensor the inclusion and quotient maps of the ideals by
the flat $R$-algebra $S$ and use
$(\mathfrak m_RS)^2=\mathfrak m_R^2 S$.
Thus the embedding dimensions also agree. In particular the
two conditions of \ref{res:SN-09} transport from $R$ to $S$.
For the converse, faithful flatness supplies a prime of $B$
above every prime of $A$ by taking a prime of the nonzero
fiber, so the same numerical comparisons transport the
conditions in the other direction. Apply \ref{res:SN-09}.
\end{proof}
\chapter{Gorenstein Rings and Canonical Modules}\label{ch:gorenstein}

Throughout the local arguments, $(R,\mathfrak m,k)$ is a nonzero
Noetherian local ring and $d=\dim R$. We use finite-module depth
from Chapter~\ref{ch:depth}. The homological preliminaries are Baer's
criterion, injective envelopes, injective resolutions, and the two
long exact sequences for $\operatorname{Ext}$, proved in P1-03
\cite[Chapters 2, 10, and 11]{p103}. The canonical module here is
the canonical module of a Cohen--Macaulay local ring as in
\cite[Chapter 3]{brunsherzog}.

\section{Injective dimension and reduction by a regular element}

\begin{definition}[Socle, type, and the local dualizing conditions]
\label{res:GC-01}
For a module $M$, its socle is $0:_M\mathfrak m\simeq
\operatorname{Hom}_R(k,M)$. For $M$ finite and nonzero define
\[
 r_R(M)=\dim_k
 \operatorname{Ext}^{\operatorname{depth}_R M}_R(k,M).
\]
The type of $R$ is $r_R(R)$. Its Gorenstein condition means
that $R$ has finite injective dimension as an $R$-module.
A Noetherian ring is Gorenstein if all its prime localizations
are Gorenstein local rings. For a Cohen--Macaulay local $R$,
a canonical module is a finite maximal Cohen--Macaulay module
of type one and finite injective dimension over $R$.
\end{definition}

\begin{proposition}[Testing injective dimension at the residue field]
\label{res:GC-02}
Let $N$ be finite and $n\ge0$. If
$\operatorname{Ext}^i_R(k,N)=0$ for every $i>n$, then
$\operatorname{injdim}_R N\le n$. The converse holds, and
for finite nonzero $N$ the finite injective dimension, if it
exists, is the greatest $i$ with $\operatorname{Ext}^i_R(k,N)\ne0$.
\end{proposition}
\begin{proof}
For finite modules $M,N$, the groups $\operatorname{Ext}^i_R(M,N)$
are finite: resolve $M$ by finite free modules, taking each next
kernel finite by Noetherianity. Every finite module has a finite
prime filtration. Indeed, in a nonzero remaining quotient choose
an associated prime and the corresponding cyclic submodule, and
repeat. The resulting increasing chain terminates by Noetherianity.

We prove vanishing above $n$ for all finite $M$ by induction on
the dimension of the support. Dimension zero modules have a
finite filtration with factors $k$, for which use the hypothesis
and the long exact sequence. In positive dimension, prime
filtrations reduce to $M=R/\mathfrak p$ of at most that dimension.
If $\mathfrak p\ne\mathfrak m$, choose $x\in\mathfrak m
\setminus\mathfrak p$. Multiplication by $x$ embeds $R/\mathfrak p$
in itself. Its cokernel has strictly smaller dimension, since in
any supporting prime chain one can prepend $\mathfrak p$. The long
exact sequence shows that multiplication by $x$ on
$\operatorname{Ext}^i_R(R/\mathfrak p,N)$ is surjective for $i>n$.
The map is scalar multiplication by commutativity and the scalar
calculation in \ref{res:DP-03}. Nakayama gives vanishing.

Take an injective resolution of $N$. Its $n$th cosyzygy $C$
(with $C=N$ if $n=0$) satisfies
$\operatorname{Ext}^1_R(R/I,C)=\operatorname{Ext}^{n+1}_R(R/I,N)=0$
for every ideal $I$ by dimension shifting. Applying Hom to
$0\to I\to R\to R/I\to0$ shows every map $I\to C$ extends to $R$.
Baer's criterion makes $C$ injective, giving the desired bound.
The converse follows from an injective resolution of length $n$.
For positive finite injective dimension the bound just proved
forces nonvanishing in its top degree. In dimension zero for the
resolution, $\operatorname{Ext}^0_R(k,N)\ne0$ follows from
\ref{res:DP-04}: the first nonzero residue-field Ext of a nonzero
finite module occurs in its finite depth.
\end{proof}

\begin{proposition}[Change of rings with concentrated Ext]
\label{res:GC-03}
Let $S\to R$ be a ring map and let $N$ be an $S$-module. Suppose
$\operatorname{Ext}^j_S(R,N)=0$ for $j\ne c$, and put
$W=\operatorname{Ext}^c_S(R,N)$ with its natural $R$-action.
For every $R$-module $M$, there are natural isomorphisms
\[
 \operatorname{Ext}^{i}_R(M,W)\simeq
 \operatorname{Ext}^{i+c}_S(M,N)\quad (i\ge0),
\]
and $\operatorname{Ext}^j_S(M,N)=0$ for $j<c$.
In particular, for $R=S/(x)$ with $x$ regular on $S$ and $N$,
one has $c=1$ and $W=N/xN$.
\end{proposition}
\begin{proof}
Take an injective resolution $N\to I^\bullet$ over $S$.
The complex $J^\bullet=\operatorname{Hom}_S(R,I^\bullet)$
consists of injective $R$-modules: the map sending an
$R$-map $M\to\operatorname{Hom}_S(R,I^j)$ to its value at $1$
identifies that Hom with $\operatorname{Hom}_S(M,I^j)$,
whose functor is exact on $R$-modules. The inverse sends
$g:M\to I^j$ to $m\mapsto(r\mapsto g(rm))$.

When $c>0$ the differential $J^0\to J^1$ is injective and
splits because its source is injective. The resulting isomorphic
pair in degrees zero and one is a contractible summand. Remove it.
Repeating in the next degree until degree $c$ yields a complex
of injectives starting at degree $c$, with $W$ as its first kernel
and exactness in subsequent degrees. Each removal is a split
decomposition and remains contractible after applying Hom.
Thus $\operatorname{Hom}_R(M,J^\bullet)$ computes the displayed
groups with the indicated shift; by the adjunction it also
computes $\operatorname{Ext}_S^\bullet(M,N)$. The claim is natural
in $M$ by this comparison. For the regular element compute
$\operatorname{Ext}_S^\bullet(R,N)$ using
$0\to S\xrightarrow{x}S\to R\to0$, with $x$ acting injectively
on $N$. This is also the shift of \ref{res:DP-14}.
\end{proof}

\section{Finite length and the Gorenstein criterion}

\begin{theorem}[Finite-length duality]\label{res:GC-04}
Let $E$ be an injective envelope of $k$ over an Artinian local
ring $R$. Then $\operatorname{Hom}_R(-,E)$ is an exact
contravariant duality on finite modules, preserves length,
and its natural evaluation map to the second dual is an
isomorphism. In particular $\ell_R(E)=\ell_R(R)$.
\end{theorem}
\begin{proof}
The inclusion $k\subset E$ is essential, meaning every nonzero
submodule of $E$ meets its image. Thus every simple submodule
of $E$ equals that image, and $\operatorname{Hom}_R(k,E)\simeq k$.
Injectivity makes $\operatorname{Hom}_R(-,E)$ exact
contravariantly. Apply it to a composition series to obtain
a filtration of the dual with the same number of simple factors.
For $M=k$ the evaluation map is nonzero, since some map from
$k$ to $E$ is nonzero on a given nonzero element of $k$. It
is therefore an isomorphism between simples. Apply naturality
of evaluation to a short exact sequence in a composition
series and use exactness of the twice-dual functor to
inductively obtain the assertion for $M$. At $M=R$ the first
dual is $E$, giving the last equality. Existence and essentiality
of the injective envelope are from \cite[Chapter 2]{p103}.
\end{proof}

\begin{proposition}[Self-injectivity and the socle]\label{res:GC-05}
If a Noetherian local ring is self-injective, it is Artinian.
An Artinian local ring is self-injective if and only if its
socle has dimension one. Over a Noetherian local ring of depth
zero, any finite module of finite injective dimension is
injective.
\end{proposition}
\begin{proof}
For the first claim let $x\in\mathfrak m$. Noetherianity makes
the annihilators of the powers of $x$ stabilize, so for some
$n$ the map $x^{n+1}R\to R$ assigning $x^{n+1}r\mapsto x^n r$
is well defined. Self-injectivity extends it to multiplication
by some $u\in R$. Then $(1-ux)x^n=0$ and $1-ux$ is a unit.
Every $x\in\mathfrak m$ is nilpotent. Since $\mathfrak m$ is
finitely generated, it is nilpotent, and the maximal-ideal
power filtration shows that $R$ has finite length.

If $R$ is Artinian and self-injective, every $k$-linear
endomorphism of its socle, followed by inclusion into $R$,
extends to an $R$-endomorphism of $R$. Its restriction to the
socle is multiplication by a scalar of $k$. A vector space
of dimension greater than one has nonscalar endomorphisms. Its
socle is nonzero, since any finite module of positive length
contains a simple submodule. Therefore the socle is one-dimensional.
Conversely, when $\operatorname{Soc}(R)=k$, extend its inclusion
into $E$ to $R\to E$. The kernel is zero because every nonzero
ideal contains a simple submodule and so meets the socle.
By \ref{res:GC-04}, source and target have the same length,
so $R\simeq E$ is injective.

For the last claim, depth zero provides an embedding $k\hookrightarrow R$
by \ref{res:DP-01}. If $\operatorname{injdim}_R N=n>0$, apply
$\operatorname{Ext}_R^\bullet(-,N)$ to $0\to k\to R\to R/k\to0$.
The groups $\operatorname{Ext}^n_R(R,N)$ and
$\operatorname{Ext}^{n+1}_R(R/k,N)$ are zero, forcing
$\operatorname{Ext}^n_R(k,N)=0$. This contradicts
\ref{res:GC-02}.
\end{proof}

\begin{theorem}[Gorenstein local rings and type one]\label{res:GC-06}
For a Noetherian local ring $R$ the following are equivalent:
\begin{enumerate}
\item $R$ is Gorenstein;
\item $R$ is Cohen--Macaulay of type one;
\item $\operatorname{Ext}^i_R(k,R)=0$ for $i$ sufficiently large.
\end{enumerate}
In this case $\operatorname{injdim}_R R=\dim R$, and the residue-field
Ext is zero except in degree $\dim R$, where it is $k$.
If $x\in\mathfrak m$ is regular on $R$, then $R$ is Gorenstein
exactly when $R/(x)$ is Gorenstein.
\end{theorem}
\begin{proof}
Conditions (1) and (3) are equivalent by \ref{res:GC-02}.
For a regular element of the maximal ideal, \ref{res:GC-03}
identifies residue-field Ext of $R/(x)$ with that of $R$ shifted
up by one; the degree-zero Ext for $R$ is zero. Thus eventual
vanishing is equivalent on both rings, proving the assertion
about regular quotients.

Suppose (1), and choose a maximal regular sequence of length
$s=\operatorname{depth}R$. By \ref{res:DP-05}, the quotient
$R_0$ has depth zero. It has finite injective dimension over
itself by the preceding comparison. By \ref{res:GC-05}, it is
self-injective Artinian of socle dimension one. The regular-element
dimension drop \ref{res:CM-02} applied $s$ times gives $s=d$.
Repeated shifts in \ref{res:GC-03} turn the Ext over $R_0$
into one copy of $k$ in degree $d$ over $R$, giving type one
and the asserted injective dimension by \ref{res:GC-02}.

Conversely, if $R$ is Cohen--Macaulay of type one, a system of
parameters is regular by \ref{res:CM-04}. Its quotient has socle
dimension one by \ref{res:DP-05} or by the same Ext comparison.
It is self-injective by \ref{res:GC-05}, and the regular-quotient
equivalence along the sequence gives (1).
\end{proof}

\begin{corollary}[Regular local rings, complete intersections, and localization]
\label{res:GC-07}
Regular local rings are Gorenstein. Quotients of local Gorenstein
rings by regular sequences in the maximal ideal are Gorenstein.
In particular, complete intersections presented as quotients
of regular local rings are Gorenstein. A localization of a
Gorenstein local ring at a prime is Gorenstein.
\end{corollary}
\begin{proof}
For a regular local ring, its minimal maximal-ideal generators
are a regular sequence and the quotient is the residue field
\cite[ALG-05-C03-S03]{p103}. Use \ref{res:GC-06} along this
sequence and along each other regular sequence in the statement.
For localization, let $R$ have injective dimension $n$ over itself.
For each ideal of $R_{\mathfrak p}$, choose an ideal $I\subseteq R$
whose extension gives that ideal, for instance its contraction.
The finite free resolution computation gives
\[
 \operatorname{Ext}^i_{R_{\mathfrak p}}
 (R_{\mathfrak p}/I R_{\mathfrak p},R_{\mathfrak p})
 \simeq \operatorname{Ext}^i_R(R/I,R)_{\mathfrak p}.
\]
Indeed, localization is exact and commutes with the Hom from
each finite free term. This vanishes for $i>n$.
The cosyzygy and Baer argument from \ref{res:GC-02}, applied in
$R_{\mathfrak p}$, yields injective dimension at most $n$ there.
\end{proof}

\begin{proposition}[Flat local maps, completion, and \'{e}tale maps]
\label{res:GC-08}
For a flat local map of Noetherian local rings $R\to S$ with
closed fiber $F=S/\mathfrak m_R S$, the ring $S$ is Gorenstein
if and only if $R$ and $F$ are Gorenstein. In particular a
Noetherian local ring and its completion are Gorenstein
simultaneously. Gorensteinness ascends along \'{e}tale maps
of Noetherian rings and descends along faithfully flat
\'etale maps.
\end{proposition}
\begin{proof}
The depth and dimension equalities for flat local maps
\ref{res:CM-10} give the Cohen--Macaulay equivalence.
When the three rings are Cohen--Macaulay, their types satisfy
$r_S(S)=r_R(R)r_F(F)$ by \ref{res:CM-18} with $M=R$ and
$N=S$. Each type is a positive integer by the first-nonzero
Ext characterization. Hence type one for $S$ is equivalent
to type one for both factors. Apply \ref{res:GC-06}.

The completion map is faithfully flat and local and its
closed fiber is the residue field
\cite[ALG-05-C04-S01]{p103}. For an \'{e}tale map,
the local homomorphism at a prime is flat by \ref{res:ET-12}
and its closed fiber is a field by \ref{res:ET-13}. For
faithfully flat descent every prime has a prime above it,
as seen in its nonzero fiber. These facts give the assertions
by the flat local statement.
\end{proof}

\section{Construction of a canonical module}

\begin{theorem}[Ext of a Cohen--Macaulay module over a Gorenstein base]
\label{res:GC-09}
Let $S$ be Gorenstein local of dimension $s$ and let $M$ be
a nonzero finite Cohen--Macaulay $S$-module of dimension $e$.
Put $c=s-e$. Then $\operatorname{Ext}_S^i(M,S)=0$ for $i\ne c$,
and $D=\operatorname{Ext}_S^c(M,S)$ is nonzero and
Cohen--Macaulay of dimension $e$. If $x$ is regular on $M$
and belongs to the maximal ideal, then $x$ acts injectively
on $D$ and
\[
 D/xD\simeq\operatorname{Ext}_S^{c+1}(M/xM,S).
\]
\end{theorem}
\begin{proof}
In dimension zero, $M$ has finite length. By \ref{res:GC-06}
the Ext of the residue field over $S$ is one copy of that
field in degree $s$ and zero elsewhere. The long exact
sequence along a composition series yields concentration
in degree $s$. Moreover the Ext in that degree has a
filtration of the same positive length, since the previous
and next terms of the long exact sequence vanish.

For $e>0$ choose a regular $x$ in the maximal ideal. The
module $\overline M=M/xM$ is nonzero Cohen--Macaulay of dimension
$e-1$ by \ref{res:CM-05}. The long exact sequence of
$0\to M\xrightarrow{x}M\to\overline M\to0$ contains
\[
 \operatorname{Ext}^i_S(\overline M,S)\longrightarrow
 \operatorname{Ext}^i_S(M,S)\xrightarrow{x}
 \operatorname{Ext}^i_S(M,S)\longrightarrow
 \operatorname{Ext}^{i+1}_S(\overline M,S).
\]
By induction the rightmost group vanishes for $i\ne c$.
Then multiplication by $x$ on the middle group is surjective,
so it vanishes by finiteness and Nakayama. For $i=c$ the
leftmost group vanishes and the next group
$\operatorname{Ext}^{c+1}_S(M,S)$ also vanishes. Thus $x$
is regular on $D$, its quotient is the nonzero indicated
group, and the group on the right is Cohen--Macaulay of
dimension $e-1$ by induction. The proof uses any such
regular $x$. The support of $D$ lies in the support of
$M$, by localization of the finite free resolution as
in \ref{res:GC-07}. By \ref{res:DP-05}, its depth equals
$1+\operatorname{depth}(D/xD)=e$. The depth bound
\ref{res:DP-02} and the support containment force its
dimension to equal $e$.
\end{proof}

\begin{theorem}[Canonical module of a Cohen--Macaulay quotient]
\label{res:GC-10}
Let $S\twoheadrightarrow R$ be a surjection from a Gorenstein
local ring of dimension $s$ to a nonzero Cohen--Macaulay local
ring of dimension $d$. Then
\[
 \omega_R=\operatorname{Ext}^{s-d}_S(R,S)
\]
is a canonical module. It has injective dimension $d$ over $R$,
and $\operatorname{Ext}^i_R(k,\omega_R)$ is $k$ in degree $d$
and zero in all other degrees. For an $R$-module $M$ the
change-of-rings identification is
\[
 \operatorname{Ext}^i_R(M,\omega_R)
 \simeq\operatorname{Ext}^{i+s-d}_S(M,S)\quad(i\ge0).
\]
\end{theorem}
\begin{proof}
A surjection of local rings identifies the residue fields
and makes regular sequences for a finite target module
exactly the images of regular sequences acting on it
over the source. Supporting primes likewise correspond,
so depth and support dimension agree over both rings
(\ref{res:CM-05}). Thus $R$ is Cohen--Macaulay as an
$S$-module of dimension $d$. By \ref{res:GC-09} its Ext is
concentrated in degree $s-d$ and the resulting $\omega_R$
is finite, nonzero and Cohen--Macaulay of dimension $d$.
Its natural $R$-action comes from the first argument of
Ext and agrees with the scalar $S$-action via $S\to R$,
by the scalar calculation of \ref{res:DP-03}. Hence
$\omega_R$ is maximal Cohen--Macaulay over $R$.

Apply \ref{res:GC-03} to the concentrated group to obtain
the stated identification. At $M=k$, \ref{res:GC-06}
gives the concentration in degree $d$ with dimension one.
By \ref{res:GC-02}, the injective dimension over $R$ is
exactly $d$. This proves the canonical module conditions.
\end{proof}

\section{Duality on maximal Cohen--Macaulay modules}

\begin{proposition}[Canonical reduction and Ext vanishing]
\label{res:GC-11}
Let a Cohen--Macaulay local ring $R$ have a canonical module
$\omega$. If $x$ is the first element of a system of parameters
of positive length, then $\omega/x\omega$ is canonical for
$R/(x)$. For every finite maximal Cohen--Macaulay module $M$,
\[
  \operatorname{Ext}^{i}_R(M,\omega)=0 \qquad (i>0).
\]
\end{proposition}
\begin{proof}
The parameter element $x$ is regular on $R$, $\omega$, and
every nonzero maximal Cohen--Macaulay module $M$ by
\ref{res:CM-04}. By \ref{res:CM-05}, their quotients have
the expected Cohen--Macaulay depth. The shift
\ref{res:GC-03} with $N=\omega$ identifies the residue-field
Ext of $\omega/x\omega$ over $R/(x)$ with the next degree
over $R$. It preserves type one and eventual vanishing,
so by \ref{res:GC-02} the quotient is canonical.

Prove the Ext vanishing by induction on $\dim R$.
In dimension zero, $\omega$ is injective by \ref{res:GC-05}.
In positive dimension use the sequence
$0\to M\xrightarrow{x}M\to M/xM\to0$ and the shift
\ref{res:GC-03} to identify
$\operatorname{Ext}^{i+1}_R(M/xM,\omega)$ with
$\operatorname{Ext}^i_{R/(x)}(M/xM,\omega/x\omega)$, which
is zero for $i>0$ by induction. Therefore multiplication
by $x$ on $\operatorname{Ext}^i_R(M,\omega)$ is surjective.
This group is finite and $x$ is in the maximal ideal,
so it vanishes by Nakayama.
\end{proof}

\begin{theorem}[Maximal Cohen--Macaulay duality]\label{res:GC-12}
If $R$ has a canonical module $\omega$, the functor
$D(M)=\operatorname{Hom}_R(M,\omega)$ is an exact
contravariant duality on finite maximal Cohen--Macaulay
modules, with the zero module adjoined. The natural map
$M\to D(D(M))$ is an isomorphism, and the scalar map
$R\to\operatorname{End}_R(\omega)$ is an isomorphism.
\end{theorem}
\begin{proof}
In dimension zero, $\omega$ is injective by \ref{res:GC-05}
and has simple socle by its type. Its simple socle is
essential since every nonzero submodule of a finite-length
module contains a simple submodule. Thus it is an injective
envelope of $k$ and \ref{res:GC-04} gives the assertion.

In positive dimension choose a first parameter $x$. The
element acts injectively on $D(M)$ since it acts injectively
on $\omega$. Write bars for reduction by $x$.
By applying $\operatorname{Hom}_R(M,-)$ to multiplication
by $x$ on $\omega$, and using \ref{res:GC-11}, one obtains
\[
 \overline{D(M)}
 \simeq \operatorname{Hom}_{R/(x)}(\overline M,\overline\omega).
\]
This isomorphism sends a map to its reduction, since any
map into $\overline\omega$ factors through $\overline M$.
By induction the quotient is maximal Cohen--Macaulay and
nonzero when $M\ne0$. By the regular-element depth formula
\ref{res:DP-05}, $D(M)$ is maximal Cohen--Macaulay. Applying
the same comparison to $D(M)$ identifies the reduction
of the evaluation map $M\to D^2(M)$ with evaluation over
$R/(x)$, which is an isomorphism by induction.

A homomorphism of finite modules whose reduction by $x$
is onto is itself onto by Nakayama. For such a surjection
with kernel $K$, if the target is $x$-regular then
$K\cap xM=xK$: from $xm\in K$ one obtains $m\in K$ by
injectivity of $x$ on the target. Therefore $K/xK$
injects into the reduced source and vanishes when the
reduced map is an isomorphism. Nakayama gives $K=0$.
This proves evaluation is an isomorphism. Exactness on
short exact sequences of maximal Cohen--Macaulay modules
follows by applying Hom and the vanishing of $\operatorname{Ext}^1$
in \ref{res:GC-11}. At $M=R$, identify $D(R)$ with
$\omega$ by evaluation at $1$. Then $R\to D^2(R)$ is
exactly the scalar map $R\to\operatorname{End}_R(\omega)$.
\end{proof}

\begin{corollary}[Uniqueness and the Gorenstein test]
\label{res:GC-13}
Canonical modules of a Cohen--Macaulay local ring are
isomorphic if they exist. A Cohen--Macaulay local ring
with canonical module $\omega$ is Gorenstein exactly when
$\omega\simeq R$. In particular, the isomorphism type
of the module in \ref{res:GC-10} is independent of the
chosen Gorenstein local presentation.
\end{corollary}
\begin{proof}
In dimension zero all canonical modules are injective
envelopes of $k$, as in \ref{res:GC-12}. For positive
dimension, reduce two canonical modules by a common
first parameter $x$. Their reductions are canonical by
\ref{res:GC-11}, hence isomorphic by induction. The
isomorphism of reductions lifts to a homomorphism
between the original modules: apply the Hom sequence
for multiplication by $x$ on the target, and use
$\operatorname{Ext}^1_R(\omega_1,\omega_2)=0$ from
\ref{res:GC-11}. The lift is an isomorphism by the
finite-module argument in \ref{res:GC-12}, since $x$
acts injectively on the target. If $\omega\simeq R$,
then $R$ has finite injective dimension. Conversely,
if $R$ is Gorenstein, \ref{res:GC-06} says that $R$
itself satisfies the canonical conditions, and
uniqueness gives the isomorphism.
\end{proof}

\begin{example}[Two Artinian rings and a hypersurface]\label{res:GC-14}
The Artinian local ring $k[t]/(t^n)$ for $n\ge1$ has
one-dimensional socle spanned by $t^{n-1}$, so is
Gorenstein. The Artinian local ring
$k[x,y]/(x,y)^2$ is Cohen--Macaulay of dimension zero,
with socle the two-dimensional span of $x,y$. It is
not Gorenstein by \ref{res:GC-05}, and its canonical
module is its injective envelope of $k$ by
\ref{res:GC-04}. If $S$ is regular local and $f$ is a
nonzerodivisor in its maximal ideal, the quotient
$S/(f)$ is Gorenstein by \ref{res:GC-07}. Its canonical
module in \ref{res:GC-10} is
$\operatorname{Ext}^1_S(S/(f),S)\simeq S/(f)$, computed
by the two-term resolution for $f$.
\end{example}

\section*{Passage to local duality}

Finite-length duality and the canonical module of a
Cohen--Macaulay ring give duality on maximal
Cohen--Macaulay modules here. Local cohomology and the
normalized dualizing complex extend these constructions
beyond that category. The remaining existence and
local-duality statements are developed using those
additional tools \cite[Section 3.5]{brunsherzog}.

\backmatter
\begin{thebibliography}{9}
\addcontentsline{toc}{chapter}{Bibliography}

\bibitem{eisenbud}
David Eisenbud.
\emph{Commutative Algebra with a View Toward Algebraic Geometry}.
Graduate Texts in Mathematics 150. Springer, 1995. Chapters 16--18 and 21.

\bibitem{matsumura}
Hideyuki Matsumura.
\emph{Commutative Ring Theory}. Translated by M. Reid.
Cambridge Studies in Advanced Mathematics 8.
Cambridge University Press, corrected paperback, 1989. Sections 11, 16--17, and 25--26.

\bibitem{brunsherzog}
Winfried Bruns and J\"urgen Herzog.
\emph{Cohen--Macaulay Rings}. Cambridge Studies in Advanced Mathematics 39.
Cambridge University Press, revised paperback, 1998. Chapters 1--3.

\bibitem{hartshorne}
Robin Hartshorne.
\emph{Deformation Theory}. Graduate Texts in Mathematics 257.
Springer, 2010. Sections 3--4.

\bibitem{stacks}
The Stacks Project Authors.
\emph{The Stacks Project}. Chapters on commutative algebra, more on algebra, and dualizing complexes.
Tags identify individual statements. Reference snapshot \texttt{ed88ff78}.
\url{https://stacks.math.columbia.edu/}.

\bibitem{p103}
MathGenome.
\emph{Algebra Core} (P1-03), accepted source notes and formalization.
MathGenome v0.21.0.

\bibitem{d07}
MathGenome.
\emph{Abstract Algebra} (D07), accepted source notes and formalization.
MathGenome v0.21.0.
\end{thebibliography}
\end{document}
